% dataset: math-formula-atlas
% generated_at: 2026-03-27T18:09:25.420Z
% start_url: https://www.mathsisfun.com/algebra/exponent-laws.html
% scope_path_prefix: /
% total_pages: 39
% total_formulas: 376

% formula_record_start
% formula_id: exponent-laws-html-001
% page_slug: exponent-laws-html
% page_title: Laws of Exponents
% page_url: https://www.mathsisfun.com/algebra/exponent-laws.html
% subject_slug: algebra
% subject_path: /algebra/
% section: Laws of Exponents
% subsection: 
% formula_name: 
% source: mathsisfun-large-span
% index_global: 1
% index_in_page: 1
\[x^{1} = x\]
% formula_record_end

% formula_record_start
% formula_id: exponent-laws-html-002
% page_slug: exponent-laws-html
% page_title: Laws of Exponents
% page_url: https://www.mathsisfun.com/algebra/exponent-laws.html
% subject_slug: algebra
% subject_path: /algebra/
% section: Laws of Exponents
% subsection: 
% formula_name: 
% source: mathsisfun-large-span
% index_global: 2
% index_in_page: 2
\[x^{0} = 1\]
% formula_record_end

% formula_record_start
% formula_id: exponent-laws-html-003
% page_slug: exponent-laws-html
% page_title: Laws of Exponents
% page_url: https://www.mathsisfun.com/algebra/exponent-laws.html
% subject_slug: algebra
% subject_path: /algebra/
% section: Laws of Exponents
% subsection: 
% formula_name: 
% source: mathsisfun-large-span
% index_global: 3
% index_in_page: 3
\[x^{-1} = 1/x\]
% formula_record_end

% formula_record_start
% formula_id: exponent-laws-html-004
% page_slug: exponent-laws-html
% page_title: Laws of Exponents
% page_url: https://www.mathsisfun.com/algebra/exponent-laws.html
% subject_slug: algebra
% subject_path: /algebra/
% section: Laws of Exponents
% subsection: 
% formula_name: 
% source: mathsisfun-large-span
% index_global: 4
% index_in_page: 4
\[x^{m}x^{n} = x^{m+n}\]
% formula_record_end

% formula_record_start
% formula_id: exponent-laws-html-005
% page_slug: exponent-laws-html
% page_title: Laws of Exponents
% page_url: https://www.mathsisfun.com/algebra/exponent-laws.html
% subject_slug: algebra
% subject_path: /algebra/
% section: Laws of Exponents
% subsection: 
% formula_name: 
% source: mathsisfun-large-span
% index_global: 5
% index_in_page: 5
\[x^{m}/x^{n} = x^{m-n}\]
% formula_record_end

% formula_record_start
% formula_id: exponent-laws-html-006
% page_slug: exponent-laws-html
% page_title: Laws of Exponents
% page_url: https://www.mathsisfun.com/algebra/exponent-laws.html
% subject_slug: algebra
% subject_path: /algebra/
% section: Laws of Exponents
% subsection: 
% formula_name: 
% source: mathsisfun-large-span
% index_global: 6
% index_in_page: 6
\[(x^{m})^{n} = x^{mn}\]
% formula_record_end

% formula_record_start
% formula_id: exponent-laws-html-007
% page_slug: exponent-laws-html
% page_title: Laws of Exponents
% page_url: https://www.mathsisfun.com/algebra/exponent-laws.html
% subject_slug: algebra
% subject_path: /algebra/
% section: Laws of Exponents
% subsection: 
% formula_name: 
% source: mathsisfun-large-span
% index_global: 7
% index_in_page: 7
\[(xy)^{n} = x^{n}y^{n}\]
% formula_record_end

% formula_record_start
% formula_id: exponent-laws-html-008
% page_slug: exponent-laws-html
% page_title: Laws of Exponents
% page_url: https://www.mathsisfun.com/algebra/exponent-laws.html
% subject_slug: algebra
% subject_path: /algebra/
% section: Laws of Exponents
% subsection: 
% formula_name: 
% source: mathsisfun-large-span
% index_global: 8
% index_in_page: 8
\[(x/y)^{n} = x^{n}/y^{n}\]
% formula_record_end

% formula_record_start
% formula_id: exponent-laws-html-009
% page_slug: exponent-laws-html
% page_title: Laws of Exponents
% page_url: https://www.mathsisfun.com/algebra/exponent-laws.html
% subject_slug: algebra
% subject_path: /algebra/
% section: Laws of Exponents
% subsection: 
% formula_name: 
% source: mathsisfun-large-span
% index_global: 9
% index_in_page: 9
\[x^{-n} = 1/x^{n}\]
% formula_record_end

% formula_record_start
% formula_id: exponent-laws-html-010
% page_slug: exponent-laws-html
% page_title: Laws of Exponents
% page_url: https://www.mathsisfun.com/algebra/exponent-laws.html
% subject_slug: algebra
% subject_path: /algebra/
% section: Laws of Exponents
% subsection: 
% formula_name: 
% source: mathsisfun-large-span
% index_global: 10
% index_in_page: 10
\[x^{1} = x\]
% formula_record_end

% formula_record_start
% formula_id: exponent-laws-html-011
% page_slug: exponent-laws-html
% page_title: Laws of Exponents
% page_url: https://www.mathsisfun.com/algebra/exponent-laws.html
% subject_slug: algebra
% subject_path: /algebra/
% section: Laws of Exponents
% subsection: 
% formula_name: 
% source: mathsisfun-large-span
% index_global: 11
% index_in_page: 11
\[x^{0} = 1\]
% formula_record_end

% formula_record_start
% formula_id: exponent-laws-html-012
% page_slug: exponent-laws-html
% page_title: Laws of Exponents
% page_url: https://www.mathsisfun.com/algebra/exponent-laws.html
% subject_slug: algebra
% subject_path: /algebra/
% section: Laws of Exponents
% subsection: 
% formula_name: 
% source: mathsisfun-large-span
% index_global: 12
% index_in_page: 12
\[x^{-1} = 1/x\]
% formula_record_end

% formula_record_start
% formula_id: exponent-laws-html-013
% page_slug: exponent-laws-html
% page_title: Laws of Exponents
% page_url: https://www.mathsisfun.com/algebra/exponent-laws.html
% subject_slug: algebra
% subject_path: /algebra/
% section: Laws of Exponents
% subsection: 
% formula_name: x1 = x
% source: mathsisfun-table-law
% index_global: 13
% index_in_page: 13
\textbf{Label: x1 = x}
\[x^{1} = x\]
% formula_record_end

% formula_record_start
% formula_id: exponent-laws-html-014
% page_slug: exponent-laws-html
% page_title: Laws of Exponents
% page_url: https://www.mathsisfun.com/algebra/exponent-laws.html
% subject_slug: algebra
% subject_path: /algebra/
% section: Laws of Exponents
% subsection: 
% formula_name: Example: x1 = x
% source: mathsisfun-table-example
% index_global: 14
% index_in_page: 14
\textbf{Label: Example: x1 = x}
\[6^{1}=6\]
% formula_record_end

% formula_record_start
% formula_id: exponent-laws-html-015
% page_slug: exponent-laws-html
% page_title: Laws of Exponents
% page_url: https://www.mathsisfun.com/algebra/exponent-laws.html
% subject_slug: algebra
% subject_path: /algebra/
% section: Laws of Exponents
% subsection: 
% formula_name: x0 = 1
% source: mathsisfun-table-law
% index_global: 15
% index_in_page: 15
\textbf{Label: x0 = 1}
\[x^{0} = 1\]
% formula_record_end

% formula_record_start
% formula_id: exponent-laws-html-016
% page_slug: exponent-laws-html
% page_title: Laws of Exponents
% page_url: https://www.mathsisfun.com/algebra/exponent-laws.html
% subject_slug: algebra
% subject_path: /algebra/
% section: Laws of Exponents
% subsection: 
% formula_name: Example: x0 = 1
% source: mathsisfun-table-example
% index_global: 16
% index_in_page: 16
\textbf{Label: Example: x0 = 1}
\[7^{0}=1\]
% formula_record_end

% formula_record_start
% formula_id: exponent-laws-html-017
% page_slug: exponent-laws-html
% page_title: Laws of Exponents
% page_url: https://www.mathsisfun.com/algebra/exponent-laws.html
% subject_slug: algebra
% subject_path: /algebra/
% section: Laws of Exponents
% subsection: 
% formula_name: x-1 = 1/x
% source: mathsisfun-table-law
% index_global: 17
% index_in_page: 17
\textbf{Label: x-1 = 1/x}
\[x^{-1} = 1/x\]
% formula_record_end

% formula_record_start
% formula_id: exponent-laws-html-018
% page_slug: exponent-laws-html
% page_title: Laws of Exponents
% page_url: https://www.mathsisfun.com/algebra/exponent-laws.html
% subject_slug: algebra
% subject_path: /algebra/
% section: Laws of Exponents
% subsection: 
% formula_name: Example: x-1 = 1/x
% source: mathsisfun-table-example
% index_global: 18
% index_in_page: 18
\textbf{Label: Example: x-1 = 1/x}
\[4^{-1}=\frac{1}{4}\]
% formula_record_end

% formula_record_start
% formula_id: exponent-laws-html-019
% page_slug: exponent-laws-html
% page_title: Laws of Exponents
% page_url: https://www.mathsisfun.com/algebra/exponent-laws.html
% subject_slug: algebra
% subject_path: /algebra/
% section: Laws of Exponents
% subsection: 
% formula_name: xmxn = xm+n
% source: mathsisfun-table-law
% index_global: 19
% index_in_page: 19
\textbf{Label: xmxn = xm+n}
\[x^{m}x^{n} = xm+n\]
% formula_record_end

% formula_record_start
% formula_id: exponent-laws-html-020
% page_slug: exponent-laws-html
% page_title: Laws of Exponents
% page_url: https://www.mathsisfun.com/algebra/exponent-laws.html
% subject_slug: algebra
% subject_path: /algebra/
% section: Laws of Exponents
% subsection: 
% formula_name: Example: xmxn = xm+n
% source: mathsisfun-table-example
% index_global: 20
% index_in_page: 20
\textbf{Label: Example: xmxn = xm+n}
\[x^{2}x^{3}=x^{2+3}=x^{5}\]
% formula_record_end

% formula_record_start
% formula_id: exponent-laws-html-021
% page_slug: exponent-laws-html
% page_title: Laws of Exponents
% page_url: https://www.mathsisfun.com/algebra/exponent-laws.html
% subject_slug: algebra
% subject_path: /algebra/
% section: Laws of Exponents
% subsection: 
% formula_name: xm/xn = xm-n
% source: mathsisfun-table-law
% index_global: 21
% index_in_page: 21
\textbf{Label: xm/xn = xm-n}
\[\frac{x^{m}}{x^{n}} = xm-n\]
% formula_record_end

% formula_record_start
% formula_id: exponent-laws-html-022
% page_slug: exponent-laws-html
% page_title: Laws of Exponents
% page_url: https://www.mathsisfun.com/algebra/exponent-laws.html
% subject_slug: algebra
% subject_path: /algebra/
% section: Laws of Exponents
% subsection: 
% formula_name: Example: xm/xn = xm-n
% source: mathsisfun-table-example
% index_global: 22
% index_in_page: 22
\textbf{Label: Example: xm/xn = xm-n}
\[x^{6}/x^{2} = x^{6}-2 = x^{4}\]
% formula_record_end

% formula_record_start
% formula_id: exponent-laws-html-023
% page_slug: exponent-laws-html
% page_title: Laws of Exponents
% page_url: https://www.mathsisfun.com/algebra/exponent-laws.html
% subject_slug: algebra
% subject_path: /algebra/
% section: Laws of Exponents
% subsection: 
% formula_name: (xm)n = xmn
% source: mathsisfun-table-law
% index_global: 23
% index_in_page: 23
\textbf{Label: (xm)n = xmn}
\[(x^{m})^{n} = xmn\]
% formula_record_end

% formula_record_start
% formula_id: exponent-laws-html-024
% page_slug: exponent-laws-html
% page_title: Laws of Exponents
% page_url: https://www.mathsisfun.com/algebra/exponent-laws.html
% subject_slug: algebra
% subject_path: /algebra/
% section: Laws of Exponents
% subsection: 
% formula_name: Example: (xm)n = xmn
% source: mathsisfun-table-example
% index_global: 24
% index_in_page: 24
\textbf{Label: Example: (xm)n = xmn}
\[(x^{2})3 = x^{2}\times 3 = x^{6}\]
% formula_record_end

% formula_record_start
% formula_id: exponent-laws-html-025
% page_slug: exponent-laws-html
% page_title: Laws of Exponents
% page_url: https://www.mathsisfun.com/algebra/exponent-laws.html
% subject_slug: algebra
% subject_path: /algebra/
% section: Laws of Exponents
% subsection: 
% formula_name: (xy)n = xnyn
% source: mathsisfun-table-law
% index_global: 25
% index_in_page: 25
\textbf{Label: (xy)n = xnyn}
\[(xy)^{n} = xnyn\]
% formula_record_end

% formula_record_start
% formula_id: exponent-laws-html-026
% page_slug: exponent-laws-html
% page_title: Laws of Exponents
% page_url: https://www.mathsisfun.com/algebra/exponent-laws.html
% subject_slug: algebra
% subject_path: /algebra/
% section: Laws of Exponents
% subsection: 
% formula_name: Example: (xy)n = xnyn
% source: mathsisfun-table-example
% index_global: 26
% index_in_page: 26
\textbf{Label: Example: (xy)n = xnyn}
\[(xy)3 = x^{3}y^{3}\]
% formula_record_end

% formula_record_start
% formula_id: exponent-laws-html-027
% page_slug: exponent-laws-html
% page_title: Laws of Exponents
% page_url: https://www.mathsisfun.com/algebra/exponent-laws.html
% subject_slug: algebra
% subject_path: /algebra/
% section: Laws of Exponents
% subsection: 
% formula_name: (x/y)n = xn/yn
% source: mathsisfun-table-law
% index_global: 27
% index_in_page: 27
\textbf{Label: (x/y)n = xn/yn}
\[(\frac{x}{y})^{n} = xn/yn\]
% formula_record_end

% formula_record_start
% formula_id: exponent-laws-html-028
% page_slug: exponent-laws-html
% page_title: Laws of Exponents
% page_url: https://www.mathsisfun.com/algebra/exponent-laws.html
% subject_slug: algebra
% subject_path: /algebra/
% section: Laws of Exponents
% subsection: 
% formula_name: Example: (x/y)n = xn/yn
% source: mathsisfun-table-example
% index_global: 28
% index_in_page: 28
\textbf{Label: Example: (x/y)n = xn/yn}
\[(x/y)2 = x^{2} / y^{2}\]
% formula_record_end

% formula_record_start
% formula_id: exponent-laws-html-029
% page_slug: exponent-laws-html
% page_title: Laws of Exponents
% page_url: https://www.mathsisfun.com/algebra/exponent-laws.html
% subject_slug: algebra
% subject_path: /algebra/
% section: Laws of Exponents
% subsection: 
% formula_name: x-n = 1/xn
% source: mathsisfun-table-law
% index_global: 29
% index_in_page: 29
\textbf{Label: x-n = 1/xn}
\[x^{-n} = 1/xn\]
% formula_record_end

% formula_record_start
% formula_id: exponent-laws-html-030
% page_slug: exponent-laws-html
% page_title: Laws of Exponents
% page_url: https://www.mathsisfun.com/algebra/exponent-laws.html
% subject_slug: algebra
% subject_path: /algebra/
% section: Laws of Exponents
% subsection: 
% formula_name: Example: x-n = 1/xn
% source: mathsisfun-table-example
% index_global: 30
% index_in_page: 30
\textbf{Label: Example: x-n = 1/xn}
\[x-3 = 1/x^{3}\]
% formula_record_end

% formula_record_start
% formula_id: exponent-laws-html-031
% page_slug: exponent-laws-html
% page_title: Laws of Exponents
% page_url: https://www.mathsisfun.com/algebra/exponent-laws.html
% subject_slug: algebra
% subject_path: /algebra/
% section: Laws of Exponents
% subsection: 
% formula_name: xm/n = n√xm = (n√x )m
% source: mathsisfun-table-law
% index_global: 31
% index_in_page: 31
\textbf{Label: xm/n = n√xm = (n√x )m}
\[x^{m/n} = n\sqrt{}xm = (n\sqrt{}x )m\]
% formula_record_end

% formula_record_start
% formula_id: exponent-laws-html-032
% page_slug: exponent-laws-html
% page_title: Laws of Exponents
% page_url: https://www.mathsisfun.com/algebra/exponent-laws.html
% subject_slug: algebra
% subject_path: /algebra/
% section: Laws of Exponents
% subsection: 
% formula_name: Example: xm/n = n√xm = (n√x )m
% source: mathsisfun-table-example
% index_global: 32
% index_in_page: 32
\textbf{Label: Example: xm/n = n√xm = (n√x )m}
\[x\\frac{2}{3} = 3\sqrt{}x^{2} = (3\sqrt{}x )2\]
% formula_record_end

% formula_record_start
% formula_id: exponent-laws-html-033
% page_slug: exponent-laws-html
% page_title: Laws of Exponents
% page_url: https://www.mathsisfun.com/algebra/exponent-laws.html
% subject_slug: algebra
% subject_path: /algebra/
% section: Uncategorized
% subsection: 
% formula_name: 
% source: formula-text-fallback
% index_global: 33
% index_in_page: 33
\[In this example: 8^{2} = 8 \times 8 = 64\]
% formula_record_end

% formula_record_start
% formula_id: exponent-laws-html-034
% page_slug: exponent-laws-html
% page_title: Laws of Exponents
% page_url: https://www.mathsisfun.com/algebra/exponent-laws.html
% subject_slug: algebra
% subject_path: /algebra/
% section: Laws Explained
% subsection: 
% formula_name: 
% source: formula-text-fallback
% index_global: 34
% index_in_page: 34
\[The first three laws above (x^{1} = x, x^{0} = 1 and x^{-1} = 1/x) are just part of the natural sequence of exponents. Have a look at this:\]
% formula_record_end

% formula_record_start
% formula_id: exponent-laws-html-035
% page_slug: exponent-laws-html
% page_title: Laws of Exponents
% page_url: https://www.mathsisfun.com/algebra/exponent-laws.html
% subject_slug: algebra
% subject_path: /algebra/
% section: The law that xmxn = xm+n
% subsection: 
% formula_name: 
% source: formula-text-fallback
% index_global: 35
% index_in_page: 35
\[With x^{m}x^{n}, how many times do we end up multiplying "x"? Answer: first "m" times, then by another "n" times, for a total of "m+n" times.\]
% formula_record_end

% formula_record_start
% formula_id: exponent-laws-html-036
% page_slug: exponent-laws-html
% page_title: Laws of Exponents
% page_url: https://www.mathsisfun.com/algebra/exponent-laws.html
% subject_slug: algebra
% subject_path: /algebra/
% section: The law that xm/xn = xm-n
% subsection: 
% formula_name: 
% source: formula-text-fallback
% index_global: 36
% index_in_page: 36
\[(Remember that x/x = 1, so every time you see an x "above the line" and one "below the line" you can cancel them out.)\]
% formula_record_end

% formula_record_start
% formula_id: exponent-laws-html-037
% page_slug: exponent-laws-html
% page_title: Laws of Exponents
% page_url: https://www.mathsisfun.com/algebra/exponent-laws.html
% subject_slug: algebra
% subject_path: /algebra/
% section: The law that xm/xn = xm-n
% subsection: 
% formula_name: 
% source: formula-text-fallback
% index_global: 37
% index_in_page: 37
\[This law can also show you why x^{0}=1 :\]
% formula_record_end

% formula_record_start
% formula_id: exponent-laws-html-038
% page_slug: exponent-laws-html
% page_title: Laws of Exponents
% page_url: https://www.mathsisfun.com/algebra/exponent-laws.html
% subject_slug: algebra
% subject_path: /algebra/
% section: The law that xm/n = n√xm = (n√x )m
% subsection: 
% formula_name: 
% source: formula-text-fallback
% index_global: 38
% index_in_page: 38
\[x^{1/n} = The n-th Root of x\]
% formula_record_end

% formula_record_start
% formula_id: exponent-laws-html-039
% page_slug: exponent-laws-html
% page_title: Laws of Exponents
% page_url: https://www.mathsisfun.com/algebra/exponent-laws.html
% subject_slug: algebra
% subject_path: /algebra/
% section: The law that xm/n = n√xm = (n√x )m
% subsection: 
% formula_name: 
% source: formula-text-fallback
% index_global: 39
% index_in_page: 39
\[And so a fractional exponent like 4^{\\frac{3}{2}} is really saying to do a cube (3) and a square root (\\frac{1}{2}), in any order, like this:\]
% formula_record_end

% formula_record_start
% formula_id: exponent-laws-html-040
% page_slug: exponent-laws-html
% page_title: Laws of Exponents
% page_url: https://www.mathsisfun.com/algebra/exponent-laws.html
% subject_slug: algebra
% subject_path: /algebra/
% section: The law that xm/n = n√xm = (n√x )m
% subsection: 
% formula_name: 
% source: formula-text-fallback
% index_global: 40
% index_in_page: 40
\[Just remember from fractions that m/n = m \times (1/n):\]
% formula_record_end

% formula_record_start
% formula_id: exponent-laws-html-041
% page_slug: exponent-laws-html
% page_title: Laws of Exponents
% page_url: https://www.mathsisfun.com/algebra/exponent-laws.html
% subject_slug: algebra
% subject_path: /algebra/
% section: The law that xm/n = n√xm = (n√x )m
% subsection: 
% formula_name: 
% source: formula-text-fallback
% index_global: 41
% index_in_page: 41
\[The order does not matter, so it also works for m/n = (1/n) \times m:\]
% formula_record_end

% formula_record_start
% formula_id: exponent-laws-html-042
% page_slug: exponent-laws-html
% page_title: Laws of Exponents
% page_url: https://www.mathsisfun.com/algebra/exponent-laws.html
% subject_slug: algebra
% subject_path: /algebra/
% section: And That Is It!
% subsection: 
% formula_name: 
% source: formula-text-fallback
% index_global: 42
% index_in_page: 42
\[A fractional exponent like 1/n means to take the nth root: x^{(1n)} = n\sqrt{}x\]
% formula_record_end

% formula_record_start
% formula_id: exponent-fractional-html-001
% page_slug: exponent-fractional-html
% page_title: Fractional Exponents
% page_url: https://www.mathsisfun.com/algebra/exponent-fractional.html
% subject_slug: algebra
% subject_path: /algebra/
% section: Try Another Fraction
% subsection: 
% formula_name: 
% source: formula-text-fallback
% index_global: 43
% index_in_page: 1
\[Let us try that again, but with an exponent of one-quarter (\\frac{1}{4}):\]
% formula_record_end

% formula_record_start
% formula_id: exponent-fractional-html-002
% page_slug: exponent-fractional-html
% page_title: Fractional Exponents
% page_url: https://www.mathsisfun.com/algebra/exponent-fractional.html
% subject_slug: algebra
% subject_path: /algebra/
% section: General Rule
% subsection: 
% formula_name: 
% source: formula-text-fallback
% index_global: 44
% index_in_page: 2
\[x^{1/n} = The n-th Root of x\]
% formula_record_end

% formula_record_start
% formula_id: exponent-fractional-html-003
% page_slug: exponent-fractional-html
% page_title: Fractional Exponents
% page_url: https://www.mathsisfun.com/algebra/exponent-fractional.html
% subject_slug: algebra
% subject_path: /algebra/
% section: What About More Complicated Fractions?
% subsection: 
% formula_name: 
% source: formula-text-fallback
% index_global: 45
% index_in_page: 3
\[What about a fractional exponent like 4^{\\frac{3}{2}} ?\]
% formula_record_end

% formula_record_start
% formula_id: exponent-fractional-html-004
% page_slug: exponent-fractional-html
% page_title: Fractional Exponents
% page_url: https://www.mathsisfun.com/algebra/exponent-fractional.html
% subject_slug: algebra
% subject_path: /algebra/
% section: What About More Complicated Fractions?
% subsection: 
% formula_name: 
% source: formula-text-fallback
% index_global: 46
% index_in_page: 4
\[That is really saying to do a cube (3) and a square root (\\frac{1}{2}), in any order.\]
% formula_record_end

% formula_record_start
% formula_id: exponent-fractional-html-005
% page_slug: exponent-fractional-html
% page_title: Fractional Exponents
% page_url: https://www.mathsisfun.com/algebra/exponent-fractional.html
% subject_slug: algebra
% subject_path: /algebra/
% section: What About More Complicated Fractions?
% subsection: 
% formula_name: 
% source: formula-text-fallback
% index_global: 47
% index_in_page: 5
\[a fraction (1/n) part\]
% formula_record_end

% formula_record_start
% formula_id: exponent-fractional-html-006
% page_slug: exponent-fractional-html
% page_title: Fractional Exponents
% page_url: https://www.mathsisfun.com/algebra/exponent-fractional.html
% subject_slug: algebra
% subject_path: /algebra/
% section: What About More Complicated Fractions?
% subsection: 
% formula_name: 
% source: formula-text-fallback
% index_global: 48
% index_in_page: 6
\[So, because m/n = m \times (1/n) we can do this:\]
% formula_record_end

% formula_record_start
% formula_id: exponent-fractional-html-007
% page_slug: exponent-fractional-html
% page_title: Fractional Exponents
% page_url: https://www.mathsisfun.com/algebra/exponent-fractional.html
% subject_slug: algebra
% subject_path: /algebra/
% section: What About More Complicated Fractions?
% subsection: 
% formula_name: 
% source: formula-text-fallback
% index_global: 49
% index_in_page: 7
\[x^{m/n} = x^{(m \times 1/n)} = (x^{m})^{1/n} = n\sqrt{}x^{m}\]
% formula_record_end

% formula_record_start
% formula_id: exponent-fractional-html-008
% page_slug: exponent-fractional-html
% page_title: Fractional Exponents
% page_url: https://www.mathsisfun.com/algebra/exponent-fractional.html
% subject_slug: algebra
% subject_path: /algebra/
% section: What About More Complicated Fractions?
% subsection: 
% formula_name: 
% source: formula-text-fallback
% index_global: 50
% index_in_page: 8
\[The order does not matter, so it also works for m/n = (1/n) \times m:\]
% formula_record_end

% formula_record_start
% formula_id: exponent-fractional-html-009
% page_slug: exponent-fractional-html
% page_title: Fractional Exponents
% page_url: https://www.mathsisfun.com/algebra/exponent-fractional.html
% subject_slug: algebra
% subject_path: /algebra/
% section: What About More Complicated Fractions?
% subsection: 
% formula_name: 
% source: formula-text-fallback
% index_global: 51
% index_in_page: 9
\[x^{m/n} = x^{(1/n \times m)} = (x^{1/n})^{m} = (n\sqrt{}x )^{m}\]
% formula_record_end

% formula_record_start
% formula_id: exponent-fractional-html-010
% page_slug: exponent-fractional-html
% page_title: Fractional Exponents
% page_url: https://www.mathsisfun.com/algebra/exponent-fractional.html
% subject_slug: algebra
% subject_path: /algebra/
% section: Now ... Play With The Graph!
% subsection: 
% formula_name: 
% source: formula-text-fallback
% index_global: 52
% index_in_page: 10
\[Start with m=1 and n=1, then slowly increase n so that you can see \\frac{1}{2}, \\frac{1}{3} and \\frac{1}{4}\]
% formula_record_end

% formula_record_start
% formula_id: exponent-fractional-html-011
% page_slug: exponent-fractional-html
% page_title: Fractional Exponents
% page_url: https://www.mathsisfun.com/algebra/exponent-fractional.html
% subject_slug: algebra
% subject_path: /algebra/
% section: Now ... Play With The Graph!
% subsection: 
% formula_name: 
% source: formula-text-fallback
% index_global: 53
% index_in_page: 11
\[Then try m=2 and slide n up and down to see fractions like \\frac{2}{3} and so on\]
% formula_record_end

% formula_record_start
% formula_id: nth-root-html-001
% page_slug: nth-root-html
% page_title: nth Root
% page_url: https://www.mathsisfun.com/numbers/nth-root.html
% subject_slug: numbers
% subject_path: /numbers/
% section: Laws of Exponents
% subsection: Addition and Subtraction
% formula_name: 
% source: mathsisfun-large-span
% index_global: 54
% index_in_page: 1
\[a^{2} + b^{2} = c^{2}\]
% formula_record_end

% formula_record_start
% formula_id: nth-root-html-002
% page_slug: nth-root-html
% page_title: nth Root
% page_url: https://www.mathsisfun.com/numbers/nth-root.html
% subject_slug: numbers
% subject_path: /numbers/
% section: Even Roots of Negative Values
% subsection: 
% formula_name: 
% source: formula-text-fallback
% index_global: 55
% index_in_page: 2
\[Example: \sqrt{}-4 = 2i, where i is the unit imaginary number.\]
% formula_record_end

% formula_record_start
% formula_id: nth-root-html-003
% page_slug: nth-root-html
% page_title: nth Root
% page_url: https://www.mathsisfun.com/numbers/nth-root.html
% subject_slug: numbers
% subject_path: /numbers/
% section: Properties
% subsection: Multiplication and Division
% formula_name: 
% source: formula-text-fallback
% index_global: 56
% index_in_page: 3
\[n\sqrt{}ab = n\sqrt{}a \times n\sqrt{}b\]
% formula_record_end

% formula_record_start
% formula_id: nth-root-html-004
% page_slug: nth-root-html
% page_title: nth Root
% page_url: https://www.mathsisfun.com/numbers/nth-root.html
% subject_slug: numbers
% subject_path: /numbers/
% section: Properties
% subsection: Multiplication and Division
% formula_name: 
% source: formula-text-fallback
% index_global: 57
% index_in_page: 4
\[n\sqrt{}a/b = n\sqrt{}a / n\sqrt{}b\]
% formula_record_end

% formula_record_start
% formula_id: nth-root-html-005
% page_slug: nth-root-html
% page_title: nth Root
% page_url: https://www.mathsisfun.com/numbers/nth-root.html
% subject_slug: numbers
% subject_path: /numbers/
% section: Properties
% subsection: Addition and Subtraction
% formula_name: 
% source: formula-text-fallback
% index_global: 58
% index_in_page: 5
\[n\sqrt{}a^{n} + b^{n} \ne a + b\]
% formula_record_end

% formula_record_start
% formula_id: nth-root-html-006
% page_slug: nth-root-html
% page_title: nth Root
% page_url: https://www.mathsisfun.com/numbers/nth-root.html
% subject_slug: numbers
% subject_path: /numbers/
% section: Properties
% subsection: Exponents vs Roots
% formula_name: 
% source: formula-text-fallback
% index_global: 59
% index_in_page: 6
\[An exponent on one side of "=" can be turned into a root on the other side of "=":\]
% formula_record_end

% formula_record_start
% formula_id: nth-root-html-007
% page_slug: nth-root-html
% page_title: nth Root
% page_url: https://www.mathsisfun.com/numbers/nth-root.html
% subject_slug: numbers
% subject_path: /numbers/
% section: Properties
% subsection: Exponents vs Roots
% formula_name: 
% source: formula-text-fallback
% index_global: 60
% index_in_page: 7
\[If a^{n} = b then a = n\sqrt{}b\]
% formula_record_end

% formula_record_start
% formula_id: nth-root-html-008
% page_slug: nth-root-html
% page_title: nth Root
% page_url: https://www.mathsisfun.com/numbers/nth-root.html
% subject_slug: numbers
% subject_path: /numbers/
% section: Properties
% subsection: nth Root of a-to-the-mth-Power
% formula_name: 
% source: formula-text-fallback
% index_global: 61
% index_in_page: 8
\[n\sqrt{}a^{m} = (n\sqrt{}a )^{m}\]
% formula_record_end

% formula_record_start
% formula_id: nth-root-html-009
% page_slug: nth-root-html
% page_title: nth Root
% page_url: https://www.mathsisfun.com/numbers/nth-root.html
% subject_slug: numbers
% subject_path: /numbers/
% section: Properties
% subsection: nth Root of a-to-the-mth-Power
% formula_name: 
% source: formula-text-fallback
% index_global: 62
% index_in_page: 9
\[n\sqrt{}a^{m} = a^{mn}\]
% formula_record_end

% formula_record_start
% formula_id: nth-root-html-010
% page_slug: nth-root-html
% page_title: nth Root
% page_url: https://www.mathsisfun.com/numbers/nth-root.html
% subject_slug: numbers
% subject_path: /numbers/
% section: Properties
% subsection: nth Root of a-to-the-mth-Power
% formula_name: 
% source: formula-text-fallback
% index_global: 63
% index_in_page: 10
\[This works because the nth root is the same as an exponent of (1/n)\]
% formula_record_end

% formula_record_start
% formula_id: nth-root-html-011
% page_slug: nth-root-html
% page_title: nth Root
% page_url: https://www.mathsisfun.com/numbers/nth-root.html
% subject_slug: numbers
% subject_path: /numbers/
% section: Properties
% subsection: nth Root of a-to-the-mth-Power
% formula_name: 
% source: formula-text-fallback
% index_global: 64
% index_in_page: 11
\[n\sqrt{}a = a^{1n}\]
% formula_record_end

% formula_record_start
% formula_id: complex-numbers-html-001
% page_slug: complex-numbers-html
% page_title: Complex Numbers
% page_url: https://www.mathsisfun.com/numbers/complex-numbers.html
% subject_slug: numbers
% subject_path: /numbers/
% section: Laws of Exponents
% subsection: 
% formula_name: 
% source: mathsisfun-large-span
% index_global: 65
% index_in_page: 1
\[i^{2 }= -1\]
% formula_record_end

% formula_record_start
% formula_id: complex-numbers-html-002
% page_slug: complex-numbers-html
% page_title: Complex Numbers
% page_url: https://www.mathsisfun.com/numbers/complex-numbers.html
% subject_slug: numbers
% subject_path: /numbers/
% section: Properties
% subsection: 
% formula_name: 
% source: formula-text-fallback
% index_global: 66
% index_in_page: 2
\[z = a + bi\]
% formula_record_end

% formula_record_start
% formula_id: complex-numbers-html-003
% page_slug: complex-numbers-html
% page_title: Complex Numbers
% page_url: https://www.mathsisfun.com/numbers/complex-numbers.html
% subject_slug: numbers
% subject_path: /numbers/
% section: Properties
% subsection: 
% formula_name: 
% source: formula-text-fallback
% index_global: 67
% index_in_page: 3
\[i is the unit imaginary number = \sqrt{}-1\]
% formula_record_end

% formula_record_start
% formula_id: complex-numbers-html-004
% page_slug: complex-numbers-html
% page_title: Complex Numbers
% page_url: https://www.mathsisfun.com/numbers/complex-numbers.html
% subject_slug: numbers
% subject_path: /numbers/
% section: Properties
% subsection: 
% formula_name: 
% source: formula-text-fallback
% index_global: 68
% index_in_page: 4
\[Re(z) = a\]
% formula_record_end

% formula_record_start
% formula_id: complex-numbers-html-005
% page_slug: complex-numbers-html
% page_title: Complex Numbers
% page_url: https://www.mathsisfun.com/numbers/complex-numbers.html
% subject_slug: numbers
% subject_path: /numbers/
% section: Properties
% subsection: 
% formula_name: 
% source: formula-text-fallback
% index_global: 69
% index_in_page: 5
\[Im(z) = b\]
% formula_record_end

% formula_record_start
% formula_id: complex-numbers-html-006
% page_slug: complex-numbers-html
% page_title: Complex Numbers
% page_url: https://www.mathsisfun.com/numbers/complex-numbers.html
% subject_slug: numbers
% subject_path: /numbers/
% section: Properties
% subsection: 
% formula_name: 
% source: formula-text-fallback
% index_global: 70
% index_in_page: 6
\[z^{\\cdot } = a - bi\]
% formula_record_end

% formula_record_start
% formula_id: complex-numbers-html-007
% page_slug: complex-numbers-html
% page_title: Complex Numbers
% page_url: https://www.mathsisfun.com/numbers/complex-numbers.html
% subject_slug: numbers
% subject_path: /numbers/
% section: Properties
% subsection: 
% formula_name: 
% source: formula-text-fallback
% index_global: 71
% index_in_page: 7
\[|z| = \sqrt{}(a^{2} + b^{2})\]
% formula_record_end

% formula_record_start
% formula_id: complex-numbers-html-008
% page_slug: complex-numbers-html
% page_title: Complex Numbers
% page_url: https://www.mathsisfun.com/numbers/complex-numbers.html
% subject_slug: numbers
% subject_path: /numbers/
% section: Properties
% subsection: 
% formula_name: 
% source: formula-text-fallback
% index_global: 72
% index_in_page: 8
\[Arg(z) = tan^{-1}(b/a)\]
% formula_record_end

% formula_record_start
% formula_id: complex-numbers-html-009
% page_slug: complex-numbers-html
% page_title: Complex Numbers
% page_url: https://www.mathsisfun.com/numbers/complex-numbers.html
% subject_slug: numbers
% subject_path: /numbers/
% section: Adding
% subsection: 
% formula_name: 
% source: formula-text-fallback
% index_global: 73
% index_in_page: 9
\[(a+bi) + (c+di) = (a+c) + (b+d)i\]
% formula_record_end

% formula_record_start
% formula_id: complex-numbers-html-010
% page_slug: complex-numbers-html
% page_title: Complex Numbers
% page_url: https://www.mathsisfun.com/numbers/complex-numbers.html
% subject_slug: numbers
% subject_path: /numbers/
% section: Multiplying
% subsection: But There is a Quicker Way!
% formula_name: 
% source: formula-text-fallback
% index_global: 74
% index_in_page: 10
\[(a+bi)(c+di) = (ac-bd) + (ad+bc)i\]
% formula_record_end

% formula_record_start
% formula_id: complex-numbers-html-011
% page_slug: complex-numbers-html
% page_title: Complex Numbers
% page_url: https://www.mathsisfun.com/numbers/complex-numbers.html
% subject_slug: numbers
% subject_path: /numbers/
% section: Multiplying
% subsection: Let us try i2
% formula_name: 
% source: formula-text-fallback
% index_global: 75
% index_in_page: 11
\[Just for fun, let's use the method to calculate i^{2}\]
% formula_record_end

% formula_record_start
% formula_id: complex-numbers-html-012
% page_slug: complex-numbers-html
% page_title: Complex Numbers
% page_url: https://www.mathsisfun.com/numbers/complex-numbers.html
% subject_slug: numbers
% subject_path: /numbers/
% section: Multiplying
% subsection: Let us try i2
% formula_name: 
% source: formula-text-fallback
% index_global: 76
% index_in_page: 12
\[And that agrees nicely with the definition that i^{2 }= -1\]
% formula_record_end

% formula_record_start
% formula_id: complex-numbers-html-013
% page_slug: complex-numbers-html
% page_title: Complex Numbers
% page_url: https://www.mathsisfun.com/numbers/complex-numbers.html
% subject_slug: numbers
% subject_path: /numbers/
% section: Multiplying By the Conjugate
% subsection: 
% formula_name: 
% source: formula-text-fallback
% index_global: 77
% index_in_page: 13
\[(4 - 5i)(4 + 5i) = 16 + 20i - 20i - 25i^{2}\]
% formula_record_end

% formula_record_start
% formula_id: complex-numbers-html-014
% page_slug: complex-numbers-html
% page_title: Complex Numbers
% page_url: https://www.mathsisfun.com/numbers/complex-numbers.html
% subject_slug: numbers
% subject_path: /numbers/
% section: Multiplying By the Conjugate
% subsection: 
% formula_name: 
% source: formula-text-fallback
% index_global: 78
% index_in_page: 14
\[(4 - 5i)(4 + 5i) = 16 - 25i^{2}\]
% formula_record_end

% formula_record_start
% formula_id: complex-numbers-html-015
% page_slug: complex-numbers-html
% page_title: Complex Numbers
% page_url: https://www.mathsisfun.com/numbers/complex-numbers.html
% subject_slug: numbers
% subject_path: /numbers/
% section: Multiplying By the Conjugate
% subsection: 
% formula_name: 
% source: formula-text-fallback
% index_global: 79
% index_in_page: 15
\[Also i^{2} = -1 :\]
% formula_record_end

% formula_record_start
% formula_id: complex-numbers-html-016
% page_slug: complex-numbers-html
% page_title: Complex Numbers
% page_url: https://www.mathsisfun.com/numbers/complex-numbers.html
% subject_slug: numbers
% subject_path: /numbers/
% section: Multiplying By the Conjugate
% subsection: 
% formula_name: 
% source: formula-text-fallback
% index_global: 80
% index_in_page: 16
\[(4 - 5i)(4 + 5i) = 16 + 25\]
% formula_record_end

% formula_record_start
% formula_id: complex-numbers-html-017
% page_slug: complex-numbers-html
% page_title: Complex Numbers
% page_url: https://www.mathsisfun.com/numbers/complex-numbers.html
% subject_slug: numbers
% subject_path: /numbers/
% section: Multiplying By the Conjugate
% subsection: 
% formula_name: 
% source: formula-text-fallback
% index_global: 81
% index_in_page: 17
\[(4 - 5i)(4 + 5i) = 4^{2} + 5^{2}\]
% formula_record_end

% formula_record_start
% formula_id: dividing-by-zero-html-001
% page_slug: dividing-by-zero-html
% page_title: Dividing by Zero
% page_url: https://www.mathsisfun.com/numbers/dividing-by-zero.html
% subject_slug: numbers
% subject_path: /numbers/
% section: So what is 0/0 ?
% subsection: 
% formula_name: 
% source: formula-text-fallback
% index_global: 82
% index_in_page: 1
\[\\frac{0}{0} is like asking "how many 0s in 0?"\]
% formula_record_end

% formula_record_start
% formula_id: dividing-by-zero-html-002
% page_slug: dividing-by-zero-html
% page_title: Dividing by Zero
% page_url: https://www.mathsisfun.com/numbers/dividing-by-zero.html
% subject_slug: numbers
% subject_path: /numbers/
% section: So what is 0/0 ?
% subsection: 
% formula_name: 
% source: formula-text-fallback
% index_global: 83
% index_in_page: 2
\[So \\frac{0}{0} is indeterminate (it could be any value).\]
% formula_record_end

% formula_record_start
% formula_id: exponent-html-001
% page_slug: exponent-html
% page_title: Exponents
% page_url: https://www.mathsisfun.com/exponent.html
% subject_slug: exponent.html
% subject_path: /exponent.html/
% section: Another Way of Writing It
% subsection: 
% formula_name: 
% source: mathsisfun-large-span
% index_global: 84
% index_in_page: 1
\[2^{4} = 2 \times 2 \times 2 \times 2 = 16\]
% formula_record_end

% formula_record_start
% formula_id: exponent-html-002
% page_slug: exponent-html
% page_title: Exponents
% page_url: https://www.mathsisfun.com/exponent.html
% subject_slug: exponent.html
% subject_path: /exponent.html/
% section: Uncategorized
% subsection: 
% formula_name: 
% source: formula-text-fallback
% index_global: 85
% index_in_page: 2
\[In 8^{2} the "2" says to use 8 twice in a multiplication,\]
% formula_record_end

% formula_record_start
% formula_id: exponent-html-003
% page_slug: exponent-html
% page_title: Exponents
% page_url: https://www.mathsisfun.com/exponent.html
% subject_slug: exponent.html
% subject_path: /exponent.html/
% section: Uncategorized
% subsection: 
% formula_name: 
% source: formula-text-fallback
% index_global: 86
% index_in_page: 3
\[so 8^{2} = 8 \times 8 = 64\]
% formula_record_end

% formula_record_start
% formula_id: exponent-html-004
% page_slug: exponent-html
% page_title: Exponents
% page_url: https://www.mathsisfun.com/exponent.html
% subject_slug: exponent.html
% subject_path: /exponent.html/
% section: Uncategorized
% subsection: 
% formula_name: 
% source: formula-text-fallback
% index_global: 87
% index_in_page: 4
\[In words: 8^{2} could be called "8 to the power 2" or "8 to the second power", or simply "8 squared"\]
% formula_record_end

% formula_record_start
% formula_id: exponent-html-005
% page_slug: exponent-html
% page_title: Exponents
% page_url: https://www.mathsisfun.com/exponent.html
% subject_slug: exponent.html
% subject_path: /exponent.html/
% section: Uncategorized
% subsection: 
% formula_name: 
% source: formula-text-fallback
% index_global: 88
% index_in_page: 5
\[a^{n} tells you to multiply a by itself,so there are n of those a's:\]
% formula_record_end

% formula_record_start
% formula_id: exponent-html-006
% page_slug: exponent-html
% page_title: Exponents
% page_url: https://www.mathsisfun.com/exponent.html
% subject_slug: exponent.html
% subject_path: /exponent.html/
% section: Another Way of Writing It
% subsection: 
% formula_name: 
% source: formula-text-fallback
% index_global: 89
% index_in_page: 6
\[Sometimes people use the ^ symbol (above the 6 on your keyboard), as it is easy to type.\]
% formula_record_end

% formula_record_start
% formula_id: index-notation-powers-html-001
% page_slug: index-notation-powers-html
% page_title: Index Notation - Powers of 10
% page_url: https://www.mathsisfun.com/index-notation-powers.html
% subject_slug: index-notation-powers.html
% subject_path: /index-notation-powers.html/
% section: Another Way of Writing It
% subsection: 
% formula_name: 
% source: formula-text-fallback
% index_global: 90
% index_in_page: 1
\[Sometimes people use the ^ symbol (above the 6 on your keyboard), as it is easy to type.\]
% formula_record_end

% formula_record_start
% formula_id: index-notation-powers-html-002
% page_slug: index-notation-powers-html
% page_title: Index Notation - Powers of 10
% page_url: https://www.mathsisfun.com/index-notation-powers.html
% subject_slug: index-notation-powers.html
% subject_path: /index-notation-powers.html/
% section: Another Way of Writing It
% subsection: 
% formula_name: 
% source: formula-text-fallback
% index_global: 91
% index_in_page: 2
\[Most calculators use "E" or "e", so 6E+5 means 6 \times 10^{5}. "E" stands for exponent.\]
% formula_record_end

% formula_record_start
% formula_id: calculus-index-html-001
% page_slug: calculus-index-html
% page_title: Calculus
% page_url: https://www.mathsisfun.com/calculus/index.html
% subject_slug: calculus
% subject_path: /calculus/
% section: Differential Equations
% subsection: 
% formula_name: 
% source: formula-text-fallback
% index_global: 92
% index_in_page: 1
\[= 5x" height="123" width="188">\]
% formula_record_end

% formula_record_start
% formula_id: calculator-html-001
% page_slug: calculator-html
% page_title: Calculator
% page_url: https://www.mathsisfun.com/numbers/calculator.html
% subject_slug: numbers
% subject_path: /numbers/
% section: Uncategorized
% subsection: 
% formula_name: 
% source: formula-text-fallback
% index_global: 93
% index_in_page: 1
\[Use the buttons, or type formulas like 3\\cdot \\sqrt{2} or pi/2\]
% formula_record_end

% formula_record_start
% formula_id: formula-sheet-001
% page_slug: formula-sheet
% page_title: Formula Sheet
% page_url: https://formulasheet.com/#q|l|1228
% subject_slug: formulasheet
% subject_path: /formulasheet/
% section: Area (geometry)
% subsection: Area of a circle
% formula_name: 
% source: formulasheet-latex
% index_global: 94
% index_in_page: 1
\[\begin{equation} A = \pi r^2 = \frac{\pi d^2}{4} \end{equation}\]
% formula_record_end

% formula_record_start
% formula_id: formula-sheet-002
% page_slug: formula-sheet
% page_title: Formula Sheet
% page_url: https://formulasheet.com/#q|l|1228
% subject_slug: formulasheet
% subject_path: /formulasheet/
% section: Area (geometry)
% subsection: Area of a parallelogram
% formula_name: 
% source: formulasheet-latex
% index_global: 95
% index_in_page: 2
\[\begin{equation} A = b h \end{equation}\]
% formula_record_end

% formula_record_start
% formula_id: formula-sheet-003
% page_slug: formula-sheet
% page_title: Formula Sheet
% page_url: https://formulasheet.com/#q|l|1228
% subject_slug: formulasheet
% subject_path: /formulasheet/
% section: Area (geometry)
% subsection: Area of a rectangle
% formula_name: 
% source: formulasheet-latex
% index_global: 96
% index_in_page: 3
\[\begin{equation} A = l w \end{equation}\]
% formula_record_end

% formula_record_start
% formula_id: formula-sheet-004
% page_slug: formula-sheet
% page_title: Formula Sheet
% page_url: https://formulasheet.com/#q|l|1228
% subject_slug: formulasheet
% subject_path: /formulasheet/
% section: Area (geometry)
% subsection: Area of a regular hexagon
% formula_name: 
% source: formulasheet-latex
% index_global: 97
% index_in_page: 4
\[\begin{equation} A = \frac{3 \sqrt{3}}{2} \, s^2 \approx 2.598076211 s^2 \end{equation}\]
% formula_record_end

% formula_record_start
% formula_id: formula-sheet-005
% page_slug: formula-sheet
% page_title: Formula Sheet
% page_url: https://formulasheet.com/#q|l|1228
% subject_slug: formulasheet
% subject_path: /formulasheet/
% section: Area (geometry)
% subsection: Area of a regular octagon
% formula_name: 
% source: formulasheet-latex
% index_global: 98
% index_in_page: 5
\[\begin{equation} A = 2 \left( 1 + \sqrt{2} \right) s^2 \approx 4.828427125 s^2 \end{equation}\]
% formula_record_end

% formula_record_start
% formula_id: formula-sheet-006
% page_slug: formula-sheet
% page_title: Formula Sheet
% page_url: https://formulasheet.com/#q|l|1228
% subject_slug: formulasheet
% subject_path: /formulasheet/
% section: Area (geometry)
% subsection: Area of a regular polygon
% formula_name: 
% source: formulasheet-latex
% index_global: 99
% index_in_page: 6
\[\begin{equation} A = \frac{n l^2 }{4} \cot \left( \frac{\pi}{n} \right) = \frac{p^2}{4n} \cot \left( \frac{\pi}{n} \right) \end{equation}\]
% formula_record_end

% formula_record_start
% formula_id: formula-sheet-007
% page_slug: formula-sheet
% page_title: Formula Sheet
% page_url: https://formulasheet.com/#q|l|1228
% subject_slug: formulasheet
% subject_path: /formulasheet/
% section: Area (geometry)
% subsection: Area of a square
% formula_name: 
% source: formulasheet-latex
% index_global: 100
% index_in_page: 7
\[\begin{equation} A = s^2 \end{equation}\]
% formula_record_end

% formula_record_start
% formula_id: formula-sheet-008
% page_slug: formula-sheet
% page_title: Formula Sheet
% page_url: https://formulasheet.com/#q|l|1228
% subject_slug: formulasheet
% subject_path: /formulasheet/
% section: Area (geometry)
% subsection: Area of a trapezoid
% formula_name: 
% source: formulasheet-latex
% index_global: 101
% index_in_page: 8
\[\begin{equation} A = \frac{a + b}{2} \, h \end{equation}\]
% formula_record_end

% formula_record_start
% formula_id: formula-sheet-009
% page_slug: formula-sheet
% page_title: Formula Sheet
% page_url: https://formulasheet.com/#q|l|1228
% subject_slug: formulasheet
% subject_path: /formulasheet/
% section: Area (geometry)
% subsection: Area of a triangle
% formula_name: 
% source: formulasheet-latex
% index_global: 102
% index_in_page: 9
\[\begin{equation} A = \frac{1}{2} b h \end{equation}\]
% formula_record_end

% formula_record_start
% formula_id: formula-sheet-010
% page_slug: formula-sheet
% page_title: Formula Sheet
% page_url: https://formulasheet.com/#q|l|1228
% subject_slug: formulasheet
% subject_path: /formulasheet/
% section: Area (geometry)
% subsection: Area of an ellipse
% formula_name: 
% source: formulasheet-latex
% index_global: 103
% index_in_page: 10
\[\begin{equation} A = \pi a b \end{equation}\]
% formula_record_end

% formula_record_start
% formula_id: formula-sheet-011
% page_slug: formula-sheet
% page_title: Formula Sheet
% page_url: https://formulasheet.com/#q|l|1228
% subject_slug: formulasheet
% subject_path: /formulasheet/
% section: Area (geometry)
% subsection: Surface area of a cone
% formula_name: 
% source: formulasheet-latex
% index_global: 104
% index_in_page: 11
\[\begin{equation} A = \pi r \left( r + s \right) = \pi r \left( r + \sqrt{r^2 + h^2} \right) \end{equation}\]
% formula_record_end

% formula_record_start
% formula_id: formula-sheet-012
% page_slug: formula-sheet
% page_title: Formula Sheet
% page_url: https://formulasheet.com/#q|l|1228
% subject_slug: formulasheet
% subject_path: /formulasheet/
% section: Area (geometry)
% subsection: Surface area of a cube
% formula_name: 
% source: formulasheet-latex
% index_global: 105
% index_in_page: 12
\[\begin{equation} A = 6 s^2 \end{equation}\]
% formula_record_end

% formula_record_start
% formula_id: formula-sheet-013
% page_slug: formula-sheet
% page_title: Formula Sheet
% page_url: https://formulasheet.com/#q|l|1228
% subject_slug: formulasheet
% subject_path: /formulasheet/
% section: Area (geometry)
% subsection: Surface area of a cylinder
% formula_name: 
% source: formulasheet-latex
% index_global: 106
% index_in_page: 13
\[\begin{equation} A = 2 \pi r \left( r + h \right) \end{equation}\]
% formula_record_end

% formula_record_start
% formula_id: formula-sheet-014
% page_slug: formula-sheet
% page_title: Formula Sheet
% page_url: https://formulasheet.com/#q|l|1228
% subject_slug: formulasheet
% subject_path: /formulasheet/
% section: Area (geometry)
% subsection: Surface area of a rectangular prism
% formula_name: 
% source: formulasheet-latex
% index_global: 107
% index_in_page: 14
\[\begin{equation} A = 2 \left( l w + l h + w h \right) \end{equation}\]
% formula_record_end

% formula_record_start
% formula_id: formula-sheet-015
% page_slug: formula-sheet
% page_title: Formula Sheet
% page_url: https://formulasheet.com/#q|l|1228
% subject_slug: formulasheet
% subject_path: /formulasheet/
% section: Area (geometry)
% subsection: Surface area of a regular tetrahedron
% formula_name: 
% source: formulasheet-latex
% index_global: 108
% index_in_page: 15
\[\begin{equation} A = \sqrt{3} s^2 \end{equation}\]
% formula_record_end

% formula_record_start
% formula_id: formula-sheet-016
% page_slug: formula-sheet
% page_title: Formula Sheet
% page_url: https://formulasheet.com/#q|l|1228
% subject_slug: formulasheet
% subject_path: /formulasheet/
% section: Area (geometry)
% subsection: Surface area of a sphere
% formula_name: 
% source: formulasheet-latex
% index_global: 109
% index_in_page: 16
\[\begin{equation} A = 4 \pi r^2 = \pi d^2 \end{equation}\]
% formula_record_end

% formula_record_start
% formula_id: formula-sheet-2-001
% page_slug: formula-sheet-2
% page_title: Formula Sheet
% page_url: https://formulasheet.com/#q|l|1255
% subject_slug: formulasheet
% subject_path: /formulasheet/
% section: Calculus
% subsection: Chain rule
% formula_name: 
% source: formulasheet-latex
% index_global: 110
% index_in_page: 1
\[\begin{equation} (f \circ g)' (x) = f'(g) \, g'(x) \end{equation}\]
% formula_record_end

% formula_record_start
% formula_id: formula-sheet-2-002
% page_slug: formula-sheet-2
% page_title: Formula Sheet
% page_url: https://formulasheet.com/#q|l|1255
% subject_slug: formulasheet
% subject_path: /formulasheet/
% section: Calculus
% subsection: Chain rule (Leibniz notation)
% formula_name: 
% source: formulasheet-latex
% index_global: 111
% index_in_page: 2
\[\begin{equation} \frac {du}{dx} = \frac {du}{dv} \, \frac {dv}{dx} \end{equation}\]
% formula_record_end

% formula_record_start
% formula_id: formula-sheet-2-003
% page_slug: formula-sheet-2
% page_title: Formula Sheet
% page_url: https://formulasheet.com/#q|l|1255
% subject_slug: formulasheet
% subject_path: /formulasheet/
% section: Calculus
% subsection: Derivative definition
% formula_name: 
% source: formulasheet-latex
% index_global: 112
% index_in_page: 3
\[\begin{equation} f'(x\!=\!a)=\lim_{h \to 0}\frac{f(a + h) - f(a)}{h} \end{equation}\]
% formula_record_end

% formula_record_start
% formula_id: formula-sheet-2-004
% page_slug: formula-sheet-2
% page_title: Formula Sheet
% page_url: https://formulasheet.com/#q|l|1255
% subject_slug: formulasheet
% subject_path: /formulasheet/
% section: Calculus
% subsection: Inverse power rule
% formula_name: 
% source: formulasheet-latex
% index_global: 113
% index_in_page: 4
\[\begin{equation} \int x^n dx = \frac{x^{n+1}}{n+1} + C \qquad\qquad n \neq -1 \end{equation}\]
% formula_record_end

% formula_record_start
% formula_id: formula-sheet-2-005
% page_slug: formula-sheet-2
% page_title: Formula Sheet
% page_url: https://formulasheet.com/#q|l|1255
% subject_slug: formulasheet
% subject_path: /formulasheet/
% section: Calculus
% subsection: Linearity of differentiation
% formula_name: 
% source: formulasheet-latex
% index_global: 114
% index_in_page: 5
\[\begin{align} f(x) &= a \, g(x) + b \, h(x) \\\ f'(x) &= a \, g'(x) + b \, h'(x) \end{align}\]
% formula_record_end

% formula_record_start
% formula_id: formula-sheet-2-006
% page_slug: formula-sheet-2
% page_title: Formula Sheet
% page_url: https://formulasheet.com/#q|l|1255
% subject_slug: formulasheet
% subject_path: /formulasheet/
% section: Calculus
% subsection: Linearity of differentiation (Leibniz notation)
% formula_name: 
% source: formulasheet-latex
% index_global: 115
% index_in_page: 6
\[\begin{equation} \frac{d(a \, u+ b \, v)}{dx} = a \, \frac{du}{dx} +b\frac{dv}{dx} \end{equation}\]
% formula_record_end

% formula_record_start
% formula_id: formula-sheet-2-007
% page_slug: formula-sheet-2
% page_title: Formula Sheet
% page_url: https://formulasheet.com/#q|l|1255
% subject_slug: formulasheet
% subject_path: /formulasheet/
% section: Calculus
% subsection: Linearity of integration
% formula_name: 
% source: formulasheet-latex
% index_global: 116
% index_in_page: 7
\[\begin{equation} \int (a \, g(x) + b \, h(x))\, dx = a \int g(x) \, dx + b \int h(x) \, dx \end{equation}\]
% formula_record_end

% formula_record_start
% formula_id: formula-sheet-2-008
% page_slug: formula-sheet-2
% page_title: Formula Sheet
% page_url: https://formulasheet.com/#q|l|1255
% subject_slug: formulasheet
% subject_path: /formulasheet/
% section: Calculus
% subsection: Power rule
% formula_name: 
% source: formulasheet-latex
% index_global: 117
% index_in_page: 8
\[\begin{equation} \frac{d}{dx} x^n = n x^{n-1} \qquad\qquad n \neq 0 \end{equation}\]
% formula_record_end

% formula_record_start
% formula_id: formula-sheet-2-009
% page_slug: formula-sheet-2
% page_title: Formula Sheet
% page_url: https://formulasheet.com/#q|l|1255
% subject_slug: formulasheet
% subject_path: /formulasheet/
% section: Calculus
% subsection: Product rule
% formula_name: 
% source: formulasheet-latex
% index_global: 118
% index_in_page: 9
\[\begin{equation} \left[ f(x) \, g(x) \right]'=f'(x) \, g(x) + f(x) \, g'(x) \end{equation}\]
% formula_record_end

% formula_record_start
% formula_id: formula-sheet-2-010
% page_slug: formula-sheet-2
% page_title: Formula Sheet
% page_url: https://formulasheet.com/#q|l|1255
% subject_slug: formulasheet
% subject_path: /formulasheet/
% section: Calculus
% subsection: Product rule (Leibniz notation)
% formula_name: 
% source: formulasheet-latex
% index_global: 119
% index_in_page: 10
\[\begin{equation} \dfrac{d}{dx}(u \, v) = \dfrac{du}{dx} \, v + u \, \dfrac{dv}{dx} \end{equation}\]
% formula_record_end

% formula_record_start
% formula_id: formula-sheet-2-011
% page_slug: formula-sheet-2
% page_title: Formula Sheet
% page_url: https://formulasheet.com/#q|l|1255
% subject_slug: formulasheet
% subject_path: /formulasheet/
% section: Calculus
% subsection: Quotient rule
% formula_name: 
% source: formulasheet-latex
% index_global: 120
% index_in_page: 11
\[\begin{equation} f'(x) = \left[ \frac{g(x)}{h(x)} \right]' = \frac{g'(x) \, h(x) - g(x) \, h'(x)}{[h(x)]^2} \end{equation}\]
% formula_record_end

% formula_record_start
% formula_id: formula-sheet-2-012
% page_slug: formula-sheet-2
% page_title: Formula Sheet
% page_url: https://formulasheet.com/#q|l|1255
% subject_slug: formulasheet
% subject_path: /formulasheet/
% section: Calculus
% subsection: Riemann sum
% formula_name: 
% source: formulasheet-latex
% index_global: 121
% index_in_page: 12
\[\begin{align} &S = \sum_{i=1}^{n} f(x_i^*)(x_{i} - x_{i-1}) \qquad x_{i-1}\le x_i^* \le x_i \\\ &\text{Left Riemann sum: } x_i^* = x_{i-1} \text{ for all } i \\\ &\text{Right Riemann sum: } x_i^* = x_{i} \text{ for all } i \\\ &\text{Middle Riemann sum: } x_i^* = \frac{1}{2} \left( x_i + x_{i-1} \right) \text{ for all } i \end{align}\]
% formula_record_end

% formula_record_start
% formula_id: formula-sheet-2-013
% page_slug: formula-sheet-2
% page_title: Formula Sheet
% page_url: https://formulasheet.com/#q|l|1255
% subject_slug: formulasheet
% subject_path: /formulasheet/
% section: Calculus
% subsection: Second fundamental theorem of calculus
% formula_name: 
% source: formulasheet-latex
% index_global: 122
% index_in_page: 13
\[\begin{equation} \int_a^b f(x) \, dx = F(b) - F(a) \end{equation}\]
% formula_record_end

% formula_record_start
% formula_id: formula-sheet-2-014
% page_slug: formula-sheet-2
% page_title: Formula Sheet
% page_url: https://formulasheet.com/#q|l|1255
% subject_slug: formulasheet
% subject_path: /formulasheet/
% section: Calculus
% subsection: Simpson's rule
% formula_name: 
% source: formulasheet-latex
% index_global: 123
% index_in_page: 14
\[\begin{equation} \int_a^b f(x) \, dx \approx \frac{b-a}{6} \left[ f(a) + 4f \left( \frac{a+b}{2} \right) + f(b) \right] \end{equation}\]
% formula_record_end

% formula_record_start
% formula_id: formula-sheet-3-001
% page_slug: formula-sheet-3
% page_title: Formula Sheet
% page_url: https://formulasheet.com/#q|l|6084
% subject_slug: formulasheet
% subject_path: /formulasheet/
% section: Chapter 12
% subsection: 
% formula_name: 
% source: formulasheet-latex
% index_global: 124
% index_in_page: 1
\[\begin{equation} \label{Arithmetic Sum for Finite} S_{n}=\frac{n}{2}(a_{1}+a_{n}) \end{equation}\]
% formula_record_end

% formula_record_start
% formula_id: formula-sheet-3-002
% page_slug: formula-sheet-3
% page_title: Formula Sheet
% page_url: https://formulasheet.com/#q|l|6084
% subject_slug: formulasheet
% subject_path: /formulasheet/
% section: Chapter 12
% subsection: 
% formula_name: 
% source: formulasheet-latex
% index_global: 125
% index_in_page: 2
\[\begin{equation} \label{Geometric Sum of a Finite} S_{n}=\frac{a_{1}(1-r^{n})}{1-r} \end{equation}\]
% formula_record_end

% formula_record_start
% formula_id: formula-sheet-3-003
% page_slug: formula-sheet-3
% page_title: Formula Sheet
% page_url: https://formulasheet.com/#q|l|6084
% subject_slug: formulasheet
% subject_path: /formulasheet/
% section: Chapter 12
% subsection: 
% formula_name: 
% source: formulasheet-latex
% index_global: 126
% index_in_page: 3
\[\begin{equation} \label{Geometric Sum of a Finite} S_{n}=\frac{a_{1}-r\cdot a_{n}}{1-r} \end{equation}\]
% formula_record_end

% formula_record_start
% formula_id: formula-sheet-3-004
% page_slug: formula-sheet-3
% page_title: Formula Sheet
% page_url: https://formulasheet.com/#q|l|6084
% subject_slug: formulasheet
% subject_path: /formulasheet/
% section: Chapter 12
% subsection: 
% formula_name: 
% source: formulasheet-latex
% index_global: 127
% index_in_page: 4
\[\begin{equation} \label{Geometric Sum of an Infinte} S_{n}=\frac{a_{1}}{1-r} \end{equation}\]
% formula_record_end

% formula_record_start
% formula_id: formula-sheet-3-005
% page_slug: formula-sheet-3
% page_title: Formula Sheet
% page_url: https://formulasheet.com/#q|l|6084
% subject_slug: formulasheet
% subject_path: /formulasheet/
% section: Chapter 12
% subsection: 
% formula_name: 
% source: formulasheet-latex
% index_global: 128
% index_in_page: 5
\[\begin{equation} \label{Arithmetic Formula with a1} a_{n}=a_{1}+(n-1)d \end{equation}\]
% formula_record_end

% formula_record_start
% formula_id: formula-sheet-3-006
% page_slug: formula-sheet-3
% page_title: Formula Sheet
% page_url: https://formulasheet.com/#q|l|6084
% subject_slug: formulasheet
% subject_path: /formulasheet/
% section: Chapter 12
% subsection: 
% formula_name: 
% source: formulasheet-latex
% index_global: 129
% index_in_page: 6
\[\begin{equation} \label{Arithmetic Formula with n&m} a_{n}=a_{m}+(n-m)d \end{equation}\]
% formula_record_end

% formula_record_start
% formula_id: formula-sheet-3-007
% page_slug: formula-sheet-3
% page_title: Formula Sheet
% page_url: https://formulasheet.com/#q|l|6084
% subject_slug: formulasheet
% subject_path: /formulasheet/
% section: Chapter 12
% subsection: 
% formula_name: 
% source: formulasheet-latex
% index_global: 130
% index_in_page: 7
\[\begin{equation} \label{Arithmetic Sum for Finite} S_{n}=\frac{n}{2}(2a_{1}+(n-1)d) \end{equation}\]
% formula_record_end

% formula_record_start
% formula_id: formula-sheet-3-008
% page_slug: formula-sheet-3
% page_title: Formula Sheet
% page_url: https://formulasheet.com/#q|l|6084
% subject_slug: formulasheet
% subject_path: /formulasheet/
% section: Chapter 12
% subsection: 
% formula_name: 
% source: formulasheet-latex
% index_global: 131
% index_in_page: 8
\[\begin{equation} \label{Combinations Calculation} _{n}C_{r}=\frac{n!}{r!(n-r)!} \end{equation}\]
% formula_record_end

% formula_record_start
% formula_id: formula-sheet-3-009
% page_slug: formula-sheet-3
% page_title: Formula Sheet
% page_url: https://formulasheet.com/#q|l|6084
% subject_slug: formulasheet
% subject_path: /formulasheet/
% section: Chapter 12
% subsection: 
% formula_name: 
% source: formulasheet-latex
% index_global: 132
% index_in_page: 9
\[\begin{equation} \label{Geometric Formula with a1} a_{n}=a_{1}\cdot r^{n-1} \end{equation}\]
% formula_record_end

% formula_record_start
% formula_id: formula-sheet-3-010
% page_slug: formula-sheet-3
% page_title: Formula Sheet
% page_url: https://formulasheet.com/#q|l|6084
% subject_slug: formulasheet
% subject_path: /formulasheet/
% section: Chapter 12
% subsection: 
% formula_name: 
% source: formulasheet-latex
% index_global: 133
% index_in_page: 10
\[\begin{equation} \label{Geometric Formula with n&m} a_{n}=a_{m}\cdot r^{n-1} \end{equation}\]
% formula_record_end

% formula_record_start
% formula_id: formula-sheet-4-001
% page_slug: formula-sheet-4
% page_title: Formula Sheet
% page_url: https://formulasheet.com/#q|l|4981
% subject_slug: formulasheet
% subject_path: /formulasheet/
% section: Chapter 3
% subsection: 
% formula_name: 
% source: formulasheet-latex
% index_global: 134
% index_in_page: 1
\[\begin{equation} c_r = \frac{1}{2}[f(x+) + f(x-)] \end{equation}\]
% formula_record_end

% formula_record_start
% formula_id: formula-sheet-4-002
% page_slug: formula-sheet-4
% page_title: Formula Sheet
% page_url: https://formulasheet.com/#q|l|4981
% subject_slug: formulasheet
% subject_path: /formulasheet/
% section: Chapter 3
% subsection: 
% formula_name: 
% source: formulasheet-latex
% index_global: 135
% index_in_page: 2
\[\begin{equation} \int _{-1}^{1}P_{m}(x)P_{n}(x)\,dx={\frac{2}{2n+1}}\delta _{mn} \end{equation}\]
% formula_record_end

% formula_record_start
% formula_id: formula-sheet-4-003
% page_slug: formula-sheet-4
% page_title: Formula Sheet
% page_url: https://formulasheet.com/#q|l|4981
% subject_slug: formulasheet
% subject_path: /formulasheet/
% section: Chapter 3
% subsection: 
% formula_name: 
% source: formulasheet-latex
% index_global: 136
% index_in_page: 3
\[\begin{equation} \label{legendiffeq} (1-x^{2})\frac{\mathrm{d}^2y}{\mathrm{d}x^2} - 2x\frac{\mathrm{d}y}{\mathrm{d}x} + n(n+1)y = 0 \end{equation}\]
% formula_record_end

% formula_record_start
% formula_id: formula-sheet-4-004
% page_slug: formula-sheet-4
% page_title: Formula Sheet
% page_url: https://formulasheet.com/#q|l|4981
% subject_slug: formulasheet
% subject_path: /formulasheet/
% section: Chapter 3
% subsection: 
% formula_name: 
% source: formulasheet-latex
% index_global: 137
% index_in_page: 4
\[\begin{equation} \label{legenpolydef} P_{n}(x) = \sum_{r=0}^{\left[ n/2 \right]}(-1)^{r}\frac{(2n-2r)!}{2^{n}r!(n-r)!(n-2r)!}x^{n-2r} \end{equation}\]
% formula_record_end

% formula_record_start
% formula_id: formula-sheet-4-005
% page_slug: formula-sheet-4
% page_title: Formula Sheet
% page_url: https://formulasheet.com/#q|l|4981
% subject_slug: formulasheet
% subject_path: /formulasheet/
% section: Chapter 3
% subsection: 
% formula_name: 
% source: formulasheet-latex
% index_global: 138
% index_in_page: 5
\[\begin{equation} \label{legenvals} {\begin{array}{r|r}n&P_{n}(x)\\\hline 0&1\\1&x\\2&{\frac {1}{2}}(3x^{2}-1)\\3&{\frac {1}{2}}(5x^{3}-3x)\\4&{\frac {1}{8}}(35x^{4}-30x^{2}+3) \end{array}} \end{equation}\]
% formula_record_end

% formula_record_start
% formula_id: formula-sheet-4-006
% page_slug: formula-sheet-4
% page_title: Formula Sheet
% page_url: https://formulasheet.com/#q|l|4981
% subject_slug: formulasheet
% subject_path: /formulasheet/
% section: Chapter 3
% subsection: 
% formula_name: 
% source: formulasheet-latex
% index_global: 139
% index_in_page: 6
\[\begin{equation} \label{thrm3.2} P_{n}(x)={\frac{1}{2^{n}n!}}\frac{\mathrm{d}^{n}}{\mathrm{d}x^{n}}(x^{2}-1)^{n} \end{equation}\]
% formula_record_end

% formula_record_start
% formula_id: formula-sheet-4-007
% page_slug: formula-sheet-4
% page_title: Formula Sheet
% page_url: https://formulasheet.com/#q|l|4981
% subject_slug: formulasheet
% subject_path: /formulasheet/
% section: Chapter 3
% subsection: 
% formula_name: 
% source: formulasheet-latex
% index_global: 140
% index_in_page: 7
\[\begin{equation} \label{thrm3.3} P_{n}(x) = \frac{1}{\pi}\int_{0}^{\pi}(x + \sqrt{(x^{2}-1)}\cos{\theta})^{n}\,d\theta \end{equation}\]
% formula_record_end

% formula_record_start
% formula_id: formula-sheet-4-008
% page_slug: formula-sheet-4
% page_title: Formula Sheet
% page_url: https://formulasheet.com/#q|l|4981
% subject_slug: formulasheet
% subject_path: /formulasheet/
% section: Chapter 3
% subsection: 
% formula_name: 
% source: formulasheet-latex
% index_global: 141
% index_in_page: 8
\[\begin{equation} \label{thrm3.4} {\begin{aligned} & (i)\, P_{n}(1)=1\\\ & (ii)\, P_{n}(-1) = (-1)^{n}\\\ & (iii)\, P_{n}^{\;'}(1) = \frac{1}{2}n(n+1)\\\ & (iv)\, P_{n}^{\;'}(-1) = (-1)^{n-1}\frac{1}{2}n(n+1)\\\ & (v)\, P_{2n}(0) = (-1)^{n}\frac{(2n)!}{2^{2n}(n!)^{2}}\\\ & (vi)\, P_{2n+1}(0) = 0 \end{aligned}} \end{equation}\]
% formula_record_end

% formula_record_start
% formula_id: formula-sheet-4-009
% page_slug: formula-sheet-4
% page_title: Formula Sheet
% page_url: https://formulasheet.com/#q|l|4981
% subject_slug: formulasheet
% subject_path: /formulasheet/
% section: Chapter 3
% subsection: 
% formula_name: 
% source: formulasheet-latex
% index_global: 142
% index_in_page: 9
\[\begin{equation} \label{thrm3.6} f(x) = \sum_{r=0}^{n}\, c_rP_r(x),\quad c_r = (r + \frac{1}{2})\int_{-1}^{1}f(x)P_r(x)\, dx \end{equation}\]
% formula_record_end

% formula_record_start
% formula_id: formula-sheet-4-010
% page_slug: formula-sheet-4
% page_title: Formula Sheet
% page_url: https://formulasheet.com/#q|l|4981
% subject_slug: formulasheet
% subject_path: /formulasheet/
% section: Chapter 3
% subsection: 
% formula_name: 
% source: formulasheet-latex
% index_global: 143
% index_in_page: 10
\[\begin{equation} \label{thrm3.8} {\begin{aligned} & (i)\, P_n^{\;'}(x) = \sum_{r=0}^{\left[ \frac{1}{2}(n-1) \right]}(2n-4r-1)P_{n-2r-1}(x) \\& (ii)\, xP_n(x) = \frac{n+1}{2n+1}P_{n+1}(x) + \frac{n}{2n+1}P_{n-1}(x) \\& (iii)\, (n+1)P_{n+1}(x) - (2n+1)xP_n(x) + nP_{n-1}(x) = 0 \\& (iv)\, P_{n+1}^{\; '}(x) - P_{n-1}^{\; '}(x) = (2n+1)P_n(x) \\& (v)\, xP_n^{\; '}(x) - P_{n-1}^{\; '}(x) = nP_n(x) \\& (vi)\, P_n^{\; '}(x) - xP_{n-1}^{\; '}(x) = nP_{n-1}(x) \\& (vii)\, (x^2-1)P_{n}^{\; '}(x) = nxP_n(x) - nP_{n-1}(x) \\& (viii)\, (x^2-1)P_{n}^{\; '}(x) = (n+1)P_{n+1}(x) - (n+1)xP_n(x) \\& (ix)\, \sum_{k=0}^n(2k+1)P_k(x)P_k(y) = \frac{n+1}{x-y}\left[ P_{n+1}(x)P_n(y) - P_n(x)P_{n+1}(y) \right] \end{aligned}} \end{equation}\]
% formula_record_end

% formula_record_start
% formula_id: formula-sheet-4-011
% page_slug: formula-sheet-4
% page_title: Formula Sheet
% page_url: https://formulasheet.com/#q|l|4981
% subject_slug: formulasheet
% subject_path: /formulasheet/
% section: Chapter 3
% subsection: 
% formula_name: 
% source: formulasheet-latex
% index_global: 144
% index_in_page: 11
\[\begin{equation} \label{thrm3.1} {\frac {1}{\sqrt {1-2xt+t^{2}}}}=\sum _{n=0}^{\infty }t^{n}P_{n}(x) \end{equation}\]
% formula_record_end

% formula_record_start
% formula_id: formula-sheet-5-001
% page_slug: formula-sheet-5
% page_title: Formula Sheet
% page_url: https://formulasheet.com/#q|l|4989
% subject_slug: formulasheet
% subject_path: /formulasheet/
% section: Chapter 4
% subsection: 
% formula_name: 
% source: formulasheet-latex
% index_global: 145
% index_in_page: 1
\[\begin{align} \label{thrm 4.10} All \; results \; from \; 4.8 \; hold \; for \; H_n^{(1)}(x) \; and \;H_n^{(2)}(x) \; instead \; of \; J_n(x) \end{align}\]
% formula_record_end

% formula_record_start
% formula_id: formula-sheet-5-002
% page_slug: formula-sheet-5
% page_title: Formula Sheet
% page_url: https://formulasheet.com/#q|l|4989
% subject_slug: formulasheet
% subject_path: /formulasheet/
% section: Chapter 4
% subsection: 
% formula_name: 
% source: formulasheet-latex
% index_global: 146
% index_in_page: 2
\[\begin{align} \label{thrm4.11(a)} x^2\frac{\mathrm d^2y}{\mathrm dx^2} + x\frac{\mathrm dy}{\mathrm dx} + (\lambda^2x^2-n^2)y = 0 \end{align}\]
% formula_record_end

% formula_record_start
% formula_id: formula-sheet-5-003
% page_slug: formula-sheet-5
% page_title: Formula Sheet
% page_url: https://formulasheet.com/#q|l|4989
% subject_slug: formulasheet
% subject_path: /formulasheet/
% section: Chapter 4
% subsection: 
% formula_name: 
% source: formulasheet-latex
% index_global: 147
% index_in_page: 3
\[\begin{align} \label{thrm4.11(b)} AJ_n(\lambda x) + BY_n(\lambda x) \end{align}\]
% formula_record_end

% formula_record_start
% formula_id: formula-sheet-5-004
% page_slug: formula-sheet-5
% page_title: Formula Sheet
% page_url: https://formulasheet.com/#q|l|4989
% subject_slug: formulasheet
% subject_path: /formulasheet/
% section: Chapter 4
% subsection: 
% formula_name: 
% source: formulasheet-latex
% index_global: 148
% index_in_page: 4
\[\begin{align} \label{thrm4.12(a)} x^2\frac{\mathrm d^2y}{\mathrm dx^2} + (1-2\alpha)\frac{\mathrm dy}{\mathrm dx} + (\beta^2\gamma^2x^{2\gamma}+(\alpha^2-n^2\gamma^2))y = 0 \end{align}\]
% formula_record_end

% formula_record_start
% formula_id: formula-sheet-5-005
% page_slug: formula-sheet-5
% page_title: Formula Sheet
% page_url: https://formulasheet.com/#q|l|4989
% subject_slug: formulasheet
% subject_path: /formulasheet/
% section: Chapter 4
% subsection: 
% formula_name: 
% source: formulasheet-latex
% index_global: 149
% index_in_page: 5
\[\begin{align} \label{thrm4.12(b)} Ax^{\alpha}J_n(\beta x^{\gamma})+Bx^{\alpha}Y_n(\beta x^\gamma) \end{align}\]
% formula_record_end

% formula_record_start
% formula_id: formula-sheet-5-006
% page_slug: formula-sheet-5
% page_title: Formula Sheet
% page_url: https://formulasheet.com/#q|l|4989
% subject_slug: formulasheet
% subject_path: /formulasheet/
% section: Chapter 4
% subsection: 
% formula_name: 
% source: formulasheet-latex
% index_global: 150
% index_in_page: 6
\[\begin{align} \label{thrm4.3} Y_n(x) = & \frac2\pi\left[ ln\frac x2 + \gamma -\frac 12 \sum_{r=1}^n\frac 1r \right]J_n(x) \\\ & -\frac 1\pi\sum_{s=0}^\infty (-1)^s\frac{1}{s!(n+s)!}\left( \frac x2 \right) ^{n+2s}\sum_{r=1}^s\left[ \frac 1r + \frac 1{r+n} \right] \\\ & - \frac 1\pi \sum_{s=0}^{n-1}\frac{(n-s-1)!}{s!}\left( \frac x2 \right)^{-n+2s} \end{align}\]
% formula_record_end

% formula_record_start
% formula_id: formula-sheet-5-007
% page_slug: formula-sheet-5
% page_title: Formula Sheet
% page_url: https://formulasheet.com/#q|l|4989
% subject_slug: formulasheet
% subject_path: /formulasheet/
% section: Chapter 4
% subsection: 
% formula_name: 
% source: formulasheet-latex
% index_global: 151
% index_in_page: 7
\[\begin{align} \label{thrm4.4} Y_{-n}(x) = (-1)^nY_n(x) \end{align}\]
% formula_record_end

% formula_record_start
% formula_id: formula-sheet-5-008
% page_slug: formula-sheet-5
% page_title: Formula Sheet
% page_url: https://formulasheet.com/#q|l|4989
% subject_slug: formulasheet
% subject_path: /formulasheet/
% section: Chapter 4
% subsection: 
% formula_name: 
% source: formulasheet-latex
% index_global: 152
% index_in_page: 8
\[\begin{align} \label{thrm4.5} e^{\left[ \frac 12x(t-\frac 1t) \right]} = \sum_{n= -\infty}^{\infty}t^nJ_n(x) \end{align}\]
% formula_record_end

% formula_record_start
% formula_id: formula-sheet-5-009
% page_slug: formula-sheet-5
% page_title: Formula Sheet
% page_url: https://formulasheet.com/#q|l|4989
% subject_slug: formulasheet
% subject_path: /formulasheet/
% section: Chapter 4
% subsection: 
% formula_name: 
% source: formulasheet-latex
% index_global: 153
% index_in_page: 9
\[\begin{align} \label{thrm4.7} J_n(x) = \frac{(\tfrac 12x)^n}{\sqrt{(\pi)}\, \Gamma(n+\tfrac 12)}\int_{-1}^1(1-t^2)^{n- \tfrac 12}e^{ixt}\mathrm dt,\quad (n>-\tfrac 12) \end{align}\]
% formula_record_end

% formula_record_start
% formula_id: formula-sheet-5-010
% page_slug: formula-sheet-5
% page_title: Formula Sheet
% page_url: https://formulasheet.com/#q|l|4989
% subject_slug: formulasheet
% subject_path: /formulasheet/
% section: Chapter 4
% subsection: 
% formula_name: 
% source: formulasheet-latex
% index_global: 154
% index_in_page: 10
\[\begin{align} \label{thrm4.8} \\\ & (i)\; \frac{\mathrm d}{\mathrm dx}\left[ x^nJ_n(x) \right] = x^nJ_{n-1}(x) \\\ & (ii)\; \frac{\mathrm d}{\mathrm dx}\left[ x^{-n}J_n(x) \right] = -x^{-n}J_{n+1}(x) \\\ & (iii)\; J_n^{\; '}(x) = J_{n-1}(x) - \frac nxJ_n(x) \\\ & (iv)\; J_n^{\; '}(x) =\frac nxJ_n(x) - J_{n+1}(x) \\\ & (v)\; J_n^{\; '}(x) = \tfrac 12 \left[J_{n-1}(x) - J_{n+1}(x) \right] \\\ & (vi)\; J_{n-1}(x) + J_{n+1}(x) = \frac {2n}xJ_n(x) \end{align}\]
% formula_record_end

% formula_record_start
% formula_id: formula-sheet-5-011
% page_slug: formula-sheet-5
% page_title: Formula Sheet
% page_url: https://formulasheet.com/#q|l|4989
% subject_slug: formulasheet
% subject_path: /formulasheet/
% section: Chapter 4
% subsection: 
% formula_name: 
% source: formulasheet-latex
% index_global: 155
% index_in_page: 11
\[\begin{align} \label{thrm4.9} All \; results \; from \; 4.8 \; hold \; for \; Y_n(x) \; instead \; of \; J_n(x) \end{align}\]
% formula_record_end

% formula_record_start
% formula_id: formula-sheet-5-012
% page_slug: formula-sheet-5
% page_title: Formula Sheet
% page_url: https://formulasheet.com/#q|l|4989
% subject_slug: formulasheet
% subject_path: /formulasheet/
% section: Chapter 4
% subsection: 
% formula_name: 
% source: formulasheet-latex
% index_global: 156
% index_in_page: 12
\[\begin{align} \label{hankdef} & H_n^{(1)}(x) = J_n(x) + iY_n(x)\\\ & H_n^{(2)}(x) = J_n(x) - iY_n(x) \end{align}\]
% formula_record_end

% formula_record_start
% formula_id: formula-sheet-5-013
% page_slug: formula-sheet-5
% page_title: Formula Sheet
% page_url: https://formulasheet.com/#q|l|4989
% subject_slug: formulasheet
% subject_path: /formulasheet/
% section: Chapter 4
% subsection: 
% formula_name: 
% source: formulasheet-latex
% index_global: 157
% index_in_page: 13
\[\begin{equation} J_{-n}(x)=(-1)^{n}J_{n}(x) \end{equation}\]
% formula_record_end

% formula_record_start
% formula_id: formula-sheet-5-014
% page_slug: formula-sheet-5
% page_title: Formula Sheet
% page_url: https://formulasheet.com/#q|l|4989
% subject_slug: formulasheet
% subject_path: /formulasheet/
% section: Chapter 4
% subsection: 
% formula_name: 
% source: formulasheet-latex
% index_global: 158
% index_in_page: 14
\[\begin{equation} x^{2}{\frac {d^{2}y}{dx^{2}}}+x{\frac {dy}{dx}}+(x^{2}-n ^{2})y=0 \end{equation}\]
% formula_record_end

% formula_record_start
% formula_id: formula-sheet-5-015
% page_slug: formula-sheet-5
% page_title: Formula Sheet
% page_url: https://formulasheet.com/#q|l|4989
% subject_slug: formulasheet
% subject_path: /formulasheet/
% section: Chapter 4
% subsection: 
% formula_name: 
% source: formulasheet-latex
% index_global: 159
% index_in_page: 15
\[\begin{equation} \label{besselfirst1} J_{-n}(x)=\sum _{r=0}^{\infty }(-1)^r{\frac {1}{r!\,\Gamma (-n+r +1)}}{\left({\frac {x}{2}}\right)}^{2r-n }, and \end{equation}\]
% formula_record_end

% formula_record_start
% formula_id: formula-sheet-5-016
% page_slug: formula-sheet-5
% page_title: Formula Sheet
% page_url: https://formulasheet.com/#q|l|4989
% subject_slug: formulasheet
% subject_path: /formulasheet/
% section: Chapter 4
% subsection: 
% formula_name: 
% source: formulasheet-latex
% index_global: 160
% index_in_page: 16
\[\begin{equation} \label{besselfirstn} J_{n }(x)=\sum _{r=0}^{\infty }(-1)^r{\frac {1}{r!\,\Gamma (n+r +1)}}{\left({\frac {x}{2}}\right)}^{2r+n }, \end{equation}\]
% formula_record_end

% formula_record_start
% formula_id: formula-sheet-5-017
% page_slug: formula-sheet-5
% page_title: Formula Sheet
% page_url: https://formulasheet.com/#q|l|4989
% subject_slug: formulasheet
% subject_path: /formulasheet/
% section: Chapter 4
% subsection: 
% formula_name: 
% source: formulasheet-latex
% index_global: 161
% index_in_page: 17
\[\begin{equation} \label{besselsecond} Y_{n }(x)={\frac {\cos(n \pi )J_n(x)-J_{-n }(x)}{\sin(n \pi )}} \end{equation}\]
% formula_record_end

% formula_record_start
% formula_id: formula-sheet-5-018
% page_slug: formula-sheet-5
% page_title: Formula Sheet
% page_url: https://formulasheet.com/#q|l|4989
% subject_slug: formulasheet
% subject_path: /formulasheet/
% section: Chapter 4
% subsection: 
% formula_name: 
% source: formulasheet-latex
% index_global: 162
% index_in_page: 18
\[\begin{equation} \label{thrm4.6} J_{n }(x)={\frac {1}{\pi }}\int _{0}^{\pi }\cos(n \phi -x\sin \phi )\,d\phi \end{equation}\]
% formula_record_end

% formula_record_start
% formula_id: formula-sheet-6-001
% page_slug: formula-sheet-6
% page_title: Formula Sheet
% page_url: https://formulasheet.com/#q|l|4982
% subject_slug: formulasheet
% subject_path: /formulasheet/
% section: Chapter 5
% subsection: 
% formula_name: 
% source: formulasheet-latex
% index_global: 163
% index_in_page: 1
\[\begin{align} \label{thrm5.6} \\\ & (i)\, H_n^{\; '}(x) = 2nH_{n-1}(x), \;\; (n \geq 1); \quad H_0^{\; '}(x) = 0 \\\ & (ii)\, H_{n+1}(x)=2xH_{n}(x)-2nH_{n-1}(x),\;\; (n\geq 1); \quad H_1(x) = 2xH_0(x) \end{align}\]
% formula_record_end

% formula_record_start
% formula_id: formula-sheet-6-002
% page_slug: formula-sheet-6
% page_title: Formula Sheet
% page_url: https://formulasheet.com/#q|l|4982
% subject_slug: formulasheet
% subject_path: /formulasheet/
% section: Chapter 5
% subsection: 
% formula_name: 
% source: formulasheet-latex
% index_global: 164
% index_in_page: 2
\[\begin{equation} \label{hermdef} H_{n}(x)=\sum _{r=0}^{\left[ {\frac {1}{2}n}\right ] }(-1)^r\frac{n!}{r!(n-2r)!}(2x)^{n-2r} \end{equation}\]
% formula_record_end

% formula_record_start
% formula_id: formula-sheet-6-003
% page_slug: formula-sheet-6
% page_title: Formula Sheet
% page_url: https://formulasheet.com/#q|l|4982
% subject_slug: formulasheet
% subject_path: /formulasheet/
% section: Chapter 5
% subsection: 
% formula_name: 
% source: formulasheet-latex
% index_global: 165
% index_in_page: 3
\[\begin{equation} \label{hermdiffeq} \frac{\mathrm{d}^2y}{\mathrm{d}x^2} - 2x\frac{\mathrm{d}y}{\mathrm{d}x} + 2ny = 0 \end{equation}\]
% formula_record_end

% formula_record_start
% formula_id: formula-sheet-6-004
% page_slug: formula-sheet-6
% page_title: Formula Sheet
% page_url: https://formulasheet.com/#q|l|4982
% subject_slug: formulasheet
% subject_path: /formulasheet/
% section: Chapter 5
% subsection: 
% formula_name: 
% source: formulasheet-latex
% index_global: 166
% index_in_page: 4
\[\begin{equation} \label{hermvals} {\begin{array}{r|r}n&H_{n}(x)\\\hline 0&1\\1&2x\\2&4x^{2}-2\\3&8x^{3}-12x\\4&16x^{4}-48x^{2}+12\\5&32x^5-160x^3+120x \end{array}} \end{equation}\]
% formula_record_end

% formula_record_start
% formula_id: formula-sheet-6-005
% page_slug: formula-sheet-6
% page_title: Formula Sheet
% page_url: https://formulasheet.com/#q|l|4982
% subject_slug: formulasheet
% subject_path: /formulasheet/
% section: Chapter 5
% subsection: 
% formula_name: 
% source: formulasheet-latex
% index_global: 167
% index_in_page: 5
\[\begin{equation} \label{thrm5.1} e^{2xt-t^{2}}=\sum _{n=0}^{\infty }{\frac {t^{n}}{n!}}H_n(x)\,\! \end{equation}\]
% formula_record_end

% formula_record_start
% formula_id: formula-sheet-6-006
% page_slug: formula-sheet-6
% page_title: Formula Sheet
% page_url: https://formulasheet.com/#q|l|4982
% subject_slug: formulasheet
% subject_path: /formulasheet/
% section: Chapter 5
% subsection: 
% formula_name: 
% source: formulasheet-latex
% index_global: 168
% index_in_page: 6
\[\begin{equation} \label{thrm5.3(ii)} \exp{(\frac{\mathrm{d}}{\mathrm{d}x})}f(x) = \sum_{n=0}^{\infty}\frac{1}{n!}\frac{\mathrm{d}^nf}{\mathrm{d}x^n} \end{equation}\]
% formula_record_end

% formula_record_start
% formula_id: formula-sheet-6-007
% page_slug: formula-sheet-6
% page_title: Formula Sheet
% page_url: https://formulasheet.com/#q|l|4982
% subject_slug: formulasheet
% subject_path: /formulasheet/
% section: Chapter 5
% subsection: 
% formula_name: 
% source: formulasheet-latex
% index_global: 169
% index_in_page: 7
\[\begin{equation} \label{thrm5.5} \int _{-\infty }^{\infty }e^{-x^2}H_{n}(x)H_{m}(x)\,\mathrm {d} x=2^{n}n!(\sqrt{\pi})\delta _{nm} \end{equation}\]
% formula_record_end

% formula_record_start
% formula_id: formula-sheet-6-008
% page_slug: formula-sheet-6
% page_title: Formula Sheet
% page_url: https://formulasheet.com/#q|l|4982
% subject_slug: formulasheet
% subject_path: /formulasheet/
% section: Chapter 5
% subsection: 
% formula_name: 
% source: formulasheet-latex
% index_global: 170
% index_in_page: 8
\[\begin{equation} \label{weberherm(i)} \Psi_n(x) = e^{-x^2/2}H_n(x), solves: \end{equation}\]
% formula_record_end

% formula_record_start
% formula_id: formula-sheet-6-009
% page_slug: formula-sheet-6
% page_title: Formula Sheet
% page_url: https://formulasheet.com/#q|l|4982
% subject_slug: formulasheet
% subject_path: /formulasheet/
% section: Chapter 5
% subsection: 
% formula_name: 
% source: formulasheet-latex
% index_global: 171
% index_in_page: 9
\[\begin{equation} \label{weberherm(ii)} \frac{\mathrm{d}^2y}{\mathrm{d}x^2} + (\lambda - x^2)y = 0\quad where\; \lambda = 2n+1 \end{equation}\]
% formula_record_end

% formula_record_start
% formula_id: formula-sheet-6-010
% page_slug: formula-sheet-6
% page_title: Formula Sheet
% page_url: https://formulasheet.com/#q|l|4982
% subject_slug: formulasheet
% subject_path: /formulasheet/
% section: Chapter 5
% subsection: 
% formula_name: 
% source: formulasheet-latex
% index_global: 172
% index_in_page: 10
\[\begin{equation} \label{thrm5.2} H_{n}(x)=(-1)^{n}e^{x^{2}}{\frac {d^{n}}{dx^{n}}}e^{-x^{2}} \end{equation}\]
% formula_record_end

% formula_record_start
% formula_id: formula-sheet-6-011
% page_slug: formula-sheet-6
% page_title: Formula Sheet
% page_url: https://formulasheet.com/#q|l|4982
% subject_slug: formulasheet
% subject_path: /formulasheet/
% section: Chapter 5
% subsection: 
% formula_name: 
% source: formulasheet-latex
% index_global: 173
% index_in_page: 11
\[\begin{equation} \label{thrm5.3(i)} H_n(x) = 2^n[\exp{(-\frac{1}{4}\frac{\mathrm{d}^2}{\mathrm{d}x^2})}]x^n, where \end{equation}\]
% formula_record_end

% formula_record_start
% formula_id: formula-sheet-6-012
% page_slug: formula-sheet-6
% page_title: Formula Sheet
% page_url: https://formulasheet.com/#q|l|4982
% subject_slug: formulasheet
% subject_path: /formulasheet/
% section: Chapter 5
% subsection: 
% formula_name: 
% source: formulasheet-latex
% index_global: 174
% index_in_page: 12
\[\begin{equation} \label{thrm5.4(i)} H_{2n}(0) = (-1)^n\frac{(2n)!}{n!} \end{equation}\]
% formula_record_end

% formula_record_start
% formula_id: formula-sheet-6-013
% page_slug: formula-sheet-6
% page_title: Formula Sheet
% page_url: https://formulasheet.com/#q|l|4982
% subject_slug: formulasheet
% subject_path: /formulasheet/
% section: Chapter 5
% subsection: 
% formula_name: 
% source: formulasheet-latex
% index_global: 175
% index_in_page: 13
\[\begin{equation} \label{thrm5.4(ii)} H_{2n+1}(0) = 0 \end{equation}\]
% formula_record_end

% formula_record_start
% formula_id: formula-sheet-7-001
% page_slug: formula-sheet-7
% page_title: Formula Sheet
% page_url: https://formulasheet.com/#q|l|1022
% subject_slug: formulasheet
% subject_path: /formulasheet/
% section: Integral Table [integral-table.com]
% subsection: 
% formula_name: 
% source: formulasheet-latex
% index_global: 176
% index_in_page: 1
\[\begin{align} \int x^n cos x dx &= -\frac{1}{2}(i)^{n+1}\left [ \Gamma(n+1, -ix) \right . \nonumber \\\ & \left . + (-1)^n \Gamma(n+1, ix)\right] \end{align}\]
% formula_record_end

% formula_record_start
% formula_id: formula-sheet-7-002
% page_slug: formula-sheet-7
% page_title: Formula Sheet
% page_url: https://formulasheet.com/#q|l|1022
% subject_slug: formulasheet
% subject_path: /formulasheet/
% section: Integral Table [integral-table.com]
% subsection: 
% formula_name: 
% source: formulasheet-latex
% index_global: 177
% index_in_page: 2
\[\begin{align} \int \sqrt{a^2 - x^2} dx &= \frac{1}{2} x \sqrt{a^2-x^2} %\nonumber \\\ & +\frac{1}{2}a^2\tan^{-1}\frac{x}{\sqrt{a^2-x^2}} \end{align}\]
% formula_record_end

% formula_record_start
% formula_id: formula-sheet-7-003
% page_slug: formula-sheet-7
% page_title: Formula Sheet
% page_url: https://formulasheet.com/#q|l|1022
% subject_slug: formulasheet
% subject_path: /formulasheet/
% section: Integral Table [integral-table.com]
% subsection: 
% formula_name: 
% source: formulasheet-latex
% index_global: 178
% index_in_page: 3
\[\begin{align} \int & e^{ax} \tanh bx dx = \nonumber \\\ & \begin{cases} { \frac{ e^{(a+2b)x}}{(a+2b)} {_2F_1}\left[ 1+\frac{a}{2b},1,2+\frac{a}{2b}, -e^{2bx}\right] }& \\\ { \hspace{1cm}-\frac{1}{a}e^{ax}{_2F_1}\left[ \frac{a}{2b},1,1E, -e^{2bx}\right] } & a\ne b \\\ {\frac{e^{ax}-2\tan^{-1}[e^{ax}]}{a} } & a = b \end{cases} \end{align}\]
% formula_record_end

% formula_record_start
% formula_id: formula-sheet-7-004
% page_slug: formula-sheet-7
% page_title: Formula Sheet
% page_url: https://formulasheet.com/#q|l|1022
% subject_slug: formulasheet
% subject_path: /formulasheet/
% section: Integral Table [integral-table.com]
% subsection: 
% formula_name: 
% source: formulasheet-latex
% index_global: 179
% index_in_page: 4
\[\begin{align} \int &x \sqrt{a x^2 + bx + c} = \frac{1}{48a^{5/2}}\left ( 2 \sqrt{a} \sqrt{ax^2+bx+c} \right . \nonumber \\\ & \times \left( -3b^2 + 2 abx + 8 a(c+ax^2) \right) \nonumber \\\ & \left. + 3(b^3-4abc)\ln \left|b + 2ax + 2\sqrt{a}\sqrt{ax^2+bx+x} \right| \right) \end{align}\]
% formula_record_end

% formula_record_start
% formula_id: formula-sheet-7-005
% page_slug: formula-sheet-7
% page_title: Formula Sheet
% page_url: https://formulasheet.com/#q|l|1022
% subject_slug: formulasheet
% subject_path: /formulasheet/
% section: Integral Table [integral-table.com]
% subsection: 
% formula_name: 
% source: formulasheet-latex
% index_global: 180
% index_in_page: 5
\[\begin{align} \int &x \sqrt{ax + b}dx = %\nonumber \\& \frac{2}{15 a^2}(-2b^2+abx + 3 a^2 x^2) \sqrt{ax+b} \end{align}\]
% formula_record_end

% formula_record_start
% formula_id: formula-sheet-7-006
% page_slug: formula-sheet-7
% page_title: Formula Sheet
% page_url: https://formulasheet.com/#q|l|1022
% subject_slug: formulasheet
% subject_path: /formulasheet/
% section: Integral Table [integral-table.com]
% subsection: 
% formula_name: 
% source: formulasheet-latex
% index_global: 181
% index_in_page: 6
\[\begin{align} \int &\frac{x}{\sqrt{ax^2+bx+c}}dx= \frac{1}{a}\sqrt{ax^2+bx + c} \nonumber \\& - \frac{b}{2a^{3/2}}\ln \left| 2ax+b + 2 \sqrt{a(ax^2+bx+c)} \right | \end{align}\]
% formula_record_end

% formula_record_start
% formula_id: formula-sheet-7-007
% page_slug: formula-sheet-7
% page_title: Formula Sheet
% page_url: https://formulasheet.com/#q|l|1022
% subject_slug: formulasheet
% subject_path: /formulasheet/
% section: Integral Table [integral-table.com]
% subsection: 
% formula_name: 
% source: formulasheet-latex
% index_global: 182
% index_in_page: 7
\[\begin{align} \int &\sin^n ax dx = \nonumber \\\ & -\frac{1}{a}{\cos ax} \hspace{2mm}{_2F_1}\left[ \frac{1}{2}, \frac{1-n}{2}, \frac{3}{2}, \cos^2 ax \right] \end{align}\]
% formula_record_end

% formula_record_start
% formula_id: formula-sheet-7-008
% page_slug: formula-sheet-7
% page_title: Formula Sheet
% page_url: https://formulasheet.com/#q|l|1022
% subject_slug: formulasheet
% subject_path: /formulasheet/
% section: Integral Table [integral-table.com]
% subsection: 
% formula_name: 
% source: formulasheet-latex
% index_global: 183
% index_in_page: 8
\[\begin{align} \int &\sqrt{a x^2 + b x + c} dx = \frac{b+2ax}{4a}\sqrt{ax^2+bx+c} \nonumber \\\ & + \frac{4ac-b^2}{8a^{3/2}}\ln \left| 2ax + b + 2\sqrt{a(ax^2+bx^+c)}\right | \end{align}\]
% formula_record_end

% formula_record_start
% formula_id: formula-sheet-7-009
% page_slug: formula-sheet-7
% page_title: Formula Sheet
% page_url: https://formulasheet.com/#q|l|1022
% subject_slug: formulasheet
% subject_path: /formulasheet/
% section: Integral Table [integral-table.com]
% subsection: 
% formula_name: 
% source: formulasheet-latex
% index_global: 184
% index_in_page: 9
\[\begin{align} \int &\tan^n ax dx = \frac{\tan^{n+1} ax }{a(1+n)} \times \nonumber \\\ & {_2}F_1\left( \frac{n+1}{2}, 1, \frac{n+3}{2}, -\tan^2 ax \right) \end{align}\]
% formula_record_end

% formula_record_start
% formula_id: formula-sheet-7-010
% page_slug: formula-sheet-7
% page_title: Formula Sheet
% page_url: https://formulasheet.com/#q|l|1022
% subject_slug: formulasheet
% subject_path: /formulasheet/
% section: Integral Table [integral-table.com]
% subsection: 
% formula_name: 
% source: formulasheet-latex
% index_global: 185
% index_in_page: 10
\[\begin{align} \int e^{ax} & \cosh bx dx = \nonumber \\\ & \begin{cases} {\frac{e^{ax}}{a^2-b^2} }[ a \cosh bx - b \sinh bx ] & a\ne b \\\ {\frac{e^{2ax}}{4a} + \frac{x}{2}} & a = b \end{cases} \end{align}\]
% formula_record_end

% formula_record_start
% formula_id: formula-sheet-7-011
% page_slug: formula-sheet-7
% page_title: Formula Sheet
% page_url: https://formulasheet.com/#q|l|1022
% subject_slug: formulasheet
% subject_path: /formulasheet/
% section: Integral Table [integral-table.com]
% subsection: 
% formula_name: 
% source: formulasheet-latex
% index_global: 186
% index_in_page: 11
\[\begin{align} \int e^{ax}& \sinh bx dx = \nonumber \\\ & \begin{cases} {\frac{e^{ax}}{a^2-b^2} }[ -b \cosh bx + a \sinh bx ] & a\ne b \\\ {\frac{e^{2ax}}{4a} - \frac{x}{2}} & a = b \end{cases} \end{align}\]
% formula_record_end

% formula_record_start
% formula_id: formula-sheet-7-012
% page_slug: formula-sheet-7
% page_title: Formula Sheet
% page_url: https://formulasheet.com/#q|l|1022
% subject_slug: formulasheet
% subject_path: /formulasheet/
% section: Integral Table [integral-table.com]
% subsection: 
% formula_name: 
% source: formulasheet-latex
% index_global: 187
% index_in_page: 12
\[\begin{align} \int x \ln (ax + b) dx &= \frac{bx}{2a}-\frac{1}{4}x^2 \nonumber \\& +\frac{1}{2}\left(x^2-\frac{b^2}{a^2}\right)\ln (ax+b) \end{align}\]
% formula_record_end

% formula_record_start
% formula_id: formula-sheet-7-013
% page_slug: formula-sheet-7
% page_title: Formula Sheet
% page_url: https://formulasheet.com/#q|l|1022
% subject_slug: formulasheet
% subject_path: /formulasheet/
% section: Integral Table [integral-table.com]
% subsection: 
% formula_name: 
% source: formulasheet-latex
% index_global: 188
% index_in_page: 13
\[\begin{align} \int x \ln \left ( a^2 - b^2 x^2 \right ) dx &= -\frac{1}{2}x^2+ \nonumber \\& \frac{1}{2}\left( x^2 - \frac{a^2}{b^2} \right ) \ln \left (a^2 -b^2 x^2 \right) \end{align}\]
% formula_record_end

% formula_record_start
% formula_id: formula-sheet-7-014
% page_slug: formula-sheet-7
% page_title: Formula Sheet
% page_url: https://formulasheet.com/#q|l|1022
% subject_slug: formulasheet
% subject_path: /formulasheet/
% section: Integral Table [integral-table.com]
% subsection: 
% formula_name: 
% source: formulasheet-latex
% index_global: 189
% index_in_page: 14
\[\begin{align} \int x^n cos ax dx &= \frac{1}{2}(ia)^{1-n}\left [ (-1)^n \Gamma(n+1, -iax) \right. \nonumber \\\ & \left. -\Gamma(n+1, ixa)\right] \end{align}\]
% formula_record_end

% formula_record_start
% formula_id: formula-sheet-7-015
% page_slug: formula-sheet-7
% page_title: Formula Sheet
% page_url: https://formulasheet.com/#q|l|1022
% subject_slug: formulasheet
% subject_path: /formulasheet/
% section: Integral Table [integral-table.com]
% subsection: 
% formula_name: 
% source: formulasheet-latex
% index_global: 190
% index_in_page: 15
\[\begin{align} \int x^n \sin x dx &= -\frac{1}{2}(i)^n\left[ \Gamma(n+1, -ix) %\right. \nonumber \\\ & \left. - (-1)^n\Gamma(n+1, -ix)\right] \end{align}\]
% formula_record_end

% formula_record_start
% formula_id: formula-sheet-7-016
% page_slug: formula-sheet-7
% page_title: Formula Sheet
% page_url: https://formulasheet.com/#q|l|1022
% subject_slug: formulasheet
% subject_path: /formulasheet/
% section: Integral Table [integral-table.com]
% subsection: 
% formula_name: 
% source: formulasheet-latex
% index_global: 191
% index_in_page: 16
\[\begin{align} \int \cos ax \cosh bx dx &= \frac{1}{a^2 + b^2} \left[ a \sin ax \cosh bx \right . \nonumber \\\ & \left. + b \cos ax \sinh bx \right] \end{align}\]
% formula_record_end

% formula_record_start
% formula_id: formula-sheet-7-017
% page_slug: formula-sheet-7
% page_title: Formula Sheet
% page_url: https://formulasheet.com/#q|l|1022
% subject_slug: formulasheet
% subject_path: /formulasheet/
% section: Integral Table [integral-table.com]
% subsection: 
% formula_name: 
% source: formulasheet-latex
% index_global: 192
% index_in_page: 17
\[\begin{align} \int \cos ax \sin bx dx &= \frac{\cos[(a-b) x]}{2(a-b)} - %\nonumber \\\ & \frac{\cos[(a+b)x]}{2(a+b)} , a\ne b \end{align}\]
% formula_record_end

% formula_record_start
% formula_id: formula-sheet-7-018
% page_slug: formula-sheet-7
% page_title: Formula Sheet
% page_url: https://formulasheet.com/#q|l|1022
% subject_slug: formulasheet
% subject_path: /formulasheet/
% section: Integral Table [integral-table.com]
% subsection: 
% formula_name: 
% source: formulasheet-latex
% index_global: 193
% index_in_page: 18
\[\begin{align} \int \cos ax \sinh bx dx& = \frac{1}{a^2 + b^2} \left[ b \cos ax \cosh bx + \right . \nonumber \\\ & \left . a \sin ax \sinh bx \right] \end{align}\]
% formula_record_end

% formula_record_start
% formula_id: formula-sheet-7-019
% page_slug: formula-sheet-7
% page_title: Formula Sheet
% page_url: https://formulasheet.com/#q|l|1022
% subject_slug: formulasheet
% subject_path: /formulasheet/
% section: Integral Table [integral-table.com]
% subsection: 
% formula_name: 
% source: formulasheet-latex
% index_global: 194
% index_in_page: 19
\[\begin{align} \int \cos^2 ax \sin bx dx &= \frac{\cos[(2a-b)x]}{4(2a-b)} - \frac{\cos bx}{2b} \nonumber \\\ & - \frac{\cos[(2a+b)x]}{4(2a+b)} \end{align}\]
% formula_record_end

% formula_record_start
% formula_id: formula-sheet-7-020
% page_slug: formula-sheet-7
% page_title: Formula Sheet
% page_url: https://formulasheet.com/#q|l|1022
% subject_slug: formulasheet
% subject_path: /formulasheet/
% section: Integral Table [integral-table.com]
% subsection: 
% formula_name: 
% source: formulasheet-latex
% index_global: 195
% index_in_page: 20
\[\begin{align} \int \cos^p ax dx & = -\frac{1}{a(1+p)}{\cos^{1+p} ax} \times \nonumber \\\ & {_2F_1}\left[ \frac{1+p}{2}, \frac{1}{2}, \frac{3+p}{2}, \cos^2 ax \right] \end{align}\]
% formula_record_end

% formula_record_start
% formula_id: formula-sheet-7-021
% page_slug: formula-sheet-7
% page_title: Formula Sheet
% page_url: https://formulasheet.com/#q|l|1022
% subject_slug: formulasheet
% subject_path: /formulasheet/
% section: Integral Table [integral-table.com]
% subsection: 
% formula_name: 
% source: formulasheet-latex
% index_global: 196
% index_in_page: 21
\[\begin{align} \int \frac{x^2}{\sqrt{x^2 \pm a^2}} dx &= \frac{1}{2}x\sqrt{x^2 \pm a^2} %\nonumber \\\ & \mp \frac{1}{2}a^2 \ln \left| x + \sqrt{x^2\pm a^2} \right | \end{align}\]
% formula_record_end

% formula_record_start
% formula_id: formula-sheet-7-022
% page_slug: formula-sheet-7
% page_title: Formula Sheet
% page_url: https://formulasheet.com/#q|l|1022
% subject_slug: formulasheet
% subject_path: /formulasheet/
% section: Integral Table [integral-table.com]
% subsection: 
% formula_name: 
% source: formulasheet-latex
% index_global: 197
% index_in_page: 22
\[\begin{align} \int \frac{x}{ax^2+bx+c}dx &= \frac{1}{2a}\ln|ax^2+bx+c| \nonumber \\&-\frac{b}{a\sqrt{4ac-b^2}}\tan^{-1}\frac{2ax+b}{\sqrt{4ac-b^2}} \end{align}\]
% formula_record_end

% formula_record_start
% formula_id: formula-sheet-7-023
% page_slug: formula-sheet-7
% page_title: Formula Sheet
% page_url: https://formulasheet.com/#q|l|1022
% subject_slug: formulasheet
% subject_path: /formulasheet/
% section: Integral Table [integral-table.com]
% subsection: 
% formula_name: 
% source: formulasheet-latex
% index_global: 198
% index_in_page: 23
\[\begin{align} \int \ln ( x^2 + a^2 )\hspace{.5ex}\text{dx} = x \ln (x^2 + a^2 ) +2a\tan^{-1} \frac{x}{a} - 2x \end{align}\]
% formula_record_end

% formula_record_start
% formula_id: formula-sheet-7-024
% page_slug: formula-sheet-7
% page_title: Formula Sheet
% page_url: https://formulasheet.com/#q|l|1022
% subject_slug: formulasheet
% subject_path: /formulasheet/
% section: Integral Table [integral-table.com]
% subsection: 
% formula_name: 
% source: formulasheet-latex
% index_global: 199
% index_in_page: 24
\[\begin{align} \int \ln ( x^2 - a^2 )\hspace{.5ex}\text{dx} = x \ln (x^2 - a^2 ) +a\ln \frac{x+a}{x-a} - 2x \end{align}\]
% formula_record_end

% formula_record_start
% formula_id: formula-sheet-7-025
% page_slug: formula-sheet-7
% page_title: Formula Sheet
% page_url: https://formulasheet.com/#q|l|1022
% subject_slug: formulasheet
% subject_path: /formulasheet/
% section: Integral Table [integral-table.com]
% subsection: 
% formula_name: 
% source: formulasheet-latex
% index_global: 200
% index_in_page: 25
\[\begin{align} \int \ln & \left ( ax^2 + bx + c\right) dx = \frac{1}{a}\sqrt{4ac-b^2}\tan^{-1}\frac{2ax+b}{\sqrt{4ac-b^2}} \nonumber \\\ & -2x + \left( \frac{b}{2a}+x \right )\ln \left (ax^2+bx+c \right) \end{align}\]
% formula_record_end

% formula_record_start
% formula_id: formula-sheet-7-026
% page_slug: formula-sheet-7
% page_title: Formula Sheet
% page_url: https://formulasheet.com/#q|l|1022
% subject_slug: formulasheet
% subject_path: /formulasheet/
% section: Integral Table [integral-table.com]
% subsection: 
% formula_name: 
% source: formulasheet-latex
% index_global: 201
% index_in_page: 26
\[\begin{align} \int \sec x dx &= \ln | \sec x + \tan x | = 2 \tanh^{-1} \left (\tan \frac{x}{2} \right) \end{align}\]
% formula_record_end

% formula_record_start
% formula_id: formula-sheet-7-027
% page_slug: formula-sheet-7
% page_title: Formula Sheet
% page_url: https://formulasheet.com/#q|l|1022
% subject_slug: formulasheet
% subject_path: /formulasheet/
% section: Integral Table [integral-table.com]
% subsection: 
% formula_name: 
% source: formulasheet-latex
% index_global: 202
% index_in_page: 27
\[\begin{align} \int \sin ax \cosh bx dx &= \frac{1}{a^2 + b^2} \left[ -a \cos ax \cosh bx + \right . \nonumber \\\ & \left . b \sin ax \sinh bx \right] \end{align}\]
% formula_record_end

% formula_record_start
% formula_id: formula-sheet-7-028
% page_slug: formula-sheet-7
% page_title: Formula Sheet
% page_url: https://formulasheet.com/#q|l|1022
% subject_slug: formulasheet
% subject_path: /formulasheet/
% section: Integral Table [integral-table.com]
% subsection: 
% formula_name: 
% source: formulasheet-latex
% index_global: 203
% index_in_page: 28
\[\begin{align} \int \sin ax \sinh bx dx &= \frac{1}{a^2 + b^2} \left[ b \cosh bx \sin ax - \right . \nonumber \\\ & \left . a \cos ax \sinh bx \right] \end{align}\]
% formula_record_end

% formula_record_start
% formula_id: formula-sheet-7-029
% page_slug: formula-sheet-7
% page_title: Formula Sheet
% page_url: https://formulasheet.com/#q|l|1022
% subject_slug: formulasheet
% subject_path: /formulasheet/
% section: Integral Table [integral-table.com]
% subsection: 
% formula_name: 
% source: formulasheet-latex
% index_global: 204
% index_in_page: 29
\[\begin{align} \int \sinh ax \cosh bx dx&= \frac{1}{b^2-a^2}\left[ b \cosh bx \sinh ax \right . \nonumber \\\ & \left . - a \cosh ax \sinh bx \right] \end{align}\]
% formula_record_end

% formula_record_start
% formula_id: formula-sheet-7-030
% page_slug: formula-sheet-7
% page_title: Formula Sheet
% page_url: https://formulasheet.com/#q|l|1022
% subject_slug: formulasheet
% subject_path: /formulasheet/
% section: Integral Table [integral-table.com]
% subsection: 
% formula_name: 
% source: formulasheet-latex
% index_global: 205
% index_in_page: 30
\[\begin{align} \int \sin^2 ax \cos bx dx &= -\frac{\sin[(2a-b)x]}{4(2a-b)} \nonumber \\\ & + \frac{\sin bx}{2b} - \frac{\sin[(2a+b)x]}{4(2a+b)} \end{align}\]
% formula_record_end

% formula_record_start
% formula_id: formula-sheet-7-031
% page_slug: formula-sheet-7
% page_title: Formula Sheet
% page_url: https://formulasheet.com/#q|l|1022
% subject_slug: formulasheet
% subject_path: /formulasheet/
% section: Integral Table [integral-table.com]
% subsection: 
% formula_name: 
% source: formulasheet-latex
% index_global: 206
% index_in_page: 31
\[\begin{align} \int \sin^2 ax \cos^2 bx dx &= \frac{x}{4} -\frac{\sin 2ax}{8a}- \frac{\sin[2(a-b)x]}{16(a-b)} \nonumber \\\ & +\frac{\sin 2bx}{8b}- \frac{\sin[2(a+b)x]}{16(a+b)} \end{align}\]
% formula_record_end

% formula_record_start
% formula_id: formula-sheet-7-032
% page_slug: formula-sheet-7
% page_title: Formula Sheet
% page_url: https://formulasheet.com/#q|l|1022
% subject_slug: formulasheet
% subject_path: /formulasheet/
% section: Integral Table [integral-table.com]
% subsection: 
% formula_name: 
% source: formulasheet-latex
% index_global: 207
% index_in_page: 32
\[\begin{align} \int \sqrt{x(ax+b)} dx &= \frac{1}{4a^{3/2}}\left[(2ax + b)\sqrt{ax(ax+b)} \right. \nonumber \\& \left. -b^2 \ln \left| a\sqrt{x} + \sqrt{a(ax+b)} \right| \right ] \end{align}\]
% formula_record_end

% formula_record_start
% formula_id: formula-sheet-7-033
% page_slug: formula-sheet-7
% page_title: Formula Sheet
% page_url: https://formulasheet.com/#q|l|1022
% subject_slug: formulasheet
% subject_path: /formulasheet/
% section: Integral Table [integral-table.com]
% subsection: 
% formula_name: 
% source: formulasheet-latex
% index_global: 208
% index_in_page: 33
\[\begin{align} \int \sqrt{x^3(ax+b)} dx &=\left [ \frac{b}{12a}- \frac{b^2}{8a^2x}+ \frac{x}{3}\right] \sqrt{x^3(ax+b)} \nonumber \\& + \frac{b^3}{8a^{5/2}}\ln \left | a\sqrt{x} + \sqrt{a(ax+b)} \right | \end{align}\]
% formula_record_end

% formula_record_start
% formula_id: formula-sheet-7-034
% page_slug: formula-sheet-7
% page_title: Formula Sheet
% page_url: https://formulasheet.com/#q|l|1022
% subject_slug: formulasheet
% subject_path: /formulasheet/
% section: Integral Table [integral-table.com]
% subsection: 
% formula_name: 
% source: formulasheet-latex
% index_global: 209
% index_in_page: 34
\[\begin{align} \int \sqrt{x} e^{ax} dx &= \frac{1}{a}\sqrt{x}e^{ax} +\frac{i\sqrt{\pi}}{2a^{3/2}} \text{erf}\left(i\sqrt{ax}\right), \nonumber \\& \text{ where erf}(x)=\frac{2}{\sqrt{\pi}}\int_0^x e^{-t^2}dt \end{align}\]
% formula_record_end

% formula_record_start
% formula_id: formula-sheet-7-035
% page_slug: formula-sheet-7
% page_title: Formula Sheet
% page_url: https://formulasheet.com/#q|l|1022
% subject_slug: formulasheet
% subject_path: /formulasheet/
% section: Integral Table [integral-table.com]
% subsection: 
% formula_name: 
% source: formulasheet-latex
% index_global: 210
% index_in_page: 35
\[\begin{align} \int \sqrt{\frac{x}{a+x}}dx &= \sqrt{x(a+x)} %\nonumber \\& -a\ln \left [ \sqrt{x} + \sqrt{x+a}\right] \end{align}\]
% formula_record_end

% formula_record_start
% formula_id: formula-sheet-7-036
% page_slug: formula-sheet-7
% page_title: Formula Sheet
% page_url: https://formulasheet.com/#q|l|1022
% subject_slug: formulasheet
% subject_path: /formulasheet/
% section: Integral Table [integral-table.com]
% subsection: 
% formula_name: 
% source: formulasheet-latex
% index_global: 211
% index_in_page: 36
\[\begin{align} \int \sqrt{\frac{x}{a-x}}dx &= -\sqrt{x(a-x)} %\nonumber \\& -a\tan^{-1}\frac{\sqrt{x(a-x)}}{x-a} \end{align}\]
% formula_record_end

% formula_record_start
% formula_id: formula-sheet-7-037
% page_slug: formula-sheet-7
% page_title: Formula Sheet
% page_url: https://formulasheet.com/#q|l|1022
% subject_slug: formulasheet
% subject_path: /formulasheet/
% section: Integral Table [integral-table.com]
% subsection: 
% formula_name: 
% source: formulasheet-latex
% index_global: 212
% index_in_page: 37
\[\begin{align} \int&\frac{1}{\sqrt{ax^2+bx+c}}dx= %==\nonumber \\& \frac{1}{\sqrt{a}}\ln \left| 2ax+b + 2 \sqrt{a(ax^2+bx+c)} \right | \end{align}\]
% formula_record_end

% formula_record_start
% formula_id: formula-sheet-7-038
% page_slug: formula-sheet-7
% page_title: Formula Sheet
% page_url: https://formulasheet.com/#q|l|1022
% subject_slug: formulasheet
% subject_path: /formulasheet/
% section: Integral Table [integral-table.com]
% subsection: 
% formula_name: 
% source: formulasheet-latex
% index_global: 213
% index_in_page: 38
\[\begin{align} \int\sqrt{x^2 \pm a^2} dx &= \frac{1}{2}x\sqrt{x^2\pm a^2} %\nonumber \\& \pm\frac{1}{2}a^2 \ln \left | x + \sqrt{x^2\pm a^2} \right | \end{align}\]
% formula_record_end

% formula_record_start
% formula_id: formula-sheet-7-039
% page_slug: formula-sheet-7
% page_title: Formula Sheet
% page_url: https://formulasheet.com/#q|l|1022
% subject_slug: formulasheet
% subject_path: /formulasheet/
% section: Integral Table [integral-table.com]
% subsection: 
% formula_name: 
% source: formulasheet-latex
% index_global: 214
% index_in_page: 39
\[\begin{equation} \begin{split} \int x^n e^{ax}\hspace{2pt}\text{d}x = \frac{(-1)^n}{a^{n+1}}\Gamma[1+n,-ax], \\\ \text{ where } \Gamma(a,x)=\int_x^{\infty} t^{a-1}e^{-t}\hspace{2pt}\text{d}t \end{split} \end{equation}\]
% formula_record_end

% formula_record_start
% formula_id: formula-sheet-7-040
% page_slug: formula-sheet-7
% page_title: Formula Sheet
% page_url: https://formulasheet.com/#q|l|1022
% subject_slug: formulasheet
% subject_path: /formulasheet/
% section: Integral Table [integral-table.com]
% subsection: 
% formula_name: 
% source: formulasheet-latex
% index_global: 215
% index_in_page: 40
\[\begin{equation} \int x \sqrt{x^2 \pm a^2} dx= \frac{1}{3}\left ( x^2 \pm a^2 \right)^{3/2} \end{equation}\]
% formula_record_end

% formula_record_start
% formula_id: formula-sheet-7-041
% page_slug: formula-sheet-7
% page_title: Formula Sheet
% page_url: https://formulasheet.com/#q|l|1022
% subject_slug: formulasheet
% subject_path: /formulasheet/
% section: Integral Table [integral-table.com]
% subsection: 
% formula_name: 
% source: formulasheet-latex
% index_global: 216
% index_in_page: 41
\[\begin{equation} \int \tanh ax\hspace{1.5pt} dx =\frac{1}{a} \ln \cosh ax \end{equation}\]
% formula_record_end

% formula_record_start
% formula_id: formula-sheet-7-042
% page_slug: formula-sheet-7
% page_title: Formula Sheet
% page_url: https://formulasheet.com/#q|l|1022
% subject_slug: formulasheet
% subject_path: /formulasheet/
% section: Integral Table [integral-table.com]
% subsection: 
% formula_name: 
% source: formulasheet-latex
% index_global: 217
% index_in_page: 42
\[\begin{equation} \int (ax+b)^{3/2}dx =\frac{2}{5a}(ax+b)^{5/2} \end{equation}\]
% formula_record_end

% formula_record_start
% formula_id: formula-sheet-7-043
% page_slug: formula-sheet-7
% page_title: Formula Sheet
% page_url: https://formulasheet.com/#q|l|1022
% subject_slug: formulasheet
% subject_path: /formulasheet/
% section: Integral Table [integral-table.com]
% subsection: 
% formula_name: 
% source: formulasheet-latex
% index_global: 218
% index_in_page: 43
\[\begin{equation} \int (x+a)^n dx = \frac{(x+a)^{n+1}}{n+1}+c, n\ne -1 \end{equation}\]
% formula_record_end

% formula_record_start
% formula_id: formula-sheet-7-044
% page_slug: formula-sheet-7
% page_title: Formula Sheet
% page_url: https://formulasheet.com/#q|l|1022
% subject_slug: formulasheet
% subject_path: /formulasheet/
% section: Integral Table [integral-table.com]
% subsection: 
% formula_name: 
% source: formulasheet-latex
% index_global: 219
% index_in_page: 44
\[\begin{equation} \int e^x \cos x dx = \frac{1}{2}e^x (\sin x + \cos x) \end{equation}\]
% formula_record_end

% formula_record_start
% formula_id: formula-sheet-7-045
% page_slug: formula-sheet-7
% page_title: Formula Sheet
% page_url: https://formulasheet.com/#q|l|1022
% subject_slug: formulasheet
% subject_path: /formulasheet/
% section: Integral Table [integral-table.com]
% subsection: 
% formula_name: 
% source: formulasheet-latex
% index_global: 220
% index_in_page: 45
\[\begin{equation} \int e^x \sin x dx = \frac{1}{2}e^x (\sin x - \cos x) \end{equation}\]
% formula_record_end

% formula_record_start
% formula_id: formula-sheet-7-046
% page_slug: formula-sheet-7
% page_title: Formula Sheet
% page_url: https://formulasheet.com/#q|l|1022
% subject_slug: formulasheet
% subject_path: /formulasheet/
% section: Integral Table [integral-table.com]
% subsection: 
% formula_name: 
% source: formulasheet-latex
% index_global: 221
% index_in_page: 46
\[\begin{equation} \int e^{-ax^2}\hspace{1pt}\text{d}x = \frac{\sqrt{\pi}}{2\sqrt{a}}\text{erf}\left(x\sqrt{a}\right) \end{equation}\]
% formula_record_end

% formula_record_start
% formula_id: formula-sheet-7-047
% page_slug: formula-sheet-7
% page_title: Formula Sheet
% page_url: https://formulasheet.com/#q|l|1022
% subject_slug: formulasheet
% subject_path: /formulasheet/
% section: Integral Table [integral-table.com]
% subsection: 
% formula_name: 
% source: formulasheet-latex
% index_global: 222
% index_in_page: 47
\[\begin{equation} \int e^{ax^2}\hspace{1pt}\text{d}x = -\frac{i\sqrt{\pi}}{2\sqrt{a}}\text{erf}\left(ix\sqrt{a}\right) \end{equation}\]
% formula_record_end

% formula_record_start
% formula_id: formula-sheet-7-048
% page_slug: formula-sheet-7
% page_title: Formula Sheet
% page_url: https://formulasheet.com/#q|l|1022
% subject_slug: formulasheet
% subject_path: /formulasheet/
% section: Integral Table [integral-table.com]
% subsection: 
% formula_name: 
% source: formulasheet-latex
% index_global: 223
% index_in_page: 48
\[\begin{equation} \int e^{ax} dx = \frac{1}{a}e^{ax} \end{equation}\]
% formula_record_end

% formula_record_start
% formula_id: formula-sheet-7-049
% page_slug: formula-sheet-7
% page_title: Formula Sheet
% page_url: https://formulasheet.com/#q|l|1022
% subject_slug: formulasheet
% subject_path: /formulasheet/
% section: Integral Table [integral-table.com]
% subsection: 
% formula_name: 
% source: formulasheet-latex
% index_global: 224
% index_in_page: 49
\[\begin{equation} \int e^{bx} \cos ax dx = \frac{1}{a^2 + b^2} e^{bx} ( a \sin ax + b \cos ax ) \end{equation}\]
% formula_record_end

% formula_record_start
% formula_id: formula-sheet-7-050
% page_slug: formula-sheet-7
% page_title: Formula Sheet
% page_url: https://formulasheet.com/#q|l|1022
% subject_slug: formulasheet
% subject_path: /formulasheet/
% section: Integral Table [integral-table.com]
% subsection: 
% formula_name: 
% source: formulasheet-latex
% index_global: 225
% index_in_page: 50
\[\begin{equation} \int e^{bx} \sin ax dx = \frac{1}{a^2+b^2}e^{bx} (b\sin ax - a\cos ax) \end{equation}\]
% formula_record_end

% formula_record_start
% formula_id: formula-sheet-7-051
% page_slug: formula-sheet-7
% page_title: Formula Sheet
% page_url: https://formulasheet.com/#q|l|1022
% subject_slug: formulasheet
% subject_path: /formulasheet/
% section: Integral Table [integral-table.com]
% subsection: 
% formula_name: 
% source: formulasheet-latex
% index_global: 226
% index_in_page: 51
\[\begin{equation} \int u dv = uv - \int v du \end{equation}\]
% formula_record_end

% formula_record_start
% formula_id: formula-sheet-7-052
% page_slug: formula-sheet-7
% page_title: Formula Sheet
% page_url: https://formulasheet.com/#q|l|1022
% subject_slug: formulasheet
% subject_path: /formulasheet/
% section: Integral Table [integral-table.com]
% subsection: 
% formula_name: 
% source: formulasheet-latex
% index_global: 227
% index_in_page: 52
\[\begin{equation} \int x e^x dx = (x-1) e^x \end{equation}\]
% formula_record_end

% formula_record_start
% formula_id: formula-sheet-7-053
% page_slug: formula-sheet-7
% page_title: Formula Sheet
% page_url: https://formulasheet.com/#q|l|1022
% subject_slug: formulasheet
% subject_path: /formulasheet/
% section: Integral Table [integral-table.com]
% subsection: 
% formula_name: 
% source: formulasheet-latex
% index_global: 228
% index_in_page: 53
\[\begin{equation} \int x e^x \cos x dx = \frac{1}{2}e^x (x \cos x - \sin x + x \sin x) \end{equation}\]
% formula_record_end

% formula_record_start
% formula_id: formula-sheet-7-054
% page_slug: formula-sheet-7
% page_title: Formula Sheet
% page_url: https://formulasheet.com/#q|l|1022
% subject_slug: formulasheet
% subject_path: /formulasheet/
% section: Integral Table [integral-table.com]
% subsection: 
% formula_name: 
% source: formulasheet-latex
% index_global: 229
% index_in_page: 54
\[\begin{equation} \int x e^x \sin x dx = \frac{1}{2}e^x (\cos x - x \cos x + x \sin x) \end{equation}\]
% formula_record_end

% formula_record_start
% formula_id: formula-sheet-7-055
% page_slug: formula-sheet-7
% page_title: Formula Sheet
% page_url: https://formulasheet.com/#q|l|1022
% subject_slug: formulasheet
% subject_path: /formulasheet/
% section: Integral Table [integral-table.com]
% subsection: 
% formula_name: 
% source: formulasheet-latex
% index_global: 230
% index_in_page: 55
\[\begin{equation} \int x e^{-ax^2}\\ \text{dx} = -\dfrac{1}{2a}e^{-ax^2} \end{equation}\]
% formula_record_end

% formula_record_start
% formula_id: formula-sheet-7-056
% page_slug: formula-sheet-7
% page_title: Formula Sheet
% page_url: https://formulasheet.com/#q|l|1022
% subject_slug: formulasheet
% subject_path: /formulasheet/
% section: Integral Table [integral-table.com]
% subsection: 
% formula_name: 
% source: formulasheet-latex
% index_global: 231
% index_in_page: 56
\[\begin{equation} \int x e^{ax} dx = \left(\frac{x}{a}-\frac{1}{a^2}\right) e^{ax} \end{equation}\]
% formula_record_end

% formula_record_start
% formula_id: formula-sheet-7-057
% page_slug: formula-sheet-7
% page_title: Formula Sheet
% page_url: https://formulasheet.com/#q|l|1022
% subject_slug: formulasheet
% subject_path: /formulasheet/
% section: Integral Table [integral-table.com]
% subsection: 
% formula_name: 
% source: formulasheet-latex
% index_global: 232
% index_in_page: 57
\[\begin{equation} \int x \cos ax dx = \frac{1}{a^2} \cos ax + \frac{x}{a} \sin ax \end{equation}\]
% formula_record_end

% formula_record_start
% formula_id: formula-sheet-7-058
% page_slug: formula-sheet-7
% page_title: Formula Sheet
% page_url: https://formulasheet.com/#q|l|1022
% subject_slug: formulasheet
% subject_path: /formulasheet/
% section: Integral Table [integral-table.com]
% subsection: 
% formula_name: 
% source: formulasheet-latex
% index_global: 233
% index_in_page: 58
\[\begin{equation} \int x \cos x dx = \cos x + x \sin x \end{equation}\]
% formula_record_end

% formula_record_start
% formula_id: formula-sheet-7-059
% page_slug: formula-sheet-7
% page_title: Formula Sheet
% page_url: https://formulasheet.com/#q|l|1022
% subject_slug: formulasheet
% subject_path: /formulasheet/
% section: Integral Table [integral-table.com]
% subsection: 
% formula_name: 
% source: formulasheet-latex
% index_global: 234
% index_in_page: 59
\[\begin{equation} \int x \sin ax dx = -\frac{x \cos ax}{a} + \frac{\sin ax}{a^2} \end{equation}\]
% formula_record_end

% formula_record_start
% formula_id: formula-sheet-7-060
% page_slug: formula-sheet-7
% page_title: Formula Sheet
% page_url: https://formulasheet.com/#q|l|1022
% subject_slug: formulasheet
% subject_path: /formulasheet/
% section: Integral Table [integral-table.com]
% subsection: 
% formula_name: 
% source: formulasheet-latex
% index_global: 235
% index_in_page: 60
\[\begin{equation} \int x \sin x dx = -x \cos x + \sin x \end{equation}\]
% formula_record_end

% formula_record_start
% formula_id: formula-sheet-7-061
% page_slug: formula-sheet-7
% page_title: Formula Sheet
% page_url: https://formulasheet.com/#q|l|1022
% subject_slug: formulasheet
% subject_path: /formulasheet/
% section: Integral Table [integral-table.com]
% subsection: 
% formula_name: 
% source: formulasheet-latex
% index_global: 236
% index_in_page: 61
\[\begin{equation} \int x(x+a)^n dx = \frac{(x+a)^{n+1} ( (n+1)x-a)}{(n+1)(n+2)} \end{equation}\]
% formula_record_end

% formula_record_start
% formula_id: formula-sheet-7-062
% page_slug: formula-sheet-7
% page_title: Formula Sheet
% page_url: https://formulasheet.com/#q|l|1022
% subject_slug: formulasheet
% subject_path: /formulasheet/
% section: Integral Table [integral-table.com]
% subsection: 
% formula_name: 
% source: formulasheet-latex
% index_global: 237
% index_in_page: 62
\[\begin{equation} \int x\sqrt{x-a} dx = \frac{2}{3}a(x-a)^{3/2}+\frac{2}{5}(x-a)^{5/2} \end{equation}\]
% formula_record_end

% formula_record_start
% formula_id: formula-sheet-7-063
% page_slug: formula-sheet-7
% page_title: Formula Sheet
% page_url: https://formulasheet.com/#q|l|1022
% subject_slug: formulasheet
% subject_path: /formulasheet/
% section: Integral Table [integral-table.com]
% subsection: 
% formula_name: 
% source: formulasheet-latex
% index_global: 238
% index_in_page: 63
\[\begin{equation} \int x^2 e^{-ax^2}\\ \text{dx} = \dfrac{1}{4}\sqrt{\dfrac{\pi}{a^3}}\text{erf}(x\sqrt{a}) -\dfrac{x}{2a}e^{-ax^2} \end{equation}\]
% formula_record_end

% formula_record_start
% formula_id: formula-sheet-7-064
% page_slug: formula-sheet-7
% page_title: Formula Sheet
% page_url: https://formulasheet.com/#q|l|1022
% subject_slug: formulasheet
% subject_path: /formulasheet/
% section: Integral Table [integral-table.com]
% subsection: 
% formula_name: 
% source: formulasheet-latex
% index_global: 239
% index_in_page: 64
\[\begin{equation} \int x^2 e^{ax} dx = \left(\frac{x^2}{a}-\frac{2x}{a^2}+\frac{2}{a^3}\right) e^{ax} \end{equation}\]
% formula_record_end

% formula_record_start
% formula_id: formula-sheet-7-065
% page_slug: formula-sheet-7
% page_title: Formula Sheet
% page_url: https://formulasheet.com/#q|l|1022
% subject_slug: formulasheet
% subject_path: /formulasheet/
% section: Integral Table [integral-table.com]
% subsection: 
% formula_name: 
% source: formulasheet-latex
% index_global: 240
% index_in_page: 65
\[\begin{equation} \int x^2 e^{x} dx = \left(x^2 - 2x + 2\right) e^{x} \end{equation}\]
% formula_record_end

% formula_record_start
% formula_id: formula-sheet-7-066
% page_slug: formula-sheet-7
% page_title: Formula Sheet
% page_url: https://formulasheet.com/#q|l|1022
% subject_slug: formulasheet
% subject_path: /formulasheet/
% section: Integral Table [integral-table.com]
% subsection: 
% formula_name: 
% source: formulasheet-latex
% index_global: 241
% index_in_page: 66
\[\begin{equation} \int x^2 \cos ax dx = \frac{2 x \cos ax }{a^2} + \frac{ a^2 x^2 - 2 }{a^3} \sin ax \end{equation}\]
% formula_record_end

% formula_record_start
% formula_id: formula-sheet-7-067
% page_slug: formula-sheet-7
% page_title: Formula Sheet
% page_url: https://formulasheet.com/#q|l|1022
% subject_slug: formulasheet
% subject_path: /formulasheet/
% section: Integral Table [integral-table.com]
% subsection: 
% formula_name: 
% source: formulasheet-latex
% index_global: 242
% index_in_page: 67
\[\begin{equation} \int x^2 \cos x dx = 2 x \cos x + \left ( x^2 - 2 \right ) \sin x \end{equation}\]
% formula_record_end

% formula_record_start
% formula_id: formula-sheet-7-068
% page_slug: formula-sheet-7
% page_title: Formula Sheet
% page_url: https://formulasheet.com/#q|l|1022
% subject_slug: formulasheet
% subject_path: /formulasheet/
% section: Integral Table [integral-table.com]
% subsection: 
% formula_name: 
% source: formulasheet-latex
% index_global: 243
% index_in_page: 68
\[\begin{equation} \int x^2 \sin ax dx =\frac{2-a^2x^2}{a^3}\cos ax +\frac{ 2 x \sin ax}{a^2} \end{equation}\]
% formula_record_end

% formula_record_start
% formula_id: formula-sheet-7-069
% page_slug: formula-sheet-7
% page_title: Formula Sheet
% page_url: https://formulasheet.com/#q|l|1022
% subject_slug: formulasheet
% subject_path: /formulasheet/
% section: Integral Table [integral-table.com]
% subsection: 
% formula_name: 
% source: formulasheet-latex
% index_global: 244
% index_in_page: 69
\[\begin{equation} \int x^2 \sin x dx = \left(2-x^2\right) \cos x + 2 x \sin x \end{equation}\]
% formula_record_end

% formula_record_start
% formula_id: formula-sheet-7-070
% page_slug: formula-sheet-7
% page_title: Formula Sheet
% page_url: https://formulasheet.com/#q|l|1022
% subject_slug: formulasheet
% subject_path: /formulasheet/
% section: Integral Table [integral-table.com]
% subsection: 
% formula_name: 
% source: formulasheet-latex
% index_global: 245
% index_in_page: 70
\[\begin{equation} \int x^3 e^{x} dx = \left(x^3-3x^2 + 6x - 6\right) e^{x} \end{equation}\]
% formula_record_end

% formula_record_start
% formula_id: formula-sheet-7-071
% page_slug: formula-sheet-7
% page_title: Formula Sheet
% page_url: https://formulasheet.com/#q|l|1022
% subject_slug: formulasheet
% subject_path: /formulasheet/
% section: Integral Table [integral-table.com]
% subsection: 
% formula_name: 
% source: formulasheet-latex
% index_global: 246
% index_in_page: 71
\[\begin{equation} \int x^n dx = \frac{1}{n+1}x^{n+1} \end{equation}\]
% formula_record_end

% formula_record_start
% formula_id: formula-sheet-7-072
% page_slug: formula-sheet-7
% page_title: Formula Sheet
% page_url: https://formulasheet.com/#q|l|1022
% subject_slug: formulasheet
% subject_path: /formulasheet/
% section: Integral Table [integral-table.com]
% subsection: 
% formula_name: 
% source: formulasheet-latex
% index_global: 247
% index_in_page: 72
\[\begin{equation} \int x^n e^{ax}\hspace{1pt}\text{d}x = \dfrac{x^n e^{ax}}{a} - \dfrac{n}{a}\int x^{n-1}e^{ax}\hspace{1pt}\text{d}x \end{equation}\]
% formula_record_end

% formula_record_start
% formula_id: formula-sheet-7-073
% page_slug: formula-sheet-7
% page_title: Formula Sheet
% page_url: https://formulasheet.com/#q|l|1022
% subject_slug: formulasheet
% subject_path: /formulasheet/
% section: Integral Table [integral-table.com]
% subsection: 
% formula_name: 
% source: formulasheet-latex
% index_global: 248
% index_in_page: 73
\[\begin{equation} \int \cos ax dx= \frac{1}{a} \sin ax \end{equation}\]
% formula_record_end

% formula_record_start
% formula_id: formula-sheet-7-074
% page_slug: formula-sheet-7
% page_title: Formula Sheet
% page_url: https://formulasheet.com/#q|l|1022
% subject_slug: formulasheet
% subject_path: /formulasheet/
% section: Integral Table [integral-table.com]
% subsection: 
% formula_name: 
% source: formulasheet-latex
% index_global: 249
% index_in_page: 74
\[\begin{equation} \int \cosh ax dx =\frac{1}{a} \sinh ax \end{equation}\]
% formula_record_end

% formula_record_start
% formula_id: formula-sheet-7-075
% page_slug: formula-sheet-7
% page_title: Formula Sheet
% page_url: https://formulasheet.com/#q|l|1022
% subject_slug: formulasheet
% subject_path: /formulasheet/
% section: Integral Table [integral-table.com]
% subsection: 
% formula_name: 
% source: formulasheet-latex
% index_global: 250
% index_in_page: 75
\[\begin{equation} \int \cos^2 ax dx = \frac{x}{2}+\frac{ \sin 2ax}{4a} \end{equation}\]
% formula_record_end

% formula_record_start
% formula_id: formula-sheet-7-076
% page_slug: formula-sheet-7
% page_title: Formula Sheet
% page_url: https://formulasheet.com/#q|l|1022
% subject_slug: formulasheet
% subject_path: /formulasheet/
% section: Integral Table [integral-table.com]
% subsection: 
% formula_name: 
% source: formulasheet-latex
% index_global: 251
% index_in_page: 76
\[\begin{equation} \int \cos^2 ax \sin ax dx = -\frac{1}{3a}\cos^3{ax} \end{equation}\]
% formula_record_end

% formula_record_start
% formula_id: formula-sheet-7-077
% page_slug: formula-sheet-7
% page_title: Formula Sheet
% page_url: https://formulasheet.com/#q|l|1022
% subject_slug: formulasheet
% subject_path: /formulasheet/
% section: Integral Table [integral-table.com]
% subsection: 
% formula_name: 
% source: formulasheet-latex
% index_global: 252
% index_in_page: 77
\[\begin{equation} \int \cos^3 ax dx = \frac{3 \sin ax}{4a}+\frac{ \sin 3ax}{12a} \end{equation}\]
% formula_record_end

% formula_record_start
% formula_id: formula-sheet-7-078
% page_slug: formula-sheet-7
% page_title: Formula Sheet
% page_url: https://formulasheet.com/#q|l|1022
% subject_slug: formulasheet
% subject_path: /formulasheet/
% section: Integral Table [integral-table.com]
% subsection: 
% formula_name: 
% source: formulasheet-latex
% index_global: 253
% index_in_page: 78
\[\begin{equation} \int \csc x dx = \ln \left | \tan \frac{x}{2} \right| = \ln | \csc x - \cot x| + C \end{equation}\]
% formula_record_end

% formula_record_start
% formula_id: formula-sheet-7-079
% page_slug: formula-sheet-7
% page_title: Formula Sheet
% page_url: https://formulasheet.com/#q|l|1022
% subject_slug: formulasheet
% subject_path: /formulasheet/
% section: Integral Table [integral-table.com]
% subsection: 
% formula_name: 
% source: formulasheet-latex
% index_global: 254
% index_in_page: 79
\[\begin{equation} \int \csc^2 ax dx = -\frac{1}{a} \cot ax \end{equation}\]
% formula_record_end

% formula_record_start
% formula_id: formula-sheet-7-080
% page_slug: formula-sheet-7
% page_title: Formula Sheet
% page_url: https://formulasheet.com/#q|l|1022
% subject_slug: formulasheet
% subject_path: /formulasheet/
% section: Integral Table [integral-table.com]
% subsection: 
% formula_name: 
% source: formulasheet-latex
% index_global: 255
% index_in_page: 80
\[\begin{equation} \int \csc^3 x dx = -\frac{1}{2}\cot x \csc x + \frac{1}{2} \ln | \csc x - \cot x | \end{equation}\]
% formula_record_end

% formula_record_start
% formula_id: formula-sheet-7-081
% page_slug: formula-sheet-7
% page_title: Formula Sheet
% page_url: https://formulasheet.com/#q|l|1022
% subject_slug: formulasheet
% subject_path: /formulasheet/
% section: Integral Table [integral-table.com]
% subsection: 
% formula_name: 
% source: formulasheet-latex
% index_global: 256
% index_in_page: 81
\[\begin{equation} \int \csc^nx \cot x dx = -\frac{1}{n}\csc^n x, n\ne 0 \end{equation}\]
% formula_record_end

% formula_record_start
% formula_id: formula-sheet-7-082
% page_slug: formula-sheet-7
% page_title: Formula Sheet
% page_url: https://formulasheet.com/#q|l|1022
% subject_slug: formulasheet
% subject_path: /formulasheet/
% section: Integral Table [integral-table.com]
% subsection: 
% formula_name: 
% source: formulasheet-latex
% index_global: 257
% index_in_page: 82
\[\begin{equation} \int \frac{1}{(x+a)(x+b)}dx = \frac{1}{b-a}\ln\frac{a+x}{b+x}, \text{ } a\ne b \end{equation}\]
% formula_record_end

% formula_record_start
% formula_id: formula-sheet-7-083
% page_slug: formula-sheet-7
% page_title: Formula Sheet
% page_url: https://formulasheet.com/#q|l|1022
% subject_slug: formulasheet
% subject_path: /formulasheet/
% section: Integral Table [integral-table.com]
% subsection: 
% formula_name: 
% source: formulasheet-latex
% index_global: 258
% index_in_page: 83
\[\begin{equation} \int \frac{1}{(x+a)^2}dx = -\frac{1}{x+a} \end{equation}\]
% formula_record_end

% formula_record_start
% formula_id: formula-sheet-7-084
% page_slug: formula-sheet-7
% page_title: Formula Sheet
% page_url: https://formulasheet.com/#q|l|1022
% subject_slug: formulasheet
% subject_path: /formulasheet/
% section: Integral Table [integral-table.com]
% subsection: 
% formula_name: 
% source: formulasheet-latex
% index_global: 259
% index_in_page: 84
\[\begin{equation} \int \frac{1}{1+x^2}dx = \tan^{-1}x \end{equation}\]
% formula_record_end

% formula_record_start
% formula_id: formula-sheet-7-085
% page_slug: formula-sheet-7
% page_title: Formula Sheet
% page_url: https://formulasheet.com/#q|l|1022
% subject_slug: formulasheet
% subject_path: /formulasheet/
% section: Integral Table [integral-table.com]
% subsection: 
% formula_name: 
% source: formulasheet-latex
% index_global: 260
% index_in_page: 85
\[\begin{equation} \int \frac{1}{ax+b}dx = \frac{1}{a} \ln |ax + b| \end{equation}\]
% formula_record_end

% formula_record_start
% formula_id: formula-sheet-7-086
% page_slug: formula-sheet-7
% page_title: Formula Sheet
% page_url: https://formulasheet.com/#q|l|1022
% subject_slug: formulasheet
% subject_path: /formulasheet/
% section: Integral Table [integral-table.com]
% subsection: 
% formula_name: 
% source: formulasheet-latex
% index_global: 261
% index_in_page: 86
\[\begin{equation} \int \frac{1}{ax^2+bx+c}dx = \frac{2}{\sqrt{4ac-b^2}}\tan^{-1}\frac{2ax+b}{\sqrt{4ac-b^2}} \end{equation}\]
% formula_record_end

% formula_record_start
% formula_id: formula-sheet-7-087
% page_slug: formula-sheet-7
% page_title: Formula Sheet
% page_url: https://formulasheet.com/#q|l|1022
% subject_slug: formulasheet
% subject_path: /formulasheet/
% section: Integral Table [integral-table.com]
% subsection: 
% formula_name: 
% source: formulasheet-latex
% index_global: 262
% index_in_page: 87
\[\begin{equation} \int \frac{1}{a^2+x^2}dx = \frac{1}{a}\tan^{-1}\frac{x}{a} \end{equation}\]
% formula_record_end

% formula_record_start
% formula_id: formula-sheet-7-088
% page_slug: formula-sheet-7
% page_title: Formula Sheet
% page_url: https://formulasheet.com/#q|l|1022
% subject_slug: formulasheet
% subject_path: /formulasheet/
% section: Integral Table [integral-table.com]
% subsection: 
% formula_name: 
% source: formulasheet-latex
% index_global: 263
% index_in_page: 88
\[\begin{equation} \int \frac{1}{x}dx = \ln |x| \end{equation}\]
% formula_record_end

% formula_record_start
% formula_id: formula-sheet-7-089
% page_slug: formula-sheet-7
% page_title: Formula Sheet
% page_url: https://formulasheet.com/#q|l|1022
% subject_slug: formulasheet
% subject_path: /formulasheet/
% section: Integral Table [integral-table.com]
% subsection: 
% formula_name: 
% source: formulasheet-latex
% index_global: 264
% index_in_page: 89
\[\begin{equation} \int \frac{1}{\sqrt{a-x}} dx = -2\sqrt{a-x} \end{equation}\]
% formula_record_end

% formula_record_start
% formula_id: formula-sheet-7-090
% page_slug: formula-sheet-7
% page_title: Formula Sheet
% page_url: https://formulasheet.com/#q|l|1022
% subject_slug: formulasheet
% subject_path: /formulasheet/
% section: Integral Table [integral-table.com]
% subsection: 
% formula_name: 
% source: formulasheet-latex
% index_global: 265
% index_in_page: 90
\[\begin{equation} \int \frac{1}{\sqrt{a^2 - x^2}} dx = \sin^{-1}\frac{x}{a} \end{equation}\]
% formula_record_end

% formula_record_start
% formula_id: formula-sheet-7-091
% page_slug: formula-sheet-7
% page_title: Formula Sheet
% page_url: https://formulasheet.com/#q|l|1022
% subject_slug: formulasheet
% subject_path: /formulasheet/
% section: Integral Table [integral-table.com]
% subsection: 
% formula_name: 
% source: formulasheet-latex
% index_global: 266
% index_in_page: 91
\[\begin{equation} \int \frac{1}{\sqrt{x\pm a}} dx = 2\sqrt{x\pm a} \end{equation}\]
% formula_record_end

% formula_record_start
% formula_id: formula-sheet-7-092
% page_slug: formula-sheet-7
% page_title: Formula Sheet
% page_url: https://formulasheet.com/#q|l|1022
% subject_slug: formulasheet
% subject_path: /formulasheet/
% section: Integral Table [integral-table.com]
% subsection: 
% formula_name: 
% source: formulasheet-latex
% index_global: 267
% index_in_page: 92
\[\begin{equation} \int \frac{1}{\sqrt{x^2 \pm a^2}} dx = \ln \left | x + \sqrt{x^2 \pm a^2} \right | \end{equation}\]
% formula_record_end

% formula_record_start
% formula_id: formula-sheet-7-093
% page_slug: formula-sheet-7
% page_title: Formula Sheet
% page_url: https://formulasheet.com/#q|l|1022
% subject_slug: formulasheet
% subject_path: /formulasheet/
% section: Integral Table [integral-table.com]
% subsection: 
% formula_name: 
% source: formulasheet-latex
% index_global: 268
% index_in_page: 93
\[\begin{equation} \int \frac{x^2}{a^2+x^2}dx = x-a\tan^{-1}\frac{x}{a} \end{equation}\]
% formula_record_end

% formula_record_start
% formula_id: formula-sheet-7-094
% page_slug: formula-sheet-7
% page_title: Formula Sheet
% page_url: https://formulasheet.com/#q|l|1022
% subject_slug: formulasheet
% subject_path: /formulasheet/
% section: Integral Table [integral-table.com]
% subsection: 
% formula_name: 
% source: formulasheet-latex
% index_global: 269
% index_in_page: 94
\[\begin{equation} \int \frac{x^3}{a^2+x^2}dx = \frac{1}{2}x^2-\frac{1}{2}a^2\ln|a^2+x^2| \end{equation}\]
% formula_record_end

% formula_record_start
% formula_id: formula-sheet-7-095
% page_slug: formula-sheet-7
% page_title: Formula Sheet
% page_url: https://formulasheet.com/#q|l|1022
% subject_slug: formulasheet
% subject_path: /formulasheet/
% section: Integral Table [integral-table.com]
% subsection: 
% formula_name: 
% source: formulasheet-latex
% index_global: 270
% index_in_page: 95
\[\begin{equation} \int \frac{x}{(x+a)^2}dx = \frac{a}{a+x}+\ln |a+x| \end{equation}\]
% formula_record_end

% formula_record_start
% formula_id: formula-sheet-7-096
% page_slug: formula-sheet-7
% page_title: Formula Sheet
% page_url: https://formulasheet.com/#q|l|1022
% subject_slug: formulasheet
% subject_path: /formulasheet/
% section: Integral Table [integral-table.com]
% subsection: 
% formula_name: 
% source: formulasheet-latex
% index_global: 271
% index_in_page: 96
\[\begin{equation} \int \frac{x}{a^2+x^2}dx = \frac{1}{2}\ln|a^2+x^2| \end{equation}\]
% formula_record_end

% formula_record_start
% formula_id: formula-sheet-7-097
% page_slug: formula-sheet-7
% page_title: Formula Sheet
% page_url: https://formulasheet.com/#q|l|1022
% subject_slug: formulasheet
% subject_path: /formulasheet/
% section: Integral Table [integral-table.com]
% subsection: 
% formula_name: 
% source: formulasheet-latex
% index_global: 272
% index_in_page: 97
\[\begin{equation} \int \frac{x}{\sqrt{a^2-x^2}}dx = -\sqrt{a^2-x^2} \end{equation}\]
% formula_record_end

% formula_record_start
% formula_id: formula-sheet-7-098
% page_slug: formula-sheet-7
% page_title: Formula Sheet
% page_url: https://formulasheet.com/#q|l|1022
% subject_slug: formulasheet
% subject_path: /formulasheet/
% section: Integral Table [integral-table.com]
% subsection: 
% formula_name: 
% source: formulasheet-latex
% index_global: 273
% index_in_page: 98
\[\begin{equation} \int \frac{x}{\sqrt{x\pm a} } dx = \frac{2}{3}(x\mp 2a)\sqrt{x\pm a} \end{equation}\]
% formula_record_end

% formula_record_start
% formula_id: formula-sheet-7-099
% page_slug: formula-sheet-7
% page_title: Formula Sheet
% page_url: https://formulasheet.com/#q|l|1022
% subject_slug: formulasheet
% subject_path: /formulasheet/
% section: Integral Table [integral-table.com]
% subsection: 
% formula_name: 
% source: formulasheet-latex
% index_global: 274
% index_in_page: 99
\[\begin{equation} \int \frac{x}{\sqrt{x^2\pm a^2}}dx = \sqrt{x^2 \pm a^2} \end{equation}\]
% formula_record_end

% formula_record_start
% formula_id: formula-sheet-7-100
% page_slug: formula-sheet-7
% page_title: Formula Sheet
% page_url: https://formulasheet.com/#q|l|1022
% subject_slug: formulasheet
% subject_path: /formulasheet/
% section: Integral Table [integral-table.com]
% subsection: 
% formula_name: 
% source: formulasheet-latex
% index_global: 275
% index_in_page: 100
\[\begin{equation} \int \frac{\ln ax}{x} dx = \frac{1}{2}\left ( \ln ax \right)^2 \end{equation}\]
% formula_record_end

% formula_record_start
% formula_id: formula-sheet-7-101
% page_slug: formula-sheet-7
% page_title: Formula Sheet
% page_url: https://formulasheet.com/#q|l|1022
% subject_slug: formulasheet
% subject_path: /formulasheet/
% section: Integral Table [integral-table.com]
% subsection: 
% formula_name: 
% source: formulasheet-latex
% index_global: 276
% index_in_page: 101
\[\begin{equation} \int \ln (ax + b) dx = \left ( x + \frac{b}{a} \right) \ln (ax+b) - x , a\ne 0 \end{equation}\]
% formula_record_end

% formula_record_start
% formula_id: formula-sheet-7-102
% page_slug: formula-sheet-7
% page_title: Formula Sheet
% page_url: https://formulasheet.com/#q|l|1022
% subject_slug: formulasheet
% subject_path: /formulasheet/
% section: Integral Table [integral-table.com]
% subsection: 
% formula_name: 
% source: formulasheet-latex
% index_global: 277
% index_in_page: 102
\[\begin{equation} \int \ln ax dx = x \ln ax - x \end{equation}\]
% formula_record_end

% formula_record_start
% formula_id: formula-sheet-7-103
% page_slug: formula-sheet-7
% page_title: Formula Sheet
% page_url: https://formulasheet.com/#q|l|1022
% subject_slug: formulasheet
% subject_path: /formulasheet/
% section: Integral Table [integral-table.com]
% subsection: 
% formula_name: 
% source: formulasheet-latex
% index_global: 278
% index_in_page: 103
\[\begin{equation} \int \sec x \csc x dx = \ln | \tan x | \end{equation}\]
% formula_record_end

% formula_record_start
% formula_id: formula-sheet-7-104
% page_slug: formula-sheet-7
% page_title: Formula Sheet
% page_url: https://formulasheet.com/#q|l|1022
% subject_slug: formulasheet
% subject_path: /formulasheet/
% section: Integral Table [integral-table.com]
% subsection: 
% formula_name: 
% source: formulasheet-latex
% index_global: 279
% index_in_page: 104
\[\begin{equation} \int \sec x \tan x dx = \sec x \end{equation}\]
% formula_record_end

% formula_record_start
% formula_id: formula-sheet-7-105
% page_slug: formula-sheet-7
% page_title: Formula Sheet
% page_url: https://formulasheet.com/#q|l|1022
% subject_slug: formulasheet
% subject_path: /formulasheet/
% section: Integral Table [integral-table.com]
% subsection: 
% formula_name: 
% source: formulasheet-latex
% index_global: 280
% index_in_page: 105
\[\begin{equation} \int \sec^2 ax dx = \frac{1}{a} \tan ax \end{equation}\]
% formula_record_end

% formula_record_start
% formula_id: formula-sheet-7-106
% page_slug: formula-sheet-7
% page_title: Formula Sheet
% page_url: https://formulasheet.com/#q|l|1022
% subject_slug: formulasheet
% subject_path: /formulasheet/
% section: Integral Table [integral-table.com]
% subsection: 
% formula_name: 
% source: formulasheet-latex
% index_global: 281
% index_in_page: 106
\[\begin{equation} \int \sec^2 x \tan x dx = \frac{1}{2} \sec^2 x \end{equation}\]
% formula_record_end

% formula_record_start
% formula_id: formula-sheet-7-107
% page_slug: formula-sheet-7
% page_title: Formula Sheet
% page_url: https://formulasheet.com/#q|l|1022
% subject_slug: formulasheet
% subject_path: /formulasheet/
% section: Integral Table [integral-table.com]
% subsection: 
% formula_name: 
% source: formulasheet-latex
% index_global: 282
% index_in_page: 107
\[\begin{equation} \int \sec^3 x \hspace{2pt}\text{dx} = \frac{1}{2} \sec x \tan x + \frac{1}{2}\ln | \sec x + \tan x | \end{equation}\]
% formula_record_end

% formula_record_start
% formula_id: formula-sheet-7-108
% page_slug: formula-sheet-7
% page_title: Formula Sheet
% page_url: https://formulasheet.com/#q|l|1022
% subject_slug: formulasheet
% subject_path: /formulasheet/
% section: Integral Table [integral-table.com]
% subsection: 
% formula_name: 
% source: formulasheet-latex
% index_global: 283
% index_in_page: 108
\[\begin{equation} \int \sec^n x \tan x dx = \frac{1}{n} \sec^n x , n\ne 0 \end{equation}\]
% formula_record_end

% formula_record_start
% formula_id: formula-sheet-7-109
% page_slug: formula-sheet-7
% page_title: Formula Sheet
% page_url: https://formulasheet.com/#q|l|1022
% subject_slug: formulasheet
% subject_path: /formulasheet/
% section: Integral Table [integral-table.com]
% subsection: 
% formula_name: 
% source: formulasheet-latex
% index_global: 284
% index_in_page: 109
\[\begin{equation} \int \sin ax dx = -\frac{1}{a} \cos ax \end{equation}\]
% formula_record_end

% formula_record_start
% formula_id: formula-sheet-7-110
% page_slug: formula-sheet-7
% page_title: Formula Sheet
% page_url: https://formulasheet.com/#q|l|1022
% subject_slug: formulasheet
% subject_path: /formulasheet/
% section: Integral Table [integral-table.com]
% subsection: 
% formula_name: 
% source: formulasheet-latex
% index_global: 285
% index_in_page: 110
\[\begin{equation} \int \sinh ax dx = \frac{1}{a} \cosh ax \end{equation}\]
% formula_record_end

% formula_record_start
% formula_id: formula-sheet-7-111
% page_slug: formula-sheet-7
% page_title: Formula Sheet
% page_url: https://formulasheet.com/#q|l|1022
% subject_slug: formulasheet
% subject_path: /formulasheet/
% section: Integral Table [integral-table.com]
% subsection: 
% formula_name: 
% source: formulasheet-latex
% index_global: 286
% index_in_page: 111
\[\begin{equation} \int \sinh ax \cosh ax dx= \frac{1}{4a}\left[ -2ax + \sinh 2ax \right] \end{equation}\]
% formula_record_end

% formula_record_start
% formula_id: formula-sheet-7-112
% page_slug: formula-sheet-7
% page_title: Formula Sheet
% page_url: https://formulasheet.com/#q|l|1022
% subject_slug: formulasheet
% subject_path: /formulasheet/
% section: Integral Table [integral-table.com]
% subsection: 
% formula_name: 
% source: formulasheet-latex
% index_global: 287
% index_in_page: 112
\[\begin{equation} \int \sin^2 ax dx = \frac{x}{2} - \frac{\sin 2ax} {4a} \end{equation}\]
% formula_record_end

% formula_record_start
% formula_id: formula-sheet-7-113
% page_slug: formula-sheet-7
% page_title: Formula Sheet
% page_url: https://formulasheet.com/#q|l|1022
% subject_slug: formulasheet
% subject_path: /formulasheet/
% section: Integral Table [integral-table.com]
% subsection: 
% formula_name: 
% source: formulasheet-latex
% index_global: 288
% index_in_page: 113
\[\begin{equation} \int \sin^2 ax \cos^2 ax dx = \frac{x}{8}-\frac{\sin 4ax}{32a} \end{equation}\]
% formula_record_end

% formula_record_start
% formula_id: formula-sheet-7-114
% page_slug: formula-sheet-7
% page_title: Formula Sheet
% page_url: https://formulasheet.com/#q|l|1022
% subject_slug: formulasheet
% subject_path: /formulasheet/
% section: Integral Table [integral-table.com]
% subsection: 
% formula_name: 
% source: formulasheet-latex
% index_global: 289
% index_in_page: 114
\[\begin{equation} \int \sin^2 x \cos x dx = \frac{1}{3} \sin^3 x \end{equation}\]
% formula_record_end

% formula_record_start
% formula_id: formula-sheet-7-115
% page_slug: formula-sheet-7
% page_title: Formula Sheet
% page_url: https://formulasheet.com/#q|l|1022
% subject_slug: formulasheet
% subject_path: /formulasheet/
% section: Integral Table [integral-table.com]
% subsection: 
% formula_name: 
% source: formulasheet-latex
% index_global: 290
% index_in_page: 115
\[\begin{equation} \int \sin^3 ax dx = -\frac{3 \cos ax}{4a} + \frac{\cos 3ax} {12a} \end{equation}\]
% formula_record_end

% formula_record_start
% formula_id: formula-sheet-7-116
% page_slug: formula-sheet-7
% page_title: Formula Sheet
% page_url: https://formulasheet.com/#q|l|1022
% subject_slug: formulasheet
% subject_path: /formulasheet/
% section: Integral Table [integral-table.com]
% subsection: 
% formula_name: 
% source: formulasheet-latex
% index_global: 291
% index_in_page: 116
\[\begin{equation} \int \sqrt{ax+b}dx = \left(\frac{2b}{3a}+\frac{2x}{3}\right)\sqrt{ax+b} \end{equation}\]
% formula_record_end

% formula_record_start
% formula_id: formula-sheet-7-117
% page_slug: formula-sheet-7
% page_title: Formula Sheet
% page_url: https://formulasheet.com/#q|l|1022
% subject_slug: formulasheet
% subject_path: /formulasheet/
% section: Integral Table [integral-table.com]
% subsection: 
% formula_name: 
% source: formulasheet-latex
% index_global: 292
% index_in_page: 117
\[\begin{equation} \int \sqrt{x-a} dx = \frac{2}{3}(x-a)^{3/2} \end{equation}\]
% formula_record_end

% formula_record_start
% formula_id: formula-sheet-7-118
% page_slug: formula-sheet-7
% page_title: Formula Sheet
% page_url: https://formulasheet.com/#q|l|1022
% subject_slug: formulasheet
% subject_path: /formulasheet/
% section: Integral Table [integral-table.com]
% subsection: 
% formula_name: 
% source: formulasheet-latex
% index_global: 293
% index_in_page: 118
\[\begin{equation} \int \tan ax dx = -\frac{1}{a} \ln \cos ax \end{equation}\]
% formula_record_end

% formula_record_start
% formula_id: formula-sheet-7-119
% page_slug: formula-sheet-7
% page_title: Formula Sheet
% page_url: https://formulasheet.com/#q|l|1022
% subject_slug: formulasheet
% subject_path: /formulasheet/
% section: Integral Table [integral-table.com]
% subsection: 
% formula_name: 
% source: formulasheet-latex
% index_global: 294
% index_in_page: 119
\[\begin{equation} \int \tan^2 ax dx = -x + \frac{1}{a} \tan ax \end{equation}\]
% formula_record_end

% formula_record_start
% formula_id: formula-sheet-7-120
% page_slug: formula-sheet-7
% page_title: Formula Sheet
% page_url: https://formulasheet.com/#q|l|1022
% subject_slug: formulasheet
% subject_path: /formulasheet/
% section: Integral Table [integral-table.com]
% subsection: 
% formula_name: 
% source: formulasheet-latex
% index_global: 295
% index_in_page: 120
\[\begin{equation} \int \tan^3 ax dx = \frac{1}{a} \ln \cos ax + \frac{1}{2a}\sec^2 ax \end{equation}\]
% formula_record_end

% formula_record_start
% formula_id: formula-sheet-7-121
% page_slug: formula-sheet-7
% page_title: Formula Sheet
% page_url: https://formulasheet.com/#q|l|1022
% subject_slug: formulasheet
% subject_path: /formulasheet/
% section: Integral Table [integral-table.com]
% subsection: 
% formula_name: 
% source: formulasheet-latex
% index_global: 296
% index_in_page: 121
\[\begin{equation} \int\frac{dx}{(a^2+x^2)^{3/2}}=\frac{x}{a^2\sqrt{a^2+x^2}} \end{equation}\]
% formula_record_end

% formula_record_start
% formula_id: formula-sheet-7-122
% page_slug: formula-sheet-7
% page_title: Formula Sheet
% page_url: https://formulasheet.com/#q|l|1022
% subject_slug: formulasheet
% subject_path: /formulasheet/
% section: Integral Table [integral-table.com]
% subsection: 
% formula_name: 
% source: formulasheet-latex
% index_global: 297
% index_in_page: 122
\[\begin{split} \int x^n e^{ax}\hspace{2pt}\text{d}x = \frac{(-1)^n}{a^{n+1}}\Gamma[1+n,-ax], \\\ \text{ where } \Gamma(a,x)=\int_x^{\infty} t^{a-1}e^{-t}\hspace{2pt}\text{d}t \end{split}\]
% formula_record_end

% formula_record_start
% formula_id: formula-sheet-8-001
% page_slug: formula-sheet-8
% page_title: Formula Sheet
% page_url: https://formulasheet.com/#q|l|1227
% subject_slug: formulasheet
% subject_path: /formulasheet/
% section: Mathematical constants
% subsection: Euler's number
% formula_name: 
% source: formulasheet-latex
% index_global: 298
% index_in_page: 1
\[\begin{equation} e \approx 2.718281828459045... \end{equation}\]
% formula_record_end

% formula_record_start
% formula_id: formula-sheet-8-002
% page_slug: formula-sheet-8
% page_title: Formula Sheet
% page_url: https://formulasheet.com/#q|l|1227
% subject_slug: formulasheet
% subject_path: /formulasheet/
% section: Mathematical constants
% subsection: Golden ratio
% formula_name: 
% source: formulasheet-latex
% index_global: 299
% index_in_page: 2
\[\begin{equation} \varphi = \frac{1+\sqrt{5}}{2} \approx 1.61803398874989... \end{equation}\]
% formula_record_end

% formula_record_start
% formula_id: formula-sheet-8-003
% page_slug: formula-sheet-8
% page_title: Formula Sheet
% page_url: https://formulasheet.com/#q|l|1227
% subject_slug: formulasheet
% subject_path: /formulasheet/
% section: Mathematical constants
% subsection: Imaginary unit
% formula_name: 
% source: formulasheet-latex
% index_global: 300
% index_in_page: 3
\[\begin{equation} i = \sqrt{-1} \end{equation}\]
% formula_record_end

% formula_record_start
% formula_id: formula-sheet-8-004
% page_slug: formula-sheet-8
% page_title: Formula Sheet
% page_url: https://formulasheet.com/#q|l|1227
% subject_slug: formulasheet
% subject_path: /formulasheet/
% section: Mathematical constants
% subsection: Pi (Archimedes' constant)
% formula_name: 
% source: formulasheet-latex
% index_global: 301
% index_in_page: 4
\[\begin{equation} \pi = \frac{C}{d} \approx 3.141592653589793... \end{equation}\]
% formula_record_end

% formula_record_start
% formula_id: formula-sheet-9-001
% page_slug: formula-sheet-9
% page_title: Formula Sheet
% page_url: https://formulasheet.com/#q|l|1245
% subject_slug: formulasheet
% subject_path: /formulasheet/
% section: Trigonometric identities
% subsection: Cosecant
% formula_name: 
% source: formulasheet-latex
% index_global: 302
% index_in_page: 1
\[\begin{equation} \csc \theta = \frac{1}{\sin \theta} \end{equation}\]
% formula_record_end

% formula_record_start
% formula_id: formula-sheet-9-002
% page_slug: formula-sheet-9
% page_title: Formula Sheet
% page_url: https://formulasheet.com/#q|l|1245
% subject_slug: formulasheet
% subject_path: /formulasheet/
% section: Trigonometric identities
% subsection: Cosine double-angle formula
% formula_name: 
% source: formulasheet-latex
% index_global: 303
% index_in_page: 2
\[\begin{align} \cos 2\theta &= \cos^2 \theta - \sin^2 \theta \\\ &= 2 \cos^2 \theta - 1 \\\ &= 1 - 2 \sin^2 \theta \\\ &= \frac{1 - \tan^2 \theta} {1 + \tan^2 \theta} \end{align}\]
% formula_record_end

% formula_record_start
% formula_id: formula-sheet-9-003
% page_slug: formula-sheet-9
% page_title: Formula Sheet
% page_url: https://formulasheet.com/#q|l|1245
% subject_slug: formulasheet
% subject_path: /formulasheet/
% section: Trigonometric identities
% subsection: Cosine sum and difference identity
% formula_name: 
% source: formulasheet-latex
% index_global: 304
% index_in_page: 3
\[\begin{equation} \cos(\alpha \pm \beta) = \cos \alpha \cos \beta \mp \sin \alpha \sin \beta \end{equation}\]
% formula_record_end

% formula_record_start
% formula_id: formula-sheet-9-004
% page_slug: formula-sheet-9
% page_title: Formula Sheet
% page_url: https://formulasheet.com/#q|l|1245
% subject_slug: formulasheet
% subject_path: /formulasheet/
% section: Trigonometric identities
% subsection: Cosine triple-angle formula
% formula_name: 
% source: formulasheet-latex
% index_global: 305
% index_in_page: 4
\[\begin{align}\cos 3\theta & = \cos^3\theta - 3 \sin^2 \theta\cos \theta \\\ & = 4 \cos^3\theta - 3 \cos\theta\end{align}\]
% formula_record_end

% formula_record_start
% formula_id: formula-sheet-9-005
% page_slug: formula-sheet-9
% page_title: Formula Sheet
% page_url: https://formulasheet.com/#q|l|1245
% subject_slug: formulasheet
% subject_path: /formulasheet/
% section: Trigonometric identities
% subsection: Cotangent
% formula_name: 
% source: formulasheet-latex
% index_global: 306
% index_in_page: 5
\[\begin{equation} \cot \theta = \frac{1}{\tan \theta} = \frac{\cos \theta}{\sin \theta} \end{equation}\]
% formula_record_end

% formula_record_start
% formula_id: formula-sheet-9-006
% page_slug: formula-sheet-9
% page_title: Formula Sheet
% page_url: https://formulasheet.com/#q|l|1245
% subject_slug: formulasheet
% subject_path: /formulasheet/
% section: Trigonometric identities
% subsection: Cotangent double-angle formula
% formula_name: 
% source: formulasheet-latex
% index_global: 307
% index_in_page: 6
\[\begin{equation} \cot 2\theta = \frac{\cot^2 \theta - 1}{2 \cot \theta} \end{equation}\]
% formula_record_end

% formula_record_start
% formula_id: formula-sheet-9-007
% page_slug: formula-sheet-9
% page_title: Formula Sheet
% page_url: https://formulasheet.com/#q|l|1245
% subject_slug: formulasheet
% subject_path: /formulasheet/
% section: Trigonometric identities
% subsection: Cotangent triple-angle formula
% formula_name: 
% source: formulasheet-latex
% index_global: 308
% index_in_page: 7
\[\begin{equation} \cot 3\theta = \frac{3 \cot\theta - \cot^3\theta}{1 - 3 \cot^2\theta} \end{equation}\]
% formula_record_end

% formula_record_start
% formula_id: formula-sheet-9-008
% page_slug: formula-sheet-9
% page_title: Formula Sheet
% page_url: https://formulasheet.com/#q|l|1245
% subject_slug: formulasheet
% subject_path: /formulasheet/
% section: Trigonometric identities
% subsection: Pythagorean identity
% formula_name: 
% source: formulasheet-latex
% index_global: 309
% index_in_page: 8
\[\begin{equation} \sin^2 \theta + \cos^2 \theta = 1 \end{equation}\]
% formula_record_end

% formula_record_start
% formula_id: formula-sheet-9-009
% page_slug: formula-sheet-9
% page_title: Formula Sheet
% page_url: https://formulasheet.com/#q|l|1245
% subject_slug: formulasheet
% subject_path: /formulasheet/
% section: Trigonometric identities
% subsection: Pythagorean identity (cotangent and cosecant)
% formula_name: 
% source: formulasheet-latex
% index_global: 310
% index_in_page: 9
\[\begin{equation} 1 + \cot^2 \theta = \csc^2 \theta \end{equation}\]
% formula_record_end

% formula_record_start
% formula_id: formula-sheet-9-010
% page_slug: formula-sheet-9
% page_title: Formula Sheet
% page_url: https://formulasheet.com/#q|l|1245
% subject_slug: formulasheet
% subject_path: /formulasheet/
% section: Trigonometric identities
% subsection: Pythagorean identity (tangent and secant)
% formula_name: 
% source: formulasheet-latex
% index_global: 311
% index_in_page: 10
\[\begin{equation} \tan^2 \theta + 1 = \sec^2 \theta \end{equation}\]
% formula_record_end

% formula_record_start
% formula_id: formula-sheet-9-011
% page_slug: formula-sheet-9
% page_title: Formula Sheet
% page_url: https://formulasheet.com/#q|l|1245
% subject_slug: formulasheet
% subject_path: /formulasheet/
% section: Trigonometric identities
% subsection: Secant
% formula_name: 
% source: formulasheet-latex
% index_global: 312
% index_in_page: 11
\[\begin{equation} \sec\theta= \frac{1}{\cos \theta} \end{equation}\]
% formula_record_end

% formula_record_start
% formula_id: formula-sheet-9-012
% page_slug: formula-sheet-9
% page_title: Formula Sheet
% page_url: https://formulasheet.com/#q|l|1245
% subject_slug: formulasheet
% subject_path: /formulasheet/
% section: Trigonometric identities
% subsection: Sine double-angle formula
% formula_name: 
% source: formulasheet-latex
% index_global: 313
% index_in_page: 12
\[\begin{align} \sin 2\theta &= 2 \sin \theta \cos \theta \\\ &= \frac{2 \tan \theta} {1 + \tan^2 \theta} \end{align}\]
% formula_record_end

% formula_record_start
% formula_id: formula-sheet-9-013
% page_slug: formula-sheet-9
% page_title: Formula Sheet
% page_url: https://formulasheet.com/#q|l|1245
% subject_slug: formulasheet
% subject_path: /formulasheet/
% section: Trigonometric identities
% subsection: Sine sum and difference identity
% formula_name: 
% source: formulasheet-latex
% index_global: 314
% index_in_page: 13
\[\begin{equation} \sin(\alpha \pm \beta) = \sin \alpha \cos \beta \pm \cos \alpha \sin \beta \end{equation}\]
% formula_record_end

% formula_record_start
% formula_id: formula-sheet-9-014
% page_slug: formula-sheet-9
% page_title: Formula Sheet
% page_url: https://formulasheet.com/#q|l|1245
% subject_slug: formulasheet
% subject_path: /formulasheet/
% section: Trigonometric identities
% subsection: Sine triple-angle formula
% formula_name: 
% source: formulasheet-latex
% index_global: 315
% index_in_page: 14
\[\begin{align} \sin 3\theta & = -\sin^3\theta + 3 \cos^2\theta \sin\theta\\\ & = -4\sin^3\theta + 3\sin\theta \end{align}\]
% formula_record_end

% formula_record_start
% formula_id: formula-sheet-9-015
% page_slug: formula-sheet-9
% page_title: Formula Sheet
% page_url: https://formulasheet.com/#q|l|1245
% subject_slug: formulasheet
% subject_path: /formulasheet/
% section: Trigonometric identities
% subsection: Tangent
% formula_name: 
% source: formulasheet-latex
% index_global: 316
% index_in_page: 15
\[\begin{equation} \tan \theta = \frac{\sin \theta}{\cos \theta} \end{equation}\]
% formula_record_end

% formula_record_start
% formula_id: formula-sheet-9-016
% page_slug: formula-sheet-9
% page_title: Formula Sheet
% page_url: https://formulasheet.com/#q|l|1245
% subject_slug: formulasheet
% subject_path: /formulasheet/
% section: Trigonometric identities
% subsection: Tangent double-angle formula
% formula_name: 
% source: formulasheet-latex
% index_global: 317
% index_in_page: 16
\[\begin{equation} \tan 2\theta = \frac{2 \tan \theta} {1 - \tan^2 \theta} \end{equation}\]
% formula_record_end

% formula_record_start
% formula_id: formula-sheet-9-017
% page_slug: formula-sheet-9
% page_title: Formula Sheet
% page_url: https://formulasheet.com/#q|l|1245
% subject_slug: formulasheet
% subject_path: /formulasheet/
% section: Trigonometric identities
% subsection: Tangent sum and difference identity
% formula_name: 
% source: formulasheet-latex
% index_global: 318
% index_in_page: 17
\[\begin{equation} \tan(\alpha \pm \beta) = \frac{\tan \alpha \pm \tan \beta}{1 \mp \tan \alpha \tan \beta} \end{equation}\]
% formula_record_end

% formula_record_start
% formula_id: formula-sheet-9-018
% page_slug: formula-sheet-9
% page_title: Formula Sheet
% page_url: https://formulasheet.com/#q|l|1245
% subject_slug: formulasheet
% subject_path: /formulasheet/
% section: Trigonometric identities
% subsection: Tangent triple-angle formula
% formula_name: 
% source: formulasheet-latex
% index_global: 319
% index_in_page: 18
\[\begin{equation} \tan 3\theta = \frac{3 \tan\theta - \tan^3\theta}{1 - 3 \tan^2\theta} \end{equation}\]
% formula_record_end

% formula_record_start
% formula_id: formula-sheet-9-019
% page_slug: formula-sheet-9
% page_title: Formula Sheet
% page_url: https://formulasheet.com/#q|l|1245
% subject_slug: formulasheet
% subject_path: /formulasheet/
% section: Trigonometric identities
% subsection: Trigonometric functions shifts and periodicity
% formula_name: 
% source: formulasheet-latex
% index_global: 320
% index_in_page: 19
\[\begin{align} \sin(\theta + \tfrac{\pi}{2}) &= +\cos \theta & \sin(\theta + \pi) &= -\sin \theta & \sin(\theta + 2\pi) &= +\sin \theta \\\ \cos(\theta + \tfrac{\pi}{2}) &= -\sin \theta & \cos(\theta + \pi) &= -\cos \theta & \cos(\theta + 2\pi) &= +\cos \theta \\\ \tan(\theta + \tfrac{\pi}{2}) &= -\cot \theta & \tan(\theta + \pi) &= +\tan \theta & \tan(\theta + 2\pi) &= +\tan \theta \\\ \csc(\theta + \tfrac{\pi}{2}) &= +\sec \theta & \csc(\theta + \pi) &= -\csc \theta & \csc(\theta + 2\pi) &= +\csc \theta \\\ \sec(\theta + \tfrac{\pi}{2}) &= -\csc \theta & \sec(\theta + \pi) &= -\sec \theta & \sec(\theta + 2\pi) &= +\sec \theta \\\ \cot(\theta + \tfrac{\pi}{2}) &= -\tan \theta & \cot(\theta + \pi) &= +\cot \theta & \cot(\theta + 2\pi) &= +\cot \theta \end{align}\]
% formula_record_end

% formula_record_start
% formula_id: formula-sheet-9-020
% page_slug: formula-sheet-9
% page_title: Formula Sheet
% page_url: https://formulasheet.com/#q|l|1245
% subject_slug: formulasheet
% subject_path: /formulasheet/
% section: Trigonometric identities
% subsection: Trigonometric functions symmetry
% formula_name: 
% source: formulasheet-latex
% index_global: 321
% index_in_page: 20
\[\begin{align} \sin(-\theta) &= -\sin \theta & \sin(\tfrac{\pi}{2} - \theta) &= +\cos \theta & \sin(\pi - \theta) &= +\sin \theta \\\ \cos(-\theta) &= +\cos \theta & \cos(\tfrac{\pi}{2} - \theta) &= +\sin \theta & \cos(\pi - \theta) &= -\cos \theta \\\ \tan(-\theta) &= -\tan \theta & \tan(\tfrac{\pi}{2} - \theta) &= +\cot \theta & \tan(\pi - \theta) &= -\tan \theta \\\ \csc(-\theta) &= -\csc \theta & \csc(\tfrac{\pi}{2} - \theta) &= +\sec \theta & \csc(\pi - \theta) &= +\csc \theta \\\ \sec(-\theta) &= +\sec \theta & \sec(\tfrac{\pi}{2} - \theta) &= +\csc \theta & \sec(\pi - \theta) &= -\sec \theta \\\ \cot(-\theta) &= -\cot \theta & \cot(\tfrac{\pi}{2} - \theta) &= +\tan \theta & \cot(\pi - \theta) &= -\cot \theta \end{align}\]
% formula_record_end

% formula_record_start
% formula_id: formula-sheet-10-001
% page_slug: formula-sheet-10
% page_title: Formula Sheet
% page_url: https://formulasheet.com/#q|l|1250
% subject_slug: formulasheet
% subject_path: /formulasheet/
% section: Trigonometric function derivatives
% subsection: Derivative of arccosecant
% formula_name: 
% source: formulasheet-latex
% index_global: 322
% index_in_page: 1
\[\begin{equation} \frac{\mathrm{d}}{\mathrm{d} x} \mathrm{arccsc}\, x = \frac{-1}{\left\lvert x \right\rvert \sqrt{x^2 - 1}} \end{equation}\]
% formula_record_end

% formula_record_start
% formula_id: formula-sheet-10-002
% page_slug: formula-sheet-10
% page_title: Formula Sheet
% page_url: https://formulasheet.com/#q|l|1250
% subject_slug: formulasheet
% subject_path: /formulasheet/
% section: Trigonometric function derivatives
% subsection: Derivative of arccosine
% formula_name: 
% source: formulasheet-latex
% index_global: 323
% index_in_page: 2
\[\begin{equation} \frac{\mathrm{d}}{\mathrm{d} x} \arccos x = \frac{-1}{\sqrt{1-x^2}} \end{equation}\]
% formula_record_end

% formula_record_start
% formula_id: formula-sheet-10-003
% page_slug: formula-sheet-10
% page_title: Formula Sheet
% page_url: https://formulasheet.com/#q|l|1250
% subject_slug: formulasheet
% subject_path: /formulasheet/
% section: Trigonometric function derivatives
% subsection: Derivative of arccotangent
% formula_name: 
% source: formulasheet-latex
% index_global: 324
% index_in_page: 3
\[\begin{equation} \frac{\mathrm{d}}{\mathrm{d} x} \mathrm{arccot}\, x = \frac{-1}{1+x^2} \end{equation}\]
% formula_record_end

% formula_record_start
% formula_id: formula-sheet-10-004
% page_slug: formula-sheet-10
% page_title: Formula Sheet
% page_url: https://formulasheet.com/#q|l|1250
% subject_slug: formulasheet
% subject_path: /formulasheet/
% section: Trigonometric function derivatives
% subsection: Derivative of arcsecant
% formula_name: 
% source: formulasheet-latex
% index_global: 325
% index_in_page: 4
\[\begin{equation} \frac{\mathrm{d}}{\mathrm{d} x} \mathrm{arcsec}\, x = \frac{1}{\left\lvert x \right\rvert \sqrt{x^2 - 1}} \end{equation}\]
% formula_record_end

% formula_record_start
% formula_id: formula-sheet-10-005
% page_slug: formula-sheet-10
% page_title: Formula Sheet
% page_url: https://formulasheet.com/#q|l|1250
% subject_slug: formulasheet
% subject_path: /formulasheet/
% section: Trigonometric function derivatives
% subsection: Derivative of arcsine
% formula_name: 
% source: formulasheet-latex
% index_global: 326
% index_in_page: 5
\[\begin{equation} \frac{\mathrm{d}}{\mathrm{d} x} \arcsin x = \frac{1}{\sqrt{1-x^2}} \end{equation}\]
% formula_record_end

% formula_record_start
% formula_id: formula-sheet-10-006
% page_slug: formula-sheet-10
% page_title: Formula Sheet
% page_url: https://formulasheet.com/#q|l|1250
% subject_slug: formulasheet
% subject_path: /formulasheet/
% section: Trigonometric function derivatives
% subsection: Derivative of arctangent
% formula_name: 
% source: formulasheet-latex
% index_global: 327
% index_in_page: 6
\[\begin{equation} \frac{\mathrm{d}}{\mathrm{d} x} \arctan x = \frac{1}{1+x^2} \end{equation}\]
% formula_record_end

% formula_record_start
% formula_id: formula-sheet-10-007
% page_slug: formula-sheet-10
% page_title: Formula Sheet
% page_url: https://formulasheet.com/#q|l|1250
% subject_slug: formulasheet
% subject_path: /formulasheet/
% section: Trigonometric function derivatives
% subsection: Derivative of cosecant
% formula_name: 
% source: formulasheet-latex
% index_global: 328
% index_in_page: 7
\[\begin{equation} \frac{\mathrm{d}}{\mathrm{d} x} \csc x = -\csc x \cot x \end{equation}\]
% formula_record_end

% formula_record_start
% formula_id: formula-sheet-10-008
% page_slug: formula-sheet-10
% page_title: Formula Sheet
% page_url: https://formulasheet.com/#q|l|1250
% subject_slug: formulasheet
% subject_path: /formulasheet/
% section: Trigonometric function derivatives
% subsection: Derivative of cosine
% formula_name: 
% source: formulasheet-latex
% index_global: 329
% index_in_page: 8
\[\begin{equation} \frac{\mathrm{d}}{\mathrm{d} x} \cos x = - \sin x \end{equation}\]
% formula_record_end

% formula_record_start
% formula_id: formula-sheet-10-009
% page_slug: formula-sheet-10
% page_title: Formula Sheet
% page_url: https://formulasheet.com/#q|l|1250
% subject_slug: formulasheet
% subject_path: /formulasheet/
% section: Trigonometric function derivatives
% subsection: Derivative of cotangent
% formula_name: 
% source: formulasheet-latex
% index_global: 330
% index_in_page: 9
\[\begin{equation} \frac{\mathrm{d}}{\mathrm{d} x} \cot x = - \csc^2 x \end{equation}\]
% formula_record_end

% formula_record_start
% formula_id: formula-sheet-10-010
% page_slug: formula-sheet-10
% page_title: Formula Sheet
% page_url: https://formulasheet.com/#q|l|1250
% subject_slug: formulasheet
% subject_path: /formulasheet/
% section: Trigonometric function derivatives
% subsection: Derivative of secant
% formula_name: 
% source: formulasheet-latex
% index_global: 331
% index_in_page: 10
\[\begin{equation} \frac{\mathrm{d}}{\mathrm{d} x} \sec x = \sec x \tan x \end{equation}\]
% formula_record_end

% formula_record_start
% formula_id: formula-sheet-10-011
% page_slug: formula-sheet-10
% page_title: Formula Sheet
% page_url: https://formulasheet.com/#q|l|1250
% subject_slug: formulasheet
% subject_path: /formulasheet/
% section: Trigonometric function derivatives
% subsection: Derivative of sine
% formula_name: 
% source: formulasheet-latex
% index_global: 332
% index_in_page: 11
\[\begin{equation} \frac{\mathrm{d}}{\mathrm{d} x} \sin x = \cos x \end{equation}\]
% formula_record_end

% formula_record_start
% formula_id: formula-sheet-10-012
% page_slug: formula-sheet-10
% page_title: Formula Sheet
% page_url: https://formulasheet.com/#q|l|1250
% subject_slug: formulasheet
% subject_path: /formulasheet/
% section: Trigonometric function derivatives
% subsection: Derivative of tangent
% formula_name: 
% source: formulasheet-latex
% index_global: 333
% index_in_page: 12
\[\begin{equation} \frac{\mathrm{d}}{\mathrm{d} x} \tan x = \sec ^2 x \end{equation}\]
% formula_record_end

% formula_record_start
% formula_id: formula-sheet-11-001
% page_slug: formula-sheet-11
% page_title: Formula Sheet
% page_url: https://formulasheet.com/#q|l|1722
% subject_slug: formulasheet
% subject_path: /formulasheet/
% section: Sets
% subsection: Complement of a Set
% formula_name: 
% source: formulasheet-latex
% index_global: 334
% index_in_page: 1
\[\begin{equation} S^c = \{ x \in \Omega \, | \, x \notin S \} \end{equation}\]
% formula_record_end

% formula_record_start
% formula_id: formula-sheet-11-002
% page_slug: formula-sheet-11
% page_title: Formula Sheet
% page_url: https://formulasheet.com/#q|l|1722
% subject_slug: formulasheet
% subject_path: /formulasheet/
% section: Sets
% subsection: Countably Infinite Set
% formula_name: 
% source: formulasheet-latex
% index_global: 335
% index_in_page: 2
\[\begin{equation} S = \{x_1, x_2, \dotso \} \end{equation}\]
% formula_record_end

% formula_record_start
% formula_id: formula-sheet-11-003
% page_slug: formula-sheet-11
% page_title: Formula Sheet
% page_url: https://formulasheet.com/#q|l|1722
% subject_slug: formulasheet
% subject_path: /formulasheet/
% section: Sets
% subsection: De Morgan's Laws
% formula_name: 
% source: formulasheet-latex
% index_global: 336
% index_in_page: 3
\[\begin{align} \left( \bigcup_n S_n \right)^c &= \bigcap_n S_n^c \\\ \left( \bigcap_n S_n \right)^c &= \bigcup_n S_n^c \end{align}\]
% formula_record_end

% formula_record_start
% formula_id: formula-sheet-11-004
% page_slug: formula-sheet-11
% page_title: Formula Sheet
% page_url: https://formulasheet.com/#q|l|1722
% subject_slug: formulasheet
% subject_path: /formulasheet/
% section: Sets
% subsection: Equal Sets
% formula_name: 
% source: formulasheet-latex
% index_global: 337
% index_in_page: 4
\[\begin{equation} S = T \end{equation}\]
% formula_record_end

% formula_record_start
% formula_id: formula-sheet-11-005
% page_slug: formula-sheet-11
% page_title: Formula Sheet
% page_url: https://formulasheet.com/#q|l|1722
% subject_slug: formulasheet
% subject_path: /formulasheet/
% section: Sets
% subsection: Finite Set
% formula_name: 
% source: formulasheet-latex
% index_global: 338
% index_in_page: 5
\[\begin{equation} S = \{x_1, x_2, \dotso x_n \} \end{equation}\]
% formula_record_end

% formula_record_start
% formula_id: formula-sheet-11-006
% page_slug: formula-sheet-11
% page_title: Formula Sheet
% page_url: https://formulasheet.com/#q|l|1722
% subject_slug: formulasheet
% subject_path: /formulasheet/
% section: Sets
% subsection: Intersection of Sets
% formula_name: 
% source: formulasheet-latex
% index_global: 339
% index_in_page: 6
\[\begin{equation} \bigcap_{i=1}^N = S_1 \, \cup \, S_2 \, \cup \dotso = \{ x \, | \, x \in S_i \text{ for all } i\} \end{equation}\]
% formula_record_end

% formula_record_start
% formula_id: formula-sheet-11-007
% page_slug: formula-sheet-11
% page_title: Formula Sheet
% page_url: https://formulasheet.com/#q|l|1722
% subject_slug: formulasheet
% subject_path: /formulasheet/
% section: Sets
% subsection: Subset and Superset
% formula_name: 
% source: formulasheet-latex
% index_global: 340
% index_in_page: 7
\[\begin{equation} S \subset T \; \text{or} \; T \supset S \end{equation}\]
% formula_record_end

% formula_record_start
% formula_id: formula-sheet-11-008
% page_slug: formula-sheet-11
% page_title: Formula Sheet
% page_url: https://formulasheet.com/#q|l|1722
% subject_slug: formulasheet
% subject_path: /formulasheet/
% section: Sets
% subsection: Union of Sets
% formula_name: 
% source: formulasheet-latex
% index_global: 341
% index_in_page: 8
\[\begin{equation} \bigcup_{i=1}^N = S_1 \, \cup \, S_2 \, \cup \dotso = \{ x \, | \, x \in S_i \text{ for some } i\} \end{equation}\]
% formula_record_end

% formula_record_start
% formula_id: formula-sheet-11-009
% page_slug: formula-sheet-11
% page_title: Formula Sheet
% page_url: https://formulasheet.com/#q|l|1722
% subject_slug: formulasheet
% subject_path: /formulasheet/
% section: Sets
% subsection: Universal Set
% formula_name: 
% source: formulasheet-latex
% index_global: 342
% index_in_page: 9
\[\begin{equation} \Omega \end{equation}\]
% formula_record_end

% formula_record_start
% formula_id: formula-sheet-12-001
% page_slug: formula-sheet-12
% page_title: Formula Sheet
% page_url: https://formulasheet.com/#q|l|1271
% subject_slug: formulasheet
% subject_path: /formulasheet/
% section: Taylor series
% subsection: Arctangent Taylor series
% formula_name: 
% source: formulasheet-latex
% index_global: 343
% index_in_page: 1
\[\begin{align} \arctan x &= x - \frac{x^3}{3} + \frac{x^5}{5} - \frac{x^7}{7} + \frac{x^9}{9} - \dotso \\\ &= \sum_{n=0}^\infty (-1)^n \frac{x^{2n + 1}}{2n + 1} \end{align}\]
% formula_record_end

% formula_record_start
% formula_id: formula-sheet-12-002
% page_slug: formula-sheet-12
% page_title: Formula Sheet
% page_url: https://formulasheet.com/#q|l|1271
% subject_slug: formulasheet
% subject_path: /formulasheet/
% section: Taylor series
% subsection: Cosine Taylor series
% formula_name: 
% source: formulasheet-latex
% index_global: 344
% index_in_page: 2
\[\begin{align} \cos x &= 1 - \frac{x^2}{2!} + \frac{x^4}{4!} - \frac{x^6}{6!} + \frac{x^8}{8!} - \dotso \\\ &= \sum_{n=0}^\infty (-1)^n \frac{x^{2n}}{(2n)!} \end{align}\]
% formula_record_end

% formula_record_start
% formula_id: formula-sheet-12-003
% page_slug: formula-sheet-12
% page_title: Formula Sheet
% page_url: https://formulasheet.com/#q|l|1271
% subject_slug: formulasheet
% subject_path: /formulasheet/
% section: Taylor series
% subsection: Exponential Taylor series
% formula_name: 
% source: formulasheet-latex
% index_global: 345
% index_in_page: 3
\[\begin{align} e^x &= 1 + x + \frac{x^2}{2!} + \frac{x^3}{3!} + \frac{x^4}{4!} + \dotso \\\ &= \sum_{n=0}^\infty \frac{x^n}{n!} \end{align}\]
% formula_record_end

% formula_record_start
% formula_id: formula-sheet-12-004
% page_slug: formula-sheet-12
% page_title: Formula Sheet
% page_url: https://formulasheet.com/#q|l|1271
% subject_slug: formulasheet
% subject_path: /formulasheet/
% section: Taylor series
% subsection: Geometric Taylor series
% formula_name: 
% source: formulasheet-latex
% index_global: 346
% index_in_page: 4
\[\begin{align} \frac{1}{1-x} &= 1 + x + x^2 + x^3 + x^4 + \dotso \\\ &= \sum_{n=0}^\infty x^n \end{align}\]
% formula_record_end

% formula_record_start
% formula_id: formula-sheet-12-005
% page_slug: formula-sheet-12
% page_title: Formula Sheet
% page_url: https://formulasheet.com/#q|l|1271
% subject_slug: formulasheet
% subject_path: /formulasheet/
% section: Taylor series
% subsection: Natural logarithm Taylor series
% formula_name: 
% source: formulasheet-latex
% index_global: 347
% index_in_page: 5
\[\begin{align} \ln (1 + x) &= x - \frac{x^2}{2} + \frac{x^3}{3} - \frac{x^4}{4} + \frac{x^5}{5} - \dotso \\\ &= \sum_{n=1}^\infty (-1)^{n + 1} \frac{x^n}{n} \end{align}\]
% formula_record_end

% formula_record_start
% formula_id: formula-sheet-12-006
% page_slug: formula-sheet-12
% page_title: Formula Sheet
% page_url: https://formulasheet.com/#q|l|1271
% subject_slug: formulasheet
% subject_path: /formulasheet/
% section: Taylor series
% subsection: Sine Taylor series
% formula_name: 
% source: formulasheet-latex
% index_global: 348
% index_in_page: 6
\[\begin{align} \sin x &= x - \frac{x^3}{3!} + \frac{x^5}{5!} - \frac{x^7}{7!} + \frac{x^9}{9!} - \dotso \\\ &= \sum_{n=0}^\infty (-1)^n \frac{x^{2n + 1}}{(2n + 1)!} \end{align}\]
% formula_record_end

% formula_record_start
% formula_id: formula-sheet-13-001
% page_slug: formula-sheet-13
% page_title: Formula Sheet
% page_url: https://formulasheet.com/#q|l|1252
% subject_slug: formulasheet
% subject_path: /formulasheet/
% section: Trigonometric laws and theorems
% subsection: Law of cosines
% formula_name: 
% source: formulasheet-latex
% index_global: 349
% index_in_page: 1
\[\begin{equation} c^2 = a^2 + b^2 - 2 a b\cos\gamma \end{equation}\]
% formula_record_end

% formula_record_start
% formula_id: formula-sheet-13-002
% page_slug: formula-sheet-13
% page_title: Formula Sheet
% page_url: https://formulasheet.com/#q|l|1252
% subject_slug: formulasheet
% subject_path: /formulasheet/
% section: Trigonometric laws and theorems
% subsection: Law of cotangents
% formula_name: 
% source: formulasheet-latex
% index_global: 350
% index_in_page: 2
\[\begin{equation} \frac{\cot(\alpha / 2)}{s - a} = \frac{\cot(\beta / 2)}{s - b} = \frac{\cot(\gamma / 2)}{s - c} \end{equation}\]
% formula_record_end

% formula_record_start
% formula_id: formula-sheet-13-003
% page_slug: formula-sheet-13
% page_title: Formula Sheet
% page_url: https://formulasheet.com/#q|l|1252
% subject_slug: formulasheet
% subject_path: /formulasheet/
% section: Trigonometric laws and theorems
% subsection: Law of sines
% formula_name: 
% source: formulasheet-latex
% index_global: 351
% index_in_page: 3
\[\begin{equation} \frac{a}{\sin \alpha} = \frac{b}{\sin \beta} = \frac{c}{\sin \gamma} = D \end{equation}\]
% formula_record_end

% formula_record_start
% formula_id: formula-sheet-13-004
% page_slug: formula-sheet-13
% page_title: Formula Sheet
% page_url: https://formulasheet.com/#q|l|1252
% subject_slug: formulasheet
% subject_path: /formulasheet/
% section: Trigonometric laws and theorems
% subsection: Law of tangents
% formula_name: 
% source: formulasheet-latex
% index_global: 352
% index_in_page: 4
\[\begin{equation} \frac{a - b}{a + b} = \frac{\tan\left[ \frac{1}{2} (\alpha - \beta) \right]}{\tan \left[ \frac{1}{2} (\alpha + \beta) \right]} \end{equation}\]
% formula_record_end

% formula_record_start
% formula_id: formula-sheet-13-005
% page_slug: formula-sheet-13
% page_title: Formula Sheet
% page_url: https://formulasheet.com/#q|l|1252
% subject_slug: formulasheet
% subject_path: /formulasheet/
% section: Trigonometric laws and theorems
% subsection: Pythagorean theorem
% formula_name: 
% source: formulasheet-latex
% index_global: 353
% index_in_page: 5
\[\begin{equation} a^2 + b^2 = c^2\! \end{equation}\]
% formula_record_end

% formula_record_start
% formula_id: formula-sheet-13-006
% page_slug: formula-sheet-13
% page_title: Formula Sheet
% page_url: https://formulasheet.com/#q|l|1252
% subject_slug: formulasheet
% subject_path: /formulasheet/
% section: Trigonometric laws and theorems
% subsection: Pythagorean theorem hypotenuse solution
% formula_name: 
% source: formulasheet-latex
% index_global: 354
% index_in_page: 6
\[\begin{equation} c = \sqrt{a^2 + b^2} \end{equation}\]
% formula_record_end

% formula_record_start
% formula_id: formula-sheet-14-001
% page_slug: formula-sheet-14
% page_title: Formula Sheet
% page_url: https://formulasheet.com/#q|l|4988
% subject_slug: formulasheet
% subject_path: /formulasheet/
% section: Chapter 9
% subsection: 
% formula_name: 
% source: formulasheet-latex
% index_global: 355
% index_in_page: 1
\[\begin{align*} {}_mF_n \left[ \left\lgroup \qquad\qquad\qquad \right\rgroup x \right] \end{align*}\]
% formula_record_end

% formula_record_start
% formula_id: formula-sheet-14-002
% page_slug: formula-sheet-14
% page_title: Formula Sheet
% page_url: https://formulasheet.com/#q|l|4988
% subject_slug: formulasheet
% subject_path: /formulasheet/
% section: Chapter 9
% subsection: 
% formula_name: 
% source: formulasheet-latex
% index_global: 356
% index_in_page: 2
\[\begin{align*} \label{hypergeorows} & \alpha_1, \alpha_2,...\, \alpha_m; \\\ & \beta_1,\beta_2,...\,\; \beta_n\, ; \end{align*}\]
% formula_record_end

% formula_record_start
% formula_id: formula-sheet-14-003
% page_slug: formula-sheet-14
% page_title: Formula Sheet
% page_url: https://formulasheet.com/#q|l|4988
% subject_slug: formulasheet
% subject_path: /formulasheet/
% section: Chapter 9
% subsection: 
% formula_name: 
% source: formulasheet-latex
% index_global: 357
% index_in_page: 3
\[\begin{align} & (i)\; P_n(x) = {}_2F_1 \left(-n , n+1 ; 1 ; \frac{1-x}2 \right) \\\ & (ii)\; P_n^m(x) = \frac{(n+m)!}{(n-m)!}\frac{(1-x^2)^{m/2}}{2^mm!}{}_2F_1 \left( m-n , m+n+1 ; m+1 ; \frac{1-x}2 \right) \\\ & (iii)\; J_n(x) = \frac{e^{-ix}}{n!} \left( \frac x2 \right)^n {}_1F_1(n+\tfrac 12 ; 2n+1 ; 2ix) \\\ & (iv)\; H_{2n}(x) = (-1)^n\frac{(2n)!}{n!} {}_1F_1(-n ; \tfrac12 ; x^2) \\\ & (v)\; H_{2n+1}(x) = (-1)^n\frac{2(2n+1)!}{n!}x\; {}_1F_1(-n ; \tfrac32 ; x^2) \\\ & (vi)\; L_n(x) = {}_1F_1(-n ; 1 ; x) \\\ & (vii)\; L_n^k(x) = \frac{\Gamma(n+k+1)}{n!\Gamma(k+1)}{}_1F_1(-n ; k+1 ; x) \\\ & (viii)\; T_n(x) = {}_2F_1 \left(-n , n ; \frac12 ; \frac{1-x}2 \right) \\\ & (ix)\; U_n(x) = \sqrt{(1-x^2)}\; n\; {}_2F_1 \left(-n+1 , n+1 ; \frac32 ; \frac{1-x}2\right) \\\ & (x)\; C_n^\lambda(x) = \frac{\Gamma(n+2\lambda)}{n!\Gamma(2\lambda)}{}_2F_1 \left( -n , n+2\lambda ; \lambda +\tfrac12 ; \frac{1-x}2 \right) \\\ & (xi)\; P_n^{(\alpha,\beta)}(x) = \frac{\Gamma(n+\alpha+1)}{n!\Gamma(\alpha+1)} {}_2F_1\left(-n , n+\alpha+\beta+1 ; \alpha+1 ; \frac{1-x}2\right) \end{align}\]
% formula_record_end

% formula_record_start
% formula_id: formula-sheet-14-004
% page_slug: formula-sheet-14
% page_title: Formula Sheet
% page_url: https://formulasheet.com/#q|l|4988
% subject_slug: formulasheet
% subject_path: /formulasheet/
% section: Chapter 9
% subsection: 
% formula_name: 
% source: formulasheet-latex
% index_global: 358
% index_in_page: 4
\[\begin{align} \label{confluenthypergeo} {}_1F_1(\alpha ; \beta ; x) \end{align}\]
% formula_record_end

% formula_record_start
% formula_id: formula-sheet-14-005
% page_slug: formula-sheet-14
% page_title: Formula Sheet
% page_url: https://formulasheet.com/#q|l|4988
% subject_slug: formulasheet
% subject_path: /formulasheet/
% section: Chapter 9
% subsection: 
% formula_name: 
% source: formulasheet-latex
% index_global: 359
% index_in_page: 5
\[\begin{align} \label{diffeqsolution} {}_2F_1(\alpha , \beta ; \gamma ; x) \end{align}\]
% formula_record_end

% formula_record_start
% formula_id: formula-sheet-14-006
% page_slug: formula-sheet-14
% page_title: Formula Sheet
% page_url: https://formulasheet.com/#q|l|4988
% subject_slug: formulasheet
% subject_path: /formulasheet/
% section: Chapter 9
% subsection: 
% formula_name: 
% source: formulasheet-latex
% index_global: 360
% index_in_page: 6
\[\begin{align} \label{thrm9.2} {}_2F_1(\alpha , \beta ; \gamma ; x) = {}_2F_1(\beta , \alpha ; \gamma ; x) \end{align}\]
% formula_record_end

% formula_record_start
% formula_id: formula-sheet-14-007
% page_slug: formula-sheet-14
% page_title: Formula Sheet
% page_url: https://formulasheet.com/#q|l|4988
% subject_slug: formulasheet
% subject_path: /formulasheet/
% section: Chapter 9
% subsection: 
% formula_name: 
% source: formulasheet-latex
% index_global: 361
% index_in_page: 7
\[\begin{align} \label{thrm9.3} x(1-x)\frac{\mathrm d^2 y}{\mathrm d x^2} + [\gamma - (\alpha + \beta + 1)x]\frac{\mathrm d y}{\mathrm d x} - \alpha\beta y = 0 \end{align}\]
% formula_record_end

% formula_record_start
% formula_id: formula-sheet-14-008
% page_slug: formula-sheet-14
% page_title: Formula Sheet
% page_url: https://formulasheet.com/#q|l|4988
% subject_slug: formulasheet
% subject_path: /formulasheet/
% section: Chapter 9
% subsection: 
% formula_name: 
% source: formulasheet-latex
% index_global: 362
% index_in_page: 8
\[\begin{align} \label{thrm9.4} {}_2F_1(\alpha , \beta ; \gamma ; x) = \frac{\Gamma(\gamma)}{\Gamma(\beta)\Gamma(\gamma-\beta)}\int_0^1t^{\beta-1}(1-t)^{\gamma-\beta-1}(1-xt)^{-\alpha}\; \mathrm dt,\quad \gamma>\beta>0 \end{align}\]
% formula_record_end

% formula_record_start
% formula_id: formula-sheet-14-009
% page_slug: formula-sheet-14
% page_title: Formula Sheet
% page_url: https://formulasheet.com/#q|l|4988
% subject_slug: formulasheet
% subject_path: /formulasheet/
% section: Chapter 9
% subsection: 
% formula_name: 
% source: formulasheet-latex
% index_global: 363
% index_in_page: 9
\[\begin{align} \label{thrm9.6} x^2\frac{\mathrm d^2 y}{\mathrm d x^2} + (\beta-x)\frac{\mathrm d y}{\mathrm d x} - \alpha y = 0 \end{align}\]
% formula_record_end

% formula_record_start
% formula_id: formula-sheet-14-010
% page_slug: formula-sheet-14
% page_title: Formula Sheet
% page_url: https://formulasheet.com/#q|l|4988
% subject_slug: formulasheet
% subject_path: /formulasheet/
% section: Chapter 9
% subsection: 
% formula_name: 
% source: formulasheet-latex
% index_global: 364
% index_in_page: 10
\[\begin{align} \label{thrm9.7} {}_1F_1(\alpha ; \beta ; x) = \frac{\Gamma(\beta)}{\Gamma(\alpha)\Gamma(\beta-\alpha)}\int_0^1(1-t)^{\beta-\alpha-1}t^{\alpha-1}e^{xt}\; \mathrm dt,\quad \beta>\alpha>0 \end{align}\]
% formula_record_end

% formula_record_start
% formula_id: formula-sheet-14-011
% page_slug: formula-sheet-14
% page_title: Formula Sheet
% page_url: https://formulasheet.com/#q|l|4988
% subject_slug: formulasheet
% subject_path: /formulasheet/
% section: Chapter 9
% subsection: 
% formula_name: 
% source: formulasheet-latex
% index_global: 365
% index_in_page: 11
\[\begin{equation} \label{confhypergeo} _1F_1(\alpha;\beta;x) = \sum_{r=0}^\infty\frac{(\alpha)_r\,x^r}{(\beta)_rr!} \end{equation}\]
% formula_record_end

% formula_record_start
% formula_id: formula-sheet-14-012
% page_slug: formula-sheet-14
% page_title: Formula Sheet
% page_url: https://formulasheet.com/#q|l|4988
% subject_slug: formulasheet
% subject_path: /formulasheet/
% section: Chapter 9
% subsection: 
% formula_name: 
% source: formulasheet-latex
% index_global: 366
% index_in_page: 12
\[\begin{equation} \label{hypergeodef} _{m}F_{n}(\alpha_1,\alpha_2,...\alpha_m;\, \beta_1,\beta_2,...\beta_n;\, x)=\sum _{r=0}^\infty \frac {(\alpha_1)_r(\alpha_2)_r...(\alpha_m)_r\;x^r}{(\beta_1)_r(\beta_2)_r...(\beta_n)_r\;r!} \end{equation}\]
% formula_record_end

% formula_record_start
% formula_id: formula-sheet-14-013
% page_slug: formula-sheet-14
% page_title: Formula Sheet
% page_url: https://formulasheet.com/#q|l|4988
% subject_slug: formulasheet
% subject_path: /formulasheet/
% section: Chapter 9
% subsection: 
% formula_name: 
% source: formulasheet-latex
% index_global: 367
% index_in_page: 13
\[\begin{equation} \label{simplehypergeo} _2F_1(\alpha,\beta;\gamma;x)=\sum _{r=0}^{\infty }\frac {(\alpha)_{r}(\beta)_{r}\,x^r}{(\gamma)_r\,r!} \end{equation}\]
% formula_record_end

% formula_record_start
% formula_id: formula-sheet-14-014
% page_slug: formula-sheet-14
% page_title: Formula Sheet
% page_url: https://formulasheet.com/#q|l|4988
% subject_slug: formulasheet
% subject_path: /formulasheet/
% section: Chapter 9
% subsection: 
% formula_name: 
% source: formulasheet-latex
% index_global: 368
% index_in_page: 14
\[\begin{equation} \label{poch0} (\alpha)_0 = 1 \end{equation}\]
% formula_record_end

% formula_record_start
% formula_id: formula-sheet-14-015
% page_slug: formula-sheet-14
% page_title: Formula Sheet
% page_url: https://formulasheet.com/#q|l|4988
% subject_slug: formulasheet
% subject_path: /formulasheet/
% section: Chapter 9
% subsection: 
% formula_name: 
% source: formulasheet-latex
% index_global: 369
% index_in_page: 15
\[\begin{equation} \label{pochhammer} (\alpha)_r=\alpha(\alpha+1)(\alpha+2) ... (\alpha+r-1) = \frac{\Gamma(\alpha + r)}{\Gamma(\alpha)} \end{equation}\]
% formula_record_end

% formula_record_start
% formula_id: formula-sheet-15-001
% page_slug: formula-sheet-15
% page_title: Formula Sheet
% page_url: https://formulasheet.com/#q|l|1229
% subject_slug: formulasheet
% subject_path: /formulasheet/
% section: Volume (geometry)
% subsection: Volume of a cone
% formula_name: 
% source: formulasheet-latex
% index_global: 370
% index_in_page: 1
\[\begin{equation} V = \frac{1}{3} \pi r^2 h \end{equation}\]
% formula_record_end

% formula_record_start
% formula_id: formula-sheet-15-002
% page_slug: formula-sheet-15
% page_title: Formula Sheet
% page_url: https://formulasheet.com/#q|l|1229
% subject_slug: formulasheet
% subject_path: /formulasheet/
% section: Volume (geometry)
% subsection: Volume of a cube
% formula_name: 
% source: formulasheet-latex
% index_global: 371
% index_in_page: 2
\[\begin{equation} V = s^3 \end{equation}\]
% formula_record_end

% formula_record_start
% formula_id: formula-sheet-15-003
% page_slug: formula-sheet-15
% page_title: Formula Sheet
% page_url: https://formulasheet.com/#q|l|1229
% subject_slug: formulasheet
% subject_path: /formulasheet/
% section: Volume (geometry)
% subsection: Volume of a cylinder
% formula_name: 
% source: formulasheet-latex
% index_global: 372
% index_in_page: 3
\[\begin{equation} V = \pi r^2 h \end{equation}\]
% formula_record_end

% formula_record_start
% formula_id: formula-sheet-15-004
% page_slug: formula-sheet-15
% page_title: Formula Sheet
% page_url: https://formulasheet.com/#q|l|1229
% subject_slug: formulasheet
% subject_path: /formulasheet/
% section: Volume (geometry)
% subsection: Volume of a rectangular prism
% formula_name: 
% source: formulasheet-latex
% index_global: 373
% index_in_page: 4
\[\begin{equation} V = l w h \end{equation}\]
% formula_record_end

% formula_record_start
% formula_id: formula-sheet-15-005
% page_slug: formula-sheet-15
% page_title: Formula Sheet
% page_url: https://formulasheet.com/#q|l|1229
% subject_slug: formulasheet
% subject_path: /formulasheet/
% section: Volume (geometry)
% subsection: Volume of a regular tetrahedron
% formula_name: 
% source: formulasheet-latex
% index_global: 374
% index_in_page: 5
\[\begin{equation} V = \frac{s^3}{6 \sqrt{2}} \end{equation}\]
% formula_record_end

% formula_record_start
% formula_id: formula-sheet-15-006
% page_slug: formula-sheet-15
% page_title: Formula Sheet
% page_url: https://formulasheet.com/#q|l|1229
% subject_slug: formulasheet
% subject_path: /formulasheet/
% section: Volume (geometry)
% subsection: Volume of a sphere
% formula_name: 
% source: formulasheet-latex
% index_global: 375
% index_in_page: 6
\[\begin{equation} \!V = \frac{4}{3}\pi r^3 \end{equation}\]
% formula_record_end

% formula_record_start
% formula_id: formula-sheet-15-007
% page_slug: formula-sheet-15
% page_title: Formula Sheet
% page_url: https://formulasheet.com/#q|l|1229
% subject_slug: formulasheet
% subject_path: /formulasheet/
% section: Volume (geometry)
% subsection: Volume of an ellipsoid
% formula_name: 
% source: formulasheet-latex
% index_global: 376
% index_in_page: 7
\[\begin{equation} V = \frac{4}{3} \pi a b c \end{equation}\]
% formula_record_end